\chapter{Integral Multivariabel} \label{mi:chapter}


%%%%%%%%%%%%%%%%%%%%%%%%%%%%%%%%%%%%%%%%%%%%%%%%%%%%%%%%%%%%%%%%%%%%%%%%%%%%%%

\section{Integral Riemann pada persegi panjang}
\label{sec:rirect}

%mbxINTROSUBSECTION

\sectionnotes{2--3 kuliah}

Seperti dalam \chapterref{int:chapter}, kita mendefinisikan integral Riemann
menggunakan integral Darboux atas dan bawah. Gagasan dalam bagian ini sangat
mirip dengan pengintegralan dalam satu dimensi. Kerumitannya terutama menyangkut
notasi. Perbedaan antara satu dan beberapa dimensi akan makin tampak pada
bagian-bagian berikutnya.

\subsection{Persegi panjang dan partisi}

\begin{defn}
Misalkan $(a_1,a_2,\ldots,a_n)$ dan
$(b_1,b_2,\ldots,b_n)$ memenuhi $a_k \leq b_k$ untuk setiap $k$.
Himpunan
$[a_1,b_1] \times
[a_2,b_2] \times \cdots \times
[a_n,b_n]$ disebut \emph{\myindex{persegi panjang tertutup}}%
\index{persegi panjang}.
Kadang-kadang berguna untuk mengizinkan $a_k = b_k$, dengan
$[a_k,b_k] = \{ a_k \}$.
Jika $a_k < b_k$ untuk setiap $k$, maka himpunan
$(a_1,b_1) \times
(a_2,b_2) \times \cdots \times
(a_n,b_n)$ disebut \emph{\myindex{persegi panjang terbuka}}.

Untuk persegi panjang terbuka atau tertutup
$R \coloneqq [a_1,b_1] \times
[a_2,b_2] \times \cdots \times
[a_n,b_n] \subset \R^n$
atau
$R \coloneqq (a_1,b_1) \times
(a_2,b_2) \times \cdots \times
(a_n,b_n) \subset \R^n$,
kita definisikan
\emph{volume berdimensi $n$}%
\index{volume berdimensi n@volume berdimensi $n$!persegi panjang}%
\index{volume persegi panjang} dengan
\glsadd{not:ndimvolume}
\begin{equation*}
V(R) \coloneqq
(b_1-a_1)
(b_2-a_2)
\cdots
(b_n-a_n) .
\end{equation*}

Suatu \emph{\myindex{partisi}} $P$ dari persegi panjang tertutup
$R = [a_1,b_1] \times
[a_2,b_2] \times \cdots \times
[a_n,b_n]$
ditentukan oleh partisi-partisi $P_1,P_2,\ldots,P_n$ dari interval
$[a_1,b_1], [a_2,b_2],\ldots, [a_n,b_n]$.
Kita tulis $P=(P_1,P_2,\ldots,P_n)$.
Artinya, untuk setiap $k=1,2,\ldots,n$ terdapat bilangan bulat positif
$\ell_k$ dan koleksi berhingga terindeks
$P_k = \{ x_{k,0},x_{k,1},x_{k,2},\ldots,x_{k,\ell_k} \}$ sedemikian sehingga,
jika $a_k<b_k$,
\begin{equation*}
a_k = x_{k,0} < x_{k,1} < x_{k,2} < \cdots < x_{k,{\ell_k}-1} < x_{k,\ell_k} = b_k .
\end{equation*}
\emph{Konvensi edisi Indonesia.} Untuk kasus degenerat $a_k=b_k$, ambil
$\ell_k=1$ dan $x_{k,0}=x_{k,1}=a_k$. Dengan demikian, definisi partisi ini
juga mencakup persegi panjang dengan satu atau lebih sisi yang panjangnya nol,
sebagaimana diizinkan di atas.

Dengan memilih $n$ bilangan bulat $j_1,j_2,\ldots,j_n$ yang memenuhi
$j_k \in \{ 1,2,\ldots,\ell_k \}$, kita memperoleh
\emph{\myindex{subpersegi panjang}}
\begin{equation*}
[x_{1,j_1-1}~,~ x_{1,j_1}]
\times
[x_{2,j_2-1}~,~ x_{2,j_2}]
\times
\cdots
\times
[x_{n,j_n-1}~,~ x_{n,j_n}] .
\end{equation*}
Kita urutkan subpersegi panjang tersebut dengan cara apa pun dan menyebut
$\{R_1,R_2,\ldots,R_N\}$ sebagai subpersegi panjang yang bersesuaian dengan
partisi $P$ dari $R$, atau, secara singkat, subpersegi panjang dari $P$.
Dengan kata lain, kita membagi persegi panjang semula menjadi banyak
subpersegi panjang yang lebih kecil. Lihat \figureref{mv:figrect}.

\begin{myfigureht}
\myincludepdft{figrect}{%
Diagram sebuah persegi panjang yang dibagi oleh dua garis horizontal dan dua
garis vertikal menjadi 9 persegi panjang yang lebih kecil. Empat titik sudut
pada sumbu horizontal, dari kiri ke kanan, berlabel x subskrip 1 koma 0,
x subskrip 1 koma 1, x subskrip 1 koma 2, dan x subskrip 1 koma 3. Pada sumbu
vertikal, dari bawah ke atas, titik-titik sudut berlabel x subskrip 2 koma 0,
x subskrip 2 koma 1, x subskrip 2 koma 2, dan x subskrip 2 koma 3. Persegi
panjang yang lebih kecil diberi label mengikuti pola berkelok dari kiri atas
ke kanan bawah: R subskrip 1, R subskrip 2, dan seterusnya hingga R subskrip 9.}
\caption{Contoh partisi persegi panjang di $\R^2$. Urutan subpersegi panjang
tidak penting.\label{mv:figrect}}
\end{myfigureht}

Misalkan $R \subset \R^n$ persegi panjang tertutup dan
$f \colon R \to \R$ fungsi terbatas. Misalkan $P$ partisi dari $R$ dengan
$N$ subpersegi panjang $R_1,R_2,\ldots,R_N$.
Definisikan
\glsadd{not:lowerdarbouxsum}
\glsadd{not:upperdarbouxsum}
\begin{align*}
& m_i \coloneqq \inf \bigl\{ f(x) : x \in R_i \bigr\} , & 
& M_i \coloneqq \sup \bigl\{ f(x) : x \in R_i \bigr\} , \\
& L(P,f) \coloneqq
\sum_{i=1}^N m_i V(R_i) , &
& U(P,f) \coloneqq
\sum_{i=1}^N M_i V(R_i) .
\end{align*}
Kita menyebut $L(P,f)$ sebagai \emph{\myindex{jumlah Darboux bawah}} dan
$U(P,f)$ sebagai \emph{\myindex{jumlah Darboux atas}}\index{jumlah Darboux}.
\end{defn}

Untuk melihat kaitannya dengan notasi $\Delta$ dari definisi satu variabel,
perhatikan bahwa jika
\begin{equation*}
R_i = [x_{1,j_1-1}~,~ x_{1,j_1}]
\times
[x_{2,j_2-1}~,~ x_{2,j_2}]
\times
\cdots
\times
[x_{n,j_n-1}~,~ x_{n,j_n}] ,
\end{equation*}
maka
\begin{equation*}
V(R_i)
=
(x_{1,j_1}-x_{1,j_1-1})
(x_{2,j_2}-x_{2,j_2-1})
\cdots
(x_{n,j_n}-x_{n,j_n-1})
= 
\Delta x_{1,j_1}
\Delta x_{2,j_2}
\cdots
\Delta x_{n,j_n} .
\end{equation*}
Mudah dilihat (dan diserahkan kepada pembaca) bahwa subpersegi panjang dari $P$
menutupi $R$ semula dan jumlah volumenya sama dengan volume $R$. Artinya,
\begin{equation*}
R= \bigcup_{k=1}^N R_k , \qquad \text{dan} \qquad
V(R) = \sum_{k=1}^N V(R_k).
\end{equation*}

Pengindeksan dalam definisi mungkin tampak rumit, tetapi untungnya kita jarang
perlu kembali langsung ke definisi tersebut.

\begin{prop} \label{mv:sumulbound:prop}
Misalkan $R \subset \R^n$ persegi panjang tertutup dan
$f \colon R \to \R$ fungsi terbatas. Misalkan $m, M \in \R$ memenuhi
$m \leq f(x) \leq M$ untuk setiap $x \in R$. Maka, untuk setiap partisi
$P$ dari~$R$,
\begin{equation*}
m \, V(R) \leq
L(P,f) \leq U(P,f)
\leq M\, V(R) .
\end{equation*}
\end{prop}

\begin{proof}
Misalkan $P$ partisi dari $R$. Untuk setiap $i$, berlaku
$m \leq m_i \leq M_i \leq M$. Selain itu,
$\sum_{i=1}^N V(R_i) = V(R)$. Karena itu,
\begin{multline*}
m \, V(R) =
m \left( \sum_{i=1}^N V(R_i) \right)
=
\sum_{i=1}^N m \, V(R_i)
\leq
\sum_{i=1}^N m_i \, V(R_i)
\leq
\\
\leq
\sum_{i=1}^N M_i \, V(R_i)
\leq
\sum_{i=1}^N M \,V(R_i)
=
M \left( \sum_{i=1}^N V(R_i) \right)
=
M \,V(R) .  \qedhere
\end{multline*}
\end{proof}

\subsection{Integral atas dan bawah}

Berdasarkan \propref{mv:sumulbound:prop}, himpunan jumlah Darboux atas dan bawah
merupakan himpunan terbatas, sehingga kita dapat mengambil infimum dan
supremumnya. Seperti dalam kasus satu variabel, kita definisikan integral bawah
dan atas.

\begin{defn}
Misalkan $f \colon R \to \R$ fungsi terbatas pada persegi panjang tertutup
$R \subset \R^n$.
Definisikan
\glsadd{not:lowerdarbouxR}
\glsadd{not:upperdarbouxR}
\begin{equation*}
%mbxlatex \begin{aligned}
%mbxlatex &
\underline{\int_R} f
\coloneqq \sup \, \bigl\{ L(P,f) : P \text{ partisi dari } R \bigr\} , 
%mbxSTARTIGNORE
\qquad
%mbxENDIGNORE
%mbxlatex \\
%mbxlatex &
\overline{\int_R} f
\coloneqq \inf \, \bigl\{ U(P,f) : P \text{ partisi dari } R \bigr\} .
%mbxlatex \end{aligned}
\end{equation*}
Kita menyebut $\underline{\int}$ sebagai
\emph{\myindex{integral Darboux bawah}}\index{integral Darboux} dan
$\overline{\int}$ sebagai \emph{\myindex{integral Darboux atas}}.
\end{defn}

Untuk menghaluskan partisi suatu persegi panjang, kita menghaluskan
partisi-partisi interval yang mendasarinya.

\begin{defn}
Misalkan $R \subset \R^n$ persegi panjang tertutup.
Misalkan $P = ( P_1, P_2, \ldots, P_n )$ dan
$\widetilde{P} = ( \widetilde{P}_1, \widetilde{P}_2, \ldots, \widetilde{P}_n )$
merupakan partisi dari $R$. Kita menyebut $\widetilde{P}$ sebagai
\emph{penghalusan}\index{penghalusan partisi} dari $P$ jika, sebagai himpunan,
$P_k \subset \widetilde{P}_k$ untuk setiap $k = 1,2,\ldots,n$.
\end{defn}

Jika $\widetilde{P}$ merupakan penghalusan dari $P$, maka subpersegi panjang
dari $P$ merupakan gabungan subpersegi panjang dari $\widetilde{P}$. Secara
sederhana, dalam suatu penghalusan kita mengambil subpersegi panjang dari $P$,
memotongnya menjadi subpersegi panjang yang lebih kecil, dan menamai partisi
yang dihasilkan $\widetilde{P}$. Lihat \figureref{mv:figrectpart}.

\begin{myfigureht}
\myincludepdft{figrectpart}{%
Diagram persegi panjang terbagi seperti diagram sebelumnya, dengan tanda
tambahan. Dua pemotongan vertikal tambahan dan satu pemotongan horizontal
ditandai dengan garis putus-putus. Persegi panjang itu kini terbagi menjadi 20
persegi panjang kecil, yang diberi label R tilde subskrip 1 hingga R tilde
subskrip 20 tanpa urutan tertentu. Titik-titik pada sumbu horizontal dan
vertikal ditandai seperti sebelumnya, tetapi sekarang secara berurutan dengan
x tilde subskrip 1 koma 0 hingga x tilde subskrip 1 koma 5 dan x tilde subskrip
2 koma 0 hingga x tilde subskrip 2 koma 4. Pada pemotongan lama ditampilkan
baik tanda lama maupun tanda baru; misalnya, x subskrip 1 koma 2 juga bertanda
x tilde subskrip 1 koma 3.}
\caption{Contoh penghalusan partisi dari \figureref{mv:figrect}.
\myquote{Pemotongan} baru ditandai dengan garis putus-putus. Urutan persis
subpersegi panjang baru tidak penting.\label{mv:figrectpart}}
\end{myfigureht}

\begin{prop} \label{mv:prop:refinement}
Misalkan $R \subset \R^n$ persegi panjang tertutup, $P$ partisi dari $R$, dan
$\widetilde{P}$ penghalusan dari $P$. Jika $f \colon R \to \R$ terbatas, maka
\begin{equation*}
L(P,f) \leq L(\widetilde{P},f) 
\qquad \text{dan} \qquad
U(\widetilde{P},f) \leq U(P,f) .
\end{equation*}
\end{prop}

\begin{proof}
Kita buktikan ketaksamaan pertama; ketaksamaan kedua diperoleh dengan cara
serupa. Misalkan $R_1,R_2,\ldots,R_N$ subpersegi panjang dari $P$ dan
$\widetilde{R}_1,\widetilde{R}_2,\ldots,\widetilde{R}_{\widetilde{N}}$
subpersegi panjang dari $\widetilde{P}$. Misalkan $I_k$ himpunan semua indeks
$j$ sedemikian sehingga $\widetilde{R}_j \subset R_k$. Sebagai contoh, pada
gambar \ref{mv:figrect} dan \ref{mv:figrectpart},
$I_4 = \{ 6, 7, 8, 9 \}$ karena
$R_4 =
\widetilde{R}_6 \cup \widetilde{R}_7 \cup
\widetilde{R}_8 \cup \widetilde{R}_9$.
Maka,
\begin{equation*}
R_k = \bigcup_{j \in I_k} \widetilde{R}_j,
\qquad
V(R_k) = \sum_{j \in I_k} V(\widetilde{R}_j).
\end{equation*}

Seperti biasa, misalkan
$m_k \coloneqq \inf \bigl\{ f(x) : x \in R_k \bigr\}$ dan
$\widetilde{m}_j \coloneqq \inf \bigl\{ f(x) : x \in \widetilde{R}_j \bigr\}$.
Jika $j \in I_k$, maka $m_k \leq \widetilde{m}_j$. Dengan demikian,
\begin{equation*}
%mbxlatex \begin{aligned}
L(P,f)
=
\sum_{k=1}^N m_k V(R_k)
%mbxlatex & 
=
\sum_{k=1}^N \sum_{j\in I_k} m_k V(\widetilde{R}_j)
%mbxlatex \\
%mbxlatex & 
\leq
\sum_{k=1}^N \sum_{j\in I_k} \widetilde{m}_j V(\widetilde{R}_j)
=
\sum_{j=1}^{\widetilde{N}} \widetilde{m}_j V(\widetilde{R}_j) = L(\widetilde{P},f) .
\qedhere
%mbxlatex \end{aligned}
\end{equation*}
\end{proof}

Pokok penting proposisi berikut adalah bahwa integral Darboux bawah kurang dari
atau sama dengan integral Darboux atas.

\begin{prop} \label{mv:intulbound:prop}
Misalkan $R \subset \R^n$ persegi panjang tertutup dan
$f \colon R \to \R$ fungsi terbatas. Misalkan $m, M \in \R$ memenuhi
$m \leq f(x) \leq M$ untuk setiap $x \in R$. Maka,
\begin{equation}
\label{mv:intulbound:eq}
m \, V(R) \leq
\underline{\int_R} f \leq \overline{\int_R} f
\leq M \, V(R).
\end{equation}
\end{prop}

\begin{proof}
Untuk setiap partisi $P$, berdasarkan \propref{mv:sumulbound:prop},
\begin{equation*}
m\,V(R) \leq L(P,f) \leq U(P,f) \leq M\,V(R).
\end{equation*}
Dengan mengambil supremum $L(P,f)$ dan infimum $U(P,f)$ atas semua partisi $P$,
kita memperoleh ketaksamaan pertama dan terakhir dalam
\eqref{mv:intulbound:eq}.

Ketaksamaan kunci dalam \eqref{mv:intulbound:eq} adalah ketaksamaan tengah.
Misalkan $P=(P_1,P_2,\ldots,P_n)$ dan
$Q=(Q_1,Q_2,\ldots,Q_n)$ merupakan partisi dari $R$. Definisikan
$\widetilde{P} = ( \widetilde{P}_1,\widetilde{P}_2,\ldots,\widetilde{P}_n )$
dengan menetapkan $\widetilde{P}_k \coloneqq P_k \cup Q_k$ untuk setiap $k$.
Maka $\widetilde{P}$ merupakan partisi dari $R$ serta penghalusan dari $P$ dan
juga dari $Q$. Berdasarkan \propref{mv:prop:refinement},
$L(P,f) \leq L(\widetilde{P},f)$ dan
$U(\widetilde{P},f) \leq U(Q,f)$. Karena itu,
\begin{equation*}
L(P,f) \leq L(\widetilde{P},f) \leq U(\widetilde{P},f) \leq U(Q,f) .
\end{equation*}
Dengan kata lain, untuk dua partisi sembarang $P$ dan $Q$, berlaku
$L(P,f) \leq U(Q,f)$.
Berdasarkan \volIref{\propref*{vI-infsupineq:prop} dari jilid I}{\propref{infsupineq:prop}},
kita peroleh
\begin{equation*}
\sup \, \bigl\{ L(P,f) : P \text{ partisi dari } R \bigr\}
\leq
\inf \, \bigl\{ U(P,f) : P \text{ partisi dari } R \bigr\} .
\end{equation*}
Dengan kata lain, $\underline{\int_R} f \leq \overline{\int_R} f$.
\end{proof}

\subsection{Integral Riemann}

Sekarang kita memiliki semua yang diperlukan untuk mendefinisikan integral
Riemann dalam $n$ dimensi pada persegi panjang. Seperti dalam satu dimensi,
integral Riemann hanya didefinisikan pada kelas fungsi tertentu, yang disebut
fungsi terintegralkan secara Riemann.

\begin{defn}
Misalkan $R \subset \R^n$ persegi panjang tertutup dan
$f \colon R \to \R$ fungsi terbatas sedemikian sehingga
\begin{equation*}
\underline{\int_R} f(x)\,dx = \overline{\int_R} f(x)\,dx .
\end{equation*}
Maka $f$ dikatakan \emph{\myindex{terintegralkan secara Riemann}}, dan terkadang
kita cukup mengatakan \emph{\myindex{terintegralkan}}. Himpunan fungsi yang
terintegralkan secara Riemann pada $R$ kita lambangkan dengan
$\sR(R)$.\glsadd{not:integrablefuncR}
Untuk $f \in \sR(R)$, definisikan \emph{\myindex{integral Riemann}}
\glsadd{not:riemannintR}
\begin{equation*}
\int_R f \coloneqq 
\underline{\int_R} f = \overline{\int_R} f .
\end{equation*}
\end{defn}

Jika variabel $x \in \R^n$ perlu ditonjolkan, kita tulis
\begin{equation*}
\int_R f(x)\,dx,
\qquad
\int_R f(x_1,\ldots,x_n)\,dx_1 \cdots dx_n,
\qquad
\text{atau}
\qquad
\int_R f(x)\,dV .
\end{equation*}
Jika $R \subset \R^2$, biasanya kita mengatakan luas alih-alih volume dan
menulis
\begin{equation*}
\int_R f(x)\,dA .
\end{equation*}

\propref{mv:intulbound:prop} langsung menghasilkan proposisi berikut.

\begin{prop} \label{mv:intbound:prop}
Misalkan $f \colon R \to \R$ fungsi yang terintegralkan secara Riemann pada
persegi panjang tertutup $R \subset \R^n$. Misalkan $m, M \in \R$ memenuhi
$m \leq f(x) \leq M$ untuk setiap $x \in R$. Maka,
\begin{equation*}
m \, V(R) \leq
\int_{R} f
\leq M \, V(R) .
\end{equation*}
\end{prop}

\begin{example}
Fungsi konstan terintegralkan secara Riemann. Buktinya, andaikan
$f(x) = c$ untuk setiap $x \in R$. Maka,
\begin{equation*}
c \, V(R) \leq \underline{\int_R} f \leq \overline{\int_R} f \leq c\, V(R) .
\end{equation*}
Jadi, $f$ terintegralkan dan, lebih lanjut, $\int_R f = c\,V(R)$.
\end{example}

Pembuktian linearitas dan kemonotonan hampir sepenuhnya sama dengan
pembuktian dalam satu variabel. Dua proposisi berikut kita serahkan sebagai
latihan.

\begin{samepage}
\begin{prop}[Linearitas] \label{mv:intlinearity:prop}
\index{linearitas integral}
Misalkan $R \subset \R^n$ persegi panjang tertutup,
$f,g \in \sR(R)$, dan $\alpha \in \R$.
\begin{enumerate}[(i)]
\item $\alpha f \in \sR(R)$ dan
\begin{equation*}
\int_R \alpha f = \alpha \int_R f .
\end{equation*}
\item $f+g \in \sR(R)$ dan
\begin{equation*}
\int_R (f+g) = 
\int_R f
+
\int_R g .
\end{equation*}
\end{enumerate}
\end{prop}
\end{samepage}

\begin{prop}[Kemonotonan]
\index{kemonotonan integral}
Misalkan $R \subset \R^n$ persegi panjang tertutup, $f,g \in \sR(R)$, dan
andaikan $f(x) \leq g(x)$ untuk setiap $x \in R$. Maka,
\begin{equation*}
\int_R f 
\leq
\int_R g .
\end{equation*}
\end{prop}

Memeriksa keterintegralan langsung dari definisi sering melibatkan teknik
berikut, sama seperti dalam kasus satu variabel.

\begin{prop} \label{mv:prop:upperlowerepsilon}
Misalkan $R \subset \R^n$ persegi panjang tertutup dan
$f \colon R \to \R$ fungsi terbatas. Maka $f \in \sR(R)$ jika dan hanya jika,
untuk setiap $\epsilon > 0$, terdapat partisi $P$ dari $R$ sedemikian sehingga
\begin{equation*}
U(P,f) - L(P,f) < \epsilon .
\end{equation*}
\end{prop}

\begin{proof}
Pertama, jika $f$ terintegralkan, maka supremum $L(P,f)$ dan infimum $U(Q,f)$,
masing-masing atas semua partisi $P$ dan $Q$, bernilai sama. Karena itu,
infimum $U(P,f)-L(Q,f)$ adalah nol. Dengan mengambil penghalusan bersama
$\widetilde{P}$ dari $P$ dan $Q$, kita memperoleh
$U(\widetilde{P},f)-L(\widetilde{P},f) \leq U(P,f)-L(Q,f)$.
Jadi, infimum $U(P,f)-L(P,f)$ atas semua partisi $P$ adalah nol. Dengan
demikian, untuk setiap $\epsilon > 0$, terdapat partisi $P$ sedemikian sehingga
$U(P,f) - L(P,f) < \epsilon$.

Untuk arah sebaliknya, diberikan $\epsilon > 0$, pilih $P$ sedemikian sehingga
$U(P,f) - L(P,f) < \epsilon$.
\begin{equation*}
\overline{\int_R} f - 
\underline{\int_R} f 
\leq
U(P,f) - L(P,f)
< \epsilon .
\end{equation*}
Karena $\overline{\int_R} f \geq \underline{\int_R} f$ dan pernyataan di atas
berlaku untuk setiap $\epsilon > 0$, kita simpulkan bahwa
$\overline{\int_R} f = \underline{\int_R} f$ dan $f \in \sR(R)$.
\end{proof}

Misalkan $f \colon S \to \R$ suatu fungsi dan $R \subset S$ persegi panjang
tertutup. Jika pembatasan $f|_R$ terintegralkan, untuk ringkasnya kita katakan
bahwa $f$ \emph{terintegralkan pada $R$}, atau $f \in \sR(R)$, dan kita tulis
\begin{equation*}
\int_R f \coloneqq \int_R f|_R .
\end{equation*}

\begin{prop} \label{mv:prop:integralsmallerset}
Misalkan $S \subset \R^n$ persegi panjang tertutup. Jika
$f \colon S \to \R$ terintegralkan dan $R \subset S$ persegi panjang tertutup,
maka $f$ terintegralkan pada $R$.
\end{prop}

\begin{proof}
Jika salah satu sisi $R$ panjangnya nol, maka setiap subpersegi panjang dalam
setiap partisi $R$ bervolume nol. Karena itu semua jumlah Darboux bawah dan atas
sama dengan nol, sehingga $f|_R$ terintegralkan. Jadi, untuk sisa pembuktian,
andaikan semua sisi $R$ mempunyai panjang positif.

Diberikan $\epsilon > 0$, pilih partisi $P=(P_1,\ldots,P_n)$ dari $S$
sedemikian sehingga $U(P,f)-L(P,f) < \epsilon$. Dengan memperhalus $P$ jika
perlu, andaikan semua titik ujung $R$ termuat dalam $P$. Artinya, jika
$R = [a_1,b_1] \times [a_2,b_2] \times \cdots \times [a_n,b_n]$, maka
$a_i,b_i \in P_i$.
Misalkan $\widetilde{P} = (\widetilde{P}_1,\ldots,\widetilde{P}_n)$ partisi dari
$R$ yang diberikan oleh $\widetilde{P}_i = P_i \cap [a_i,b_i]$.
Subpersegi panjang dari $\widetilde{P}$ merupakan subpersegi panjang dari~$P$;
dengan kata lain, $R$ merupakan gabungan subpersegi panjang dari~$P$.
Bagi subpersegi panjang dari $P$ menjadi dua koleksi: misalkan
$R_1,R_2,\ldots,R_K$ subpersegi panjang dari $P$ yang juga merupakan
subpersegi panjang dari $\widetilde{P}$, dan misalkan
$R_{K+1},\ldots,R_N$ sisanya. Lihat \figureref{fig:figrectsubrect}.
Seperti biasa, misalkan $m_k$ dan $M_k$ masing-masing merupakan infimum dan
supremum $f$ pada $R_k$. Maka,
\begin{equation*}
\begin{split}
\epsilon & > 
U(P,f)-L(P,f)
=
\sum_{k=1}^K (M_k-m_k) V(R_k)
+
\sum_{k=K+1}^N (M_k-m_k) V(R_k)
\\
&
\geq
\sum_{k=1}^K (M_k-m_k) V(R_k)
=
U(\widetilde{P},f|_R)-L(\widetilde{P},f|_R) .
\end{split}
\end{equation*}
Karena itu, $f|_R$ terintegralkan.
\end{proof}
\begin{myfigureht}
\myincludepdft{figrectsubrect}{%
Sebuah persegi panjang dibagi menjadi subpersegi panjang seperti sebelumnya;
kali ini, 2 pemotongan horizontal dan 3 pemotongan vertikal menghasilkan 12
persegi panjang yang lebih kecil. Empat persegi panjang kecil pada baris kedua
dan ketiga serta kolom kedua dan ketiga diberi arsiran gelap dan berlabel
R subskrip 1, R subskrip 2, R subskrip 3, dan R subskrip 4.}
\caption{Partisi persegi panjang besar $S$ yang juga menghasilkan partisi
persegi panjang lebih kecil (diarsir dan diberi garis tepi) $R \subset S$.
Subpersegi panjang $R_1,R_2,R_3,R_4$ merupakan subpersegi panjang dari
$\widetilde{P} = \bigl( \{ x_{1,1}, x_{1,2} , x_{1,3} \} ,
\{ x_{2,1}, x_{2,2} , x_{2,3} \} \bigr)$.\label{fig:figrectsubrect}}
\end{myfigureht}

\subsection{Integral fungsi kontinu}

Walaupun kelak kita akan membuktikan hasil yang lebih umum, ada baiknya kita
mulai dengan keterintegralan fungsi kontinu. Untuk itu, kita perlu mengukur
kehalusan partisi. Dalam satu variabel, kita mengukur panjang subinterval.
Dalam beberapa variabel, kita mengukur panjang sisi subpersegi panjang.
Kita katakan persegi panjang $R = [a_1,b_1] \times
[a_2,b_2] \times \cdots \times
[a_n,b_n]$ memiliki \emph{panjang setiap sisi yang tidak melebihi $\alpha$}%
\index{sisi terpanjang} jika
$b_k-a_k \leq \alpha$ untuk setiap $k=1,2,\ldots,n$.

\begin{prop} \label{prop:diameterrectangle}
Jika panjang setiap sisi persegi panjang $R \subset \R^n$ tidak melebihi
$\alpha$, maka, untuk setiap $x,y \in R$,
\begin{equation*}
\snorm{x-y} \leq \sqrt{n} \, \alpha .
\end{equation*}
\end{prop}

\begin{proof}
\begin{equation*}
\begin{split}
\snorm{x-y} 
& =
\sqrt{
{(x_1-y_1)}^2
+
{(x_2-y_2)}^2
+ \cdots +
{(x_n-y_n)}^2
}
\\
& \leq
\sqrt{
{(b_1-a_1)}^2
+
{(b_2-a_2)}^2
+ \cdots +
{(b_n-a_n)}^2
}
\\
& \leq
\sqrt{
{\alpha}^2
+
{\alpha}^2
+ \cdots +
{\alpha}^2
}
=
\sqrt{n} \, \alpha .  \qedhere
\end{split}
\end{equation*}
\end{proof}


\begin{thm} \label{mv:thm:contintrect}
Misalkan $R \subset \R^n$ persegi panjang tertutup. Jika
$f \colon R \to \R$ kontinu, maka $f \in \sR(R)$.
\end{thm}

\begin{proof}
Pembuktian ini serupa dengan pembuktian satu variabel, meskipun terdapat
beberapa kerumitan tambahan. Himpunan $R$ tertutup dan terbatas di $\R^n$,
sehingga kompak. Oleh karena itu, $f$ terbatas pada $R$.
Karena itu, $f$ kontinu seragam berdasarkan
\volIref{\thmref*{vI-thm:Xcompactfunifcont} dari jilid I}{\thmref{thm:Xcompactfunifcont}}.

Jika $V(R)=0$, maka setiap subpersegi panjang dari setiap partisi $R$ juga
bervolume nol. Semua jumlah Darboux bawah dan atasnya sama dengan nol, sehingga
$f\in\sR(R)$. Oleh sebab itu, selanjutnya andaikan $V(R)>0$.

Diberikan $\epsilon > 0$, pilih $\delta > 0$ sedemikian sehingga
$\snorm{x-y} < \delta$ mengakibatkan
$\babs{f(x)-f(y)} < \frac{\epsilon}{V(R)}$.

Ambil partisi $P$ dari $R$ sedemikian sehingga panjang sisi terpanjang setiap
subpersegi panjang benar-benar kurang dari $\frac{\delta}{\sqrt{n}}$.
Jika $x,y \in R_k$ untuk suatu subpersegi panjang $R_k$ dari $P$ dan $\alpha_k$
menyatakan panjang sisi terpanjang $R_k$, maka, berdasarkan proposisi,
$\snorm{x-y} \leq \sqrt{n}\,\alpha_k <
\sqrt{n}\frac{\delta}{\sqrt{n}} = \delta$. Karena itu,
\begin{equation*}
f(x)-f(y) \leq \babs{f(x)-f(y)} < \frac{\epsilon}{V(R)} .
\end{equation*}
Karena $f$ kontinu pada $R_k$ yang kompak, $f$ mencapai maksimum dan minimum
pada subpersegi panjang ini. Misalkan $x$ titik tempat $f$ mencapai maksimum
dan $y$ titik tempat $f$ mencapai minimum. Maka $f(x) = M_k$ dan
$f(y) = m_k$ dalam notasi dari definisi integral. Jadi,
\begin{equation*}
M_k-m_k = f(x)-f(y) < 
\frac{\epsilon}{V(R)} .
\end{equation*}
Dengan demikian,
\begin{equation*}
\begin{split}
U(P,f) - L(P,f)
& =
\left(
\sum_{k=1}^N
M_k V(R_k)
\right)
-
\left(
\sum_{k=1}^N
m_k V(R_k)
\right)
\\
& =
\sum_{k=1}^N
(M_k-m_k) V(R_k)
\\
& <
\frac{\epsilon}{V(R)}
\sum_{k=1}^N
V(R_k)
= \epsilon.
\end{split}
\avoidbreak
\end{equation*}
\propref{mv:prop:upperlowerepsilon} kemudian menyatakan bahwa $f \in \sR(R)$.
\end{proof}

\subsection{Pengintegralan fungsi bertumpuan kompak}

Misalkan $U \subset \R^n$ himpunan terbuka dan $f \colon U \to \R$ suatu
fungsi. \emph{\myindex{Tumpuan}} $f$ adalah himpunan
\begin{equation*}
\operatorname{supp} (f) \coloneqq
\overline{
\{ x \in U : f(x) \neq 0 \}
} ,
\end{equation*}
di mana penutupan diambil terhadap topologi subruang pada $U$. Mengambil
penutupan terhadap topologi subruang sama dengan
$\overline{
\{ x \in U : f(x) \neq 0 \}
} \cap U$, dengan penutupan pada ungkapan terakhir diambil terhadap ruang Euklides
$\R^n$. Secara khusus,
$\operatorname{supp} (f) \subset U$.
Tumpuan adalah penutupan (di $U$) himpunan titik tempat fungsi tidak nol.
Komplemennya di $U$ terbuka. Jika $x \in U$ dan $x$ tidak termasuk tumpuan
$f$, maka $f$ identik nol pada suatu lingkungan dari $x$.

Fungsi $f$ dikatakan \emph{\myindex{bertumpuan kompak}} jika $\supp(f)$
merupakan himpunan kompak.

\begin{example}
Fungsi $f \colon \R^2 \to \R$ yang didefinisikan oleh
\begin{equation*}
f(x,y) \coloneqq
\begin{cases}
-x{(x^2+y^2-1)}^2 & \text{jika } \sqrt{x^2+y^2} \leq 1, \\
0                 & \text{selainnya},
\end{cases}
\end{equation*}
bersifat kontinu dan tumpuannya adalah cakram satuan tertutup
$C(0,1) = \bigl\{ (x,y) : \sqrt{x^2 + y^2} \leq 1 \bigr\}$. Karena himpunan
ini kompak, $f$ bertumpuan kompak. Perhatikan bahwa fungsi tersebut bernilai
nol pada seluruh sumbu-$y$ dan pada lingkaran satuan. Namun, setiap titik
semacam itu yang terletak dalam cakram satuan tertutup tetap termasuk tumpuan,
sebab titik tersebut berada dalam penutupan himpunan titik tempat $f$ tidak
nol. Lihat \figureref{fig:compsup}.
\begin{myfigureht}
\myincludepdft{compsup_full}{%
Di sebelah kiri terdapat grafik permukaan suatu fungsi yang bernilai nol di
luar cakram satuan. Di dalam cakram, fungsi bernilai positif pada titik-titik
dengan x negatif dan bernilai negatif pada titik-titik dengan x positif. Di
sebelah kanan terdapat diagram cakram satuan: sebuah lingkaran berjari-jari 1
yang berpusat di titik asal, dengan bagian dalam diarsir.}
\caption{Fungsi bertumpuan kompak (kiri); tumpuannya adalah cakram satuan
tertutup (kanan).\label{fig:compsup}}
\end{myfigureht}
\end{example}

Jika $U \neq \R^n$, kita harus berhati-hati mengambil penutupan di dalam $U$.
Tinjau dua contoh berikut.

\begin{example}
Misalkan $B(0,1) \subset \R^2$ cakram satuan. Fungsi
$f \colon B(0,1) \to \R$ yang didefinisikan oleh
\begin{equation*}
f(x,y) \coloneqq
\begin{cases}
0                                & \text{jika } \sqrt{x^2+y^2} > \nicefrac{1}{2}, \\
\nicefrac{1}{2} - \sqrt{x^2+y^2} & \text{jika } \sqrt{x^2+y^2} \leq \nicefrac{1}{2},
\end{cases}
\end{equation*}
bersifat kontinu pada $B(0,1)$ dan tumpuannya adalah bola tertutup yang lebih
kecil, $C(0,\nicefrac{1}{2})$. Karena himpunan ini kompak, $f$ bertumpuan
kompak.

Fungsi $g \colon B(0,1) \to \R$ yang didefinisikan oleh
\begin{equation*}
g(x,y) \coloneqq
\begin{cases}
0 & \text{jika } x \leq 0, \\
x & \text{jika } x > 0,
\end{cases}
\end{equation*}
bersifat kontinu pada $B(0,1)$, tetapi tumpuannya adalah himpunan
$\bigl\{ (x,y) \in B(0,1) : x \geq 0 \bigr\}$. Secara khusus, $g$ tidak
bertumpuan kompak.
\end{example}

Sesungguhnya kita hanya perlu meninjau kasus $U=\R^n$. Berdasarkan
\exerciseref{exercise:contcompactsupportRn}, setiap fungsi kontinu bertumpuan
kompak pada himpunan terbuka $U \subset \R^n$ dapat diperluas menjadi fungsi
kontinu bertumpuan kompak pada $\R^n$. Jadi, meninjau $U=\R^n$ bukanlah
penyederhanaan yang berlebihan.

\begin{example}
Fungsi kontinu $f \colon B(0,1) \to \R$ yang diberikan oleh
$f(x,y) \coloneqq \sin\bigl(\frac{1}{1-x^2-y^2}\bigr)$
tidak bertumpuan kompak. Untuk setiap titik di $B(0,1)$, fungsi $f$ tidak
identik nol pada lingkungan mana pun dari titik tersebut. Jadi, tumpuannya
adalah seluruh cakram $B(0,1)$.
Fungsi ini tidak dapat diperluas seperti di atas menjadi fungsi kontinu pada
$\R^2$. Bahkan, mudah ditunjukkan bahwa tidak ada cara apa pun untuk memperluas
$f$ menjadi fungsi kontinu pada seluruh $\R^2$ (masalahnya terletak pada batas
cakram).
\end{example}

\begin{prop} \label{mv:prop:rectanglessupp}
Misalkan $f \colon \R^n \to \R$ fungsi kontinu bertumpuan kompak. Jika $R$ dan
$S$ merupakan persegi panjang tertutup sedemikian sehingga
$\operatorname{supp}(f) \subset R$
dan
$\operatorname{supp}(f) \subset S$, maka
\begin{equation*}
\int_S f = \int_R f .
\end{equation*}
\end{prop}

\begin{proof}
Pilih persegi panjang tertutup $T$ yang memuat $R\cup S$. Karena $f$ kontinu,
fungsi ini terintegralkan pada $T$, $R$, dan $S$. Selain itu,
$\operatorname{supp}(f)\subset R\cap S$, sehingga $f$ identik nol pada
$T\setminus R$ maupun $T\setminus S$. Dengan menerapkan
\exerciseref{mv:zerooutside} dua kali,
$\int_R f = \int_T f = \int_S f$.
\end{proof}

Berdasarkan proposisi ini, jika $f \colon \R^n \to \R$ bertumpuan kompak dan
terintegralkan pada persegi panjang $R$ yang memuat tumpuan tersebut, kita tulis
\begin{equation*}
\int f \coloneqq \int_R f \qquad \text{atau} \qquad 
\int_{\R^n} f \coloneqq \int_R f .
\end{equation*}
Sebagai contoh, jika $f$ kontinu dan bertumpuan kompak, maka
$\int_{\R^n} f$ ada.

\subsection{Latihan}

\begin{exercise} \label{exercise:contcompactsupportRn}
Misalkan $U \subset \R^n$ terbuka dan $f \colon U \to \R$ kontinu serta
bertumpuan kompak. Tunjukkan bahwa fungsi $\widetilde{f} \colon \R^n \to
\R$, dengan
\begin{equation*}
\widetilde{f}(x) \coloneqq
\begin{cases}
f(x) & \text{jika } x \in U, \\
0 & \text{selainnya,}
\end{cases}
\avoidbreak
\end{equation*}
bersifat kontinu.
\end{exercise}


\begin{exercise}
Buktikan \propref{mv:intlinearity:prop}.
\end{exercise}


\begin{exercise}
Misalkan $R \subset \R^n$ persegi panjang tertutup yang panjang salah satu
sisinya sama dengan 0. Untuk setiap fungsi terbatas $f \colon R \to \R$, tunjukkan
bahwa $f \in \sR(R)$ dan $\int_R f = 0$.
\end{exercise}

\begin{exercise} \label{mv:zerosiderectangle}
Misalkan $R \subset \R^n$ persegi panjang tertutup yang panjang salah satu
sisinya sama dengan 0, dan misalkan $S$ persegi panjang tertutup dengan $R \subset S$.
Jika $f \colon S \to \R$ fungsi terbatas sedemikian sehingga $f(x) = 0$ untuk
$x \in S \setminus R$, tunjukkan bahwa $f \in \sR(S)$ dan $\int_S f = 0$.
\end{exercise}

\begin{exercise}
Misalkan $f\colon \R^n \to \R$ sedemikian sehingga $f(x) \coloneqq 0$ jika
$x \neq 0$ dan $f(0) \coloneqq 1$. Tunjukkan langsung dari definisi bahwa $f$
terintegralkan pada
$R \coloneqq [-1,1] \times [-1,1] \times \cdots \times [-1,1]$, lalu tentukan
$\int_R f$.
\end{exercise}

\begin{exercise} \label{mv:zeroinside}
Misalkan $R \subset \R^n$ dan $S \subset \R^n$ persegi panjang tertutup dengan
$R \subset S$. Misalkan $h \colon S \to \R$ fungsi terbatas sedemikian sehingga
$h(x) = 0$ jika $x \notin \partial R$ (batas $R$). Tunjukkan bahwa
$h \in \sR(S)$ dan
\begin{equation*}
\int_{S} h = 0 .
\avoidbreak
\end{equation*}
Petunjuk: Nyatakan $h$ sebagai jumlah fungsi-fungsi seperti pada
\exerciseref{mv:zerosiderectangle}.
\end{exercise}

\begin{exercise} \label{mv:zerooutside}
Misalkan $R$ dan $R'$ dua persegi panjang tertutup di $\R^n$ dengan
$R' \subset R$. Misalkan $f \colon R \to \R$ termasuk dalam $\sR(R')$ dan
$f(x) = 0$ untuk $x \in R \setminus R'$. Tunjukkan bahwa $f \in \sR(R)$ dan
\begin{equation*}
\int_{R'} f = \int_R f .
\end{equation*}
Kerjakan dalam langkah-langkah berikut.
\begin{enumerate}[a)]
\item
Pertama, buktikan dengan asumsi tambahan bahwa $f(x) = 0$ setiap kali
$x \in \overline{R \setminus R'}$.
\item
Tuliskan $f(x) = g(x) + h(x)$, dengan $g(x) = 0$ setiap kali
$x \in \overline{R \setminus R'}$, sedangkan $h(x)$ bernilai nol kecuali
mungkin pada $\partial R'$. Kemudian tunjukkan
$\int_R h = \int_{R'} h = 0$ (lihat \exerciseref{mv:zeroinside}).
\item
Tunjukkan bahwa $\int_{R'} f = \int_R f$.
\end{enumerate}
\end{exercise}

\begin{exercise}
Misalkan $R' \subset \R^n$ dan $R'' \subset \R^n$ dua persegi panjang tertutup
sedemikian sehingga $R = R' \cup R''$ merupakan persegi panjang tertutup,
sedangkan $R' \cap R''$ merupakan persegi panjang tertutup yang panjang salah
satu sisinya sama dengan 0
(yakni, $V(R' \cap R'') = 0$). Misalkan $f \colon R \to \R$ fungsi sedemikian
sehingga $f \in \sR(R')$ dan $f \in \sR(R'')$. Tunjukkan bahwa
$f \in \sR(R)$ dan
\begin{equation*}
\int_{R} f = \int_{R'} f + \int_{R''} f .
\avoidbreak
\end{equation*}
Petunjuk: Lihat latihan sebelumnya.
\end{exercise}

\begin{exercise}
Buktikan versi yang lebih kuat dari \propref{mv:prop:rectanglessupp}.
Misalkan $f \colon \R^n \to \R$ fungsi bertumpuan kompak, tetapi tidak harus
kontinu. Buktikan bahwa jika $R$ persegi panjang tertutup sedemikian sehingga
$\operatorname{supp}(f) \subset R$ dan $f$ terintegralkan pada $R$, maka untuk
setiap persegi panjang tertutup lain $S$ dengan
$\operatorname{supp}(f) \subset S$, fungsi $f$ terintegralkan pada $S$ dan
$\int_S f = \int_R f$.
Petunjuk: Lihat \exerciseref{mv:zerooutside}.
\end{exercise}

\begin{exercise}
Misalkan $R$ dan $S$ persegi panjang tertutup di $\R^n$. Definisikan
$f \colon \R^n \to \R$ dengan $f(x) \coloneqq 1$ jika $x \in R$, dan
$f(x) \coloneqq 0$ selainnya. Buktikan bahwa $f$ terintegralkan pada $S$ dan
hitung $\int_S f$. Petunjuk: Tangani kasus $S \cap R = \varnothing$ secara
terpisah; jika irisannya tidak kosong, tinjau $S \cap R$.
\end{exercise}

\begin{samepage}
\begin{exercise}
Misalkan $R \coloneqq [0,1] \times [0,1] \subset \R^2$.
\begin{enumerate}[a)]
\item
Misalkan $f \colon R \to \R$ didefinisikan oleh
\begin{equation*}
f(x,y) \coloneqq 
\begin{cases}
1 & \text{jika } x = y, \\
0 & \text{selainnya.}
\end{cases}
\end{equation*}
Tunjukkan bahwa $f \in \sR(R)$ dan hitung $\int_R f$.
\item
Misalkan $f \colon R \to \R$ didefinisikan oleh
\begin{equation*}
f(x,y) \coloneqq 
\begin{cases}
1 & \text{jika } x \in \Q \text{ atau } y \in \Q, \\
0 & \text{selainnya.}
\end{cases}
\end{equation*}
Tunjukkan bahwa $f \notin \sR(R)$.
\end{enumerate}
\end{exercise}
\end{samepage}

\begin{exercise}
Misalkan $R \subset \R^n$ persegi panjang tertutup, dan misalkan
$\{ S_j \}_{j=1}^\infty$ barisan persegi panjang tertutup sedemikian sehingga
$S_j \subset R$ dan $S_j \subset S_{j+1}$ untuk setiap $j$, serta
$R = \bigcup_{j=1}^\infty S_j$. Misalkan $f \colon R \to \R$ terbatas dan
$f \in \sR(S_j)$ untuk setiap $j$. Tunjukkan bahwa $f \in \sR(R)$ dan
\begin{equation*}
\lim_{j\to\infty} \int_{S_j} f = \int_R f .
\end{equation*}
\end{exercise}

\begin{exercise}
Misalkan $f\colon [-1,1] \times [-1,1] \to \R$ fungsi yang terintegralkan
Riemann dan memenuhi $f(x) = -f(-x)$. Dengan menggunakan definisi integral,
buktikan bahwa
\begin{equation*}
\int_{[-1,1] \times [-1,1]} f = 0 .
\end{equation*}
\end{exercise}

%%%%%%%%%%%%%%%%%%%%%%%%%%%%%%%%%%%%%%%%%%%%%%%%%%%%%%%%%%%%%%%%%%%%%%%%%%%%%%

\sectionnewpage
\section{Integral berulang dan teorema Fubini}
\label{sec:iteratedints}

\sectionnotes{1--2 kuliah}

Integral Riemann beberapa peubah sulit dihitung langsung dari definisinya.
Untuk integral Riemann satu dimensi, kita memiliki teorema dasar kalkulus
(lihat \volIref{\sectionref*{vI-sec:ftc} di jilid I}{\sectionref{sec:ftc}}),
yang memungkinkan kita menghitung banyak integral tanpa harus menggunakan
definisi integral. Kita akan menulis ulang integral Riemann beberapa peubah
sebagai beberapa integral Riemann satu dimensi melalui pengintegralan
berulang. Namun, jika $f \colon [0,1]^2 \to \R$ fungsi terintegralkan Riemann,
belum langsung jelas apakah ketiga ungkapan
\begin{equation*}
\int_{[0,1]^2} f ,
\qquad
\int_0^1 \int_0^1 f(x,y) \, dx \, dy ,
\qquad \text{dan}
\qquad
\int_0^1 \int_0^1 f(x,y) \, dy \, dx
\end{equation*}
bernilai sama, atau bahkan apakah dua ungkapan terakhir terdefinisi.

\begin{example}
Definisikan
\begin{equation*}
f(x,y) \coloneqq 
\begin{cases}
1 & \text{jika } x=\nicefrac{1}{2} \text{ dan } y \in \Q, \\
0 & \text{selainnya.}
\end{cases}
\end{equation*}
Fungsi $f$ terintegralkan Riemann pada $R \coloneqq [0,1]^2$ dan
$\int_R f = 0$.

Selain itu, $\int_0^1 \int_0^1 f(x,y) \, dx \, dy = 0$. Akan tetapi,
\begin{equation*}
\int_0^1 f(\nicefrac{1}{2},y) \, dy
\end{equation*}
tidak ada. Jadi, secara ketat,
$\int_0^1 \int_0^1 f(x,y) \, dy \, dx$ tidak bermakna.
Lihat \figureref{fig:fubinibad}.

\begin{myfigureht}
\myincludepdft{fubinibad}{%
Tiga diagram berbentuk persegi. Di sebelah kiri terdapat persegi putih dengan
garis abu-abu mendatar di bagian atas dan garis vertikal bertitik-titik yang
melalui tengahnya. Diagram tengah memuat garis bertitik-titik yang sama, tetapi
garis abu-abunya kini vertikal dan terletak di bagian kiri persegi. Diagram
kanan memuat garis bertitik-titik yang sama, dengan dua garis vertikal
putus-putus, masing-masing sangat dekat di salah satu sisinya. Daerah di
sebelah kiri
garis putus-putus pertama diarsir gelap, daerah di sebelah kanan garis
putus-putus kedua diarsir abu-abu sedang, dan daerah tipis di sekitar garis
bertitik-titik diarsir dengan warna paling terang.}
\caption{Kiri: $[0,1]^2$ dengan garis $x=\nicefrac{1}{2}$ ditandai
bertitik-titik dan $\int_0^1 f(x,y) \, dx$ ditandai sebagai garis penuh abu-abu
untuk suatu $y$ tetap. Tengah: Diagram serupa, tetapi
$\int_0^1 f(x,y) \, dy$ ditandai untuk suatu $x \neq \nicefrac{1}{2}$.
Kanan: Tiga persegi panjang berbeda dalam partisi yang digunakan untuk
mengintegralkan $f$, masing-masing dengan corak abu-abu berbeda.\label{fig:fubinibad}}
\end{myfigureht}

Bukti:
Kita mulai dengan membuktikan keterintegralan $f$. Untuk
$0 < \epsilon < \nicefrac{1}{2}$, tinjau partisi $[0,1]^2$ dengan partisi
arah $x$ berupa $\{ 0, \nicefrac{1}{2}-\epsilon,
\nicefrac{1}{2}+\epsilon,1\}$ dan partisi arah $y$ berupa $\{ 0, 1 \}$.
Subpersegi panjang yang bersesuaian adalah
\begin{equation*}
R_1 \coloneqq [0, \nicefrac{1}{2}-\epsilon] \times [0,1],
%mbxSTARTIGNORE
\qquad
%mbxENDIGNORE
%mbxlatex \quad
R_2 \coloneqq [\nicefrac{1}{2}-\epsilon,
\nicefrac{1}{2}+\epsilon] \times [0,1],
%mbxSTARTIGNORE
\qquad
%mbxENDIGNORE
%mbxlatex \quad
R_3 \coloneqq [\nicefrac{1}{2}+\epsilon,1] \times [0,1] .
\end{equation*}
Kita memiliki $m_1 = M_1 = 0$, $m_2 =0$, $M_2 = 1$, dan
$m_3 = M_3 = 0$. Oleh karena itu,
\begin{equation*}
L(P,f) = 
m_1 V(R_1)
+
m_2 V(R_2)
+
m_3 V(R_3)
=
0 (\nicefrac{1}{2}-\epsilon)
+
0 (2\epsilon)
+
0 (\nicefrac{1}{2}-\epsilon) = 0 ,
\end{equation*}
dan
\begin{equation*}
U(P,f) = 
M_1 V(R_1)
+
M_2 V(R_2)
+
M_3 V(R_3)
=
0 (\nicefrac{1}{2}-\epsilon)
+
1 (2\epsilon)
+
0 (\nicefrac{1}{2}-\epsilon) = 2 \epsilon .
\end{equation*}
Jumlah atas dan jumlah bawah dapat dibuat sedekat apa pun, sedangkan jumlah
bawah selalu nol. Jadi, fungsi tersebut terintegralkan dan $\int_R f = 0$.

Untuk setiap $y$ tetap, fungsi yang memetakan $x$ ke $f(x,y)$ bernilai nol
kecuali mungkin pada satu titik, yaitu $x=\nicefrac{1}{2}$. Fungsi semacam ini
terintegralkan dan $\int_0^1 f(x,y) \, dx = 0$. Oleh karena itu,
$\int_0^1 \int_0^1 f(x,y) \, dx \, dy = 0$. Akan tetapi, jika
$x=\nicefrac{1}{2}$, fungsi yang memetakan $y$ ke
$f(\nicefrac{1}{2},y)$ merupakan fungsi tak terintegralkan yang bernilai 1
pada bilangan rasional dan 0 pada bilangan irasional. Lihat
\volIref{\exampleref*{vI-example:dirichletfunc} dari jilid I}{\exampleref{example:dirichletfunc}}.
\end{example}

Kita mengatasi masalah tidak terdefinisinya integral di bagian dalam dengan
menggunakan integral atas dan integral bawah, yang selalu terdefinisi untuk
setiap fungsi terbatas.

\medskip

Kelompokkan koordinat pada $\R^{n+m}$ menjadi dua bagian. Tuliskan koordinat pada
$\R^{n+m} = \R^n \times \R^m$ sebagai $(x,y)$, dengan $x \in \R^n$ dan
$y \in \R^m$. Untuk fungsi $f(x,y)$, tulis
\begin{equation*}
f_x(y) \coloneqq f(x,y)
\end{equation*}
jika $x$ ditetapkan dan kita hendak memandangnya sebagai fungsi dari $y$.
Tulis
\begin{equation*}
f^y(x) \coloneqq f(x,y)
\end{equation*}
jika $y$ ditetapkan dan kita hendak memandangnya sebagai fungsi dari $x$.

\begin{thm}[Fubini versi A%
\footnote{Dinamai menurut matematikawan Italia
\href{https://en.wikipedia.org/wiki/Guido_Fubini}{Guido Fubini}
(1879--1943).}] \label{mv:fubinivA} \index{teorema Fubini}
Misalkan $R \times S \subset \R^n \times \R^m$ merupakan persegi panjang
tertutup dan $f \colon R \times S \to \R$ fungsi yang terintegralkan. Fungsi
$g \colon R \to \R$ dan $h \colon R \to \R$ yang didefinisikan oleh
\begin{equation*}
g(x) \coloneqq \underline{\int_S} f_x \qquad
\text{dan} \qquad
h(x) \coloneqq \overline{\int_S} f_x 
\end{equation*}
terintegralkan pada $R$ dan
\begin{equation*}
\int_R g = \int_R h = \int_{R \times S} f .
\end{equation*}
\end{thm}

Dengan kata lain,
\begin{equation*}
\int_{R \times S} f
=
 \int_R \left(
 \underline{\int_S} f(x,y) \, dy
\right) \, dx
=
 \int_R \left(
 \overline{\int_S} f(x,y) \, dy
\right) \, dx .
\end{equation*}
Jika $f_x$ terintegralkan untuk setiap $x \in R$, misalnya ketika $f$ kontinu, kita
memperoleh bentuk yang lebih lazim
\begin{equation*}
\int_{R \times S} f
=
 \int_R \int_S f(x,y) \, dy \, dx .
\end{equation*}

\begin{proof}
Partisi dari $R \times S$ dapat ditulis sebagai
$(P,P') = (P_1,P_2,\ldots,P_n,P'_1,P'_2,\ldots,P'_m)$, dengan
$P = (P_1,P_2,\ldots,P_n)$ dan
$P' = (P'_1,P'_2,\ldots,P'_m)$ berturut-turut merupakan partisi dari $R$ dan
$S$. Misalkan $R_1,R_2,\ldots,R_N$ subpersegi panjang dari $P$, sedangkan
$R'_1,R'_2,\ldots,R'_K$ subpersegi panjang dari $P'$. Subpersegi panjang dari
$(P,P')$ adalah $R_i \times R'_j$, dengan $1 \leq i \leq N$ dan
$1 \leq j \leq K$.

Misalkan
\begin{equation*}
m_{i,j} \coloneqq
\inf_{(x,y) \in R_i \times R'_j} f(x,y) .
\end{equation*}
Perhatikan bahwa $V(R_i \times R'_j) = V(R_i)V(R'_j)$, sehingga
\begin{equation*}
L\bigl((P,P'),f\bigr) =
\sum_{i=1}^N
\sum_{j=1}^K
m_{i,j} \, V(R_i \times R'_j)
=
\sum_{i=1}^N
\left(
\sum_{j=1}^K
m_{i,j} \, V(R'_j) \right) V(R_i) .
\end{equation*}
Definisikan
\begin{equation*}
m_j(x) \coloneqq \inf_{y \in R'_j} f(x,y) = \inf_{y \in R'_j} f_x(y) .
\end{equation*}
Untuk $x \in R_i$, berlaku $m_{i,j} \leq m_j(x)$. Oleh karena itu,
\begin{equation*}
\sum_{j=1}^K
m_{i,j} \, V(R'_j)
\leq \sum_{j=1}^K m_j(x) \, V(R'_j) = L(P',f_x) \leq
\underline{\int_S} f_x = g(x) .
\end{equation*}
Ketaksamaan ini berlaku untuk setiap $x \in R_i$, sehingga
\begin{equation*}
\sum_{j=1}^K
m_{i,j} \, V(R'_j)
\leq \inf_{x \in R_i} g(x) .
\end{equation*}
Dengan demikian, kita memperoleh
\begin{equation*}
L\bigl((P,P'),f\bigr) 
\leq
\sum_{i=1}^N
\left(
\inf_{x \in R_i} g(x)
\right) V(R_i) = L(P,g) .
\end{equation*}

Dengan cara serupa, $U\bigl((P,P'),f\bigr) \geq U(P,h)$; pembuktian
ketaksamaan ini diserahkan kepada pembaca sebagai latihan. Dengan menggabungkan
kedua ketaksamaan tersebut dan fakta bahwa $g(x) \leq h(x)$ untuk setiap
$x \in R$, kita memperoleh
\begin{equation*}
L\bigl((P,P'),f\bigr)
\leq
L(P,g) \leq
U(P,g) \leq
U(P,h) \leq
U\bigl((P,P'),f\bigr) .
\end{equation*}
Karena $f$ terintegralkan, untuk setiap $\eta > 0$ terdapat pasangan partisi
$P$ dari $R$ dan $P'$ dari $S$ sedemikian sehingga
$U\bigl((P,P'),f\bigr)-L\bigl((P,P'),f\bigr)<\eta$. Untuk pasangan partisi
tersebut,
\begin{equation*}
U(P,g) - L(P,g)
\leq
U\bigl((P,P'),f\bigr) -
L\bigl((P,P'),f\bigr) < \eta .
\end{equation*}
Jadi, menurut kriteria keterintegralan, $g$ terintegralkan. Rangkaian
ketaksamaan sebelumnya menempatkan baik $\int_R g$ maupun
$\int_{R \times S} f$ di antara jumlah bawah dan jumlah atas untuk $f$ yang
selisihnya kurang dari $\eta$. Karena $\eta>0$ sebarang,
$\int_R g = \int_{R \times S} f$.

Demikian pula,
\begin{equation*}
L\bigl((P,P'),f\bigr)
\leq
L(P,g) \leq
L(P,h) \leq
U(P,h) \leq
U\bigl((P,P'),f\bigr) ,
\end{equation*}
sehingga
\begin{equation*}
U(P,h) - L(P,h)
\leq
U\bigl((P,P'),f\bigr) -
L\bigl((P,P'),f\bigr) .
\end{equation*}
Dengan memilih pasangan partisi di atas untuk setiap $\eta>0$, kriteria
keterintegralan menunjukkan bahwa $h$ terintegralkan. Selain itu,
rangkaian ketaksamaan terakhir menempatkan baik $\int_R h$ maupun
$\int_{R \times S} f$ di antara dua jumlah untuk $f$ yang selisihnya kurang
dari $\eta$. Jadi, seperti sebelumnya,
$\int_R h = \int_{R \times S} f$.
\end{proof}

Kita juga dapat melakukan pengintegralan berulang dalam urutan sebaliknya.
Bukti versi ini hampir sama dengan bukti versi A
(atau dapat diperoleh dengan cepat dari versi A)\@.
Pembuktiannya diserahkan kepada pembaca sebagai latihan.

\begin{thm}[Fubini versi B]\label{mv:fubinivB}\index{teorema Fubini}
\pagebreak[2]
Misalkan $R \times S \subset \R^n \times \R^m$ merupakan persegi panjang
tertutup dan $f \colon R \times S \to \R$ fungsi yang terintegralkan.
Fungsi $g \colon S \to \R$ dan $h \colon S \to \R$ yang didefinisikan oleh
\begin{equation*}
g(y) \coloneqq \underline{\int_R} f^y \qquad
\text{dan} \qquad
h(y) \coloneqq \overline{\int_R} f^y 
\end{equation*}
terintegralkan pada $S$, dan
\begin{equation*}
\int_S g = \int_S h = \int_{R \times S} f .
\end{equation*}
\end{thm}

Dengan kata lain,
\begin{equation*}
\int_{R \times S} f
=
 \int_S \left(
 \underline{\int_R} f(x,y) \, dx
\right) \, dy
=
 \int_S \left(
 \overline{\int_R} f(x,y) \, dx
\right) \, dy .
\end{equation*}

Selanjutnya, dengan tetap mengandaikan bahwa $f$ terintegralkan pada
$R \times S$, misalkan $f_x$ dan $f^y$ terintegralkan untuk setiap $x \in R$
dan $y \in S$. Misalnya, anggaplah $f$ kontinu. Dengan menggabungkan kedua
versi tersebut, kita memperoleh rumus yang lazim
\begin{equation*}
\int_{R \times S} f
=
 \int_R 
 \int_S f(x,y) \, dy \, dx 
=
 \int_S 
 \int_R f(x,y) \, dx \, dy .
\end{equation*}

Teorema Fubini sering dinyatakan dalam dua dimensi untuk fungsi kontinu
$f \colon R \to \R$ pada persegi panjang
$R = [a,b] \times [c,d]$. Dalam hal ini, teorema Fubini menyatakan bahwa
\begin{equation*}
\int_R f = \int_a^b \int_c^d f(x,y) \,dy\,dx
=
\int_c^d \int_a^b f(x,y) \,dx\,dy .
\end{equation*}
Teorema Fubini lazim dipandang sebagai \emph{teorema} yang memungkinkan kita
menukar urutan integral berulang, meskipun teorema Fubini memiliki banyak
variasi dan baru dua di antaranya yang telah kita lihat.

Dengan menerapkan teorema Fubini berulang kali, kita memperoleh akibat berikut.
Misalkan $R \coloneqq [a_1,b_1] \times [a_2,b_2] \times \cdots \times [a_n,b_n]
\subset \R^n$ merupakan persegi panjang tertutup dan
$f \colon R \to \R$ fungsi kontinu. Maka
\begin{equation*}
\int_R f = 
\int_{a_1}^{b_1}
\int_{a_2}^{b_2}
\cdots
\int_{a_n}^{b_n}
f(x_1,x_2,\ldots,x_n)
\,
dx_n
\,
dx_{n-1}
\cdots
dx_1 .
\end{equation*}

Kita boleh mengubah urutan pengintegralan ke urutan mana pun yang dikehendaki.
Dengan tetap mensyaratkan bahwa $f$ terintegralkan pada $R$, syarat kekontinuan
dapat dilonggarkan dengan memastikan bahwa semua fungsi perantara yang muncul
dalam urutan pilihan terintegralkan. Sebagai alternatif, integral atas atau
integral bawah dapat digunakan secara tepat.

\subsection{Latihan}

\begin{exercise}
Hitung $\int_{0}^1 \int_{-1}^1 xe^{xy} \, dx \, dy$ dengan cara sederhana.
\end{exercise}

\begin{exercise}
Buktikan pernyataan $U\bigl((P,P'),f\bigr) \geq U(P,h)$ dari bukti
\thmref{mv:fubinivA}.
\end{exercise}

\begin{exercise}[Mudah]
Buktikan \thmref{mv:fubinivB}.
\end{exercise}

\begin{exercise}
Misalkan $R \coloneqq [a,b] \times [c,d]$ dan
$f \colon R \to \R$ fungsi yang terintegralkan sedemikian sehingga, untuk
setiap $y \in [c,d]$ yang ditetapkan, fungsi yang memetakan $x$ ke $f(x,y)$
bernilai nol kecuali pada berhingga banyak titik. Tunjukkan
\begin{equation*}
\int_R f = 0 .
\end{equation*}
\end{exercise}

\begin{exercise}
Misalkan $R \coloneqq [a,b] \times [c,d]$ dan
$f(x,y) \coloneqq g(x)h(y)$ untuk fungsi kontinu
$g \colon [a,b] \to \R$ dan $h \colon [c,d] \to \R$. Buktikan
\begin{equation*}
\int_R f = \left(\int_a^b g\right)\left(\int_c^d h\right) .
\end{equation*}
\end{exercise}

\begin{exercise}
Hitunglah, dengan menggunakan kalkulus,
\begin{equation*}
\int_0^1 \int_0^1 \frac{x^2-y^2}{{(x^2+y^2)}^2} \, dx \, dy
\qquad \text{dan} \qquad
\int_0^1 \int_0^1 \frac{x^2-y^2}{{(x^2+y^2)}^2} \, dy \, dx .
\end{equation*}
Dalam setiap integral berulang tersebut, hitung dahulu integral dalam untuk
nilai positif variabel luar, lalu tafsirkan integral luar sebagai integral tak
wajar $\lim_{\epsilon \to 0^+}\int_\epsilon^1$.
\end{exercise}

\begin{exercise}
Misalkan $f(x,y) \coloneqq g(x)$, dengan $g \colon [a,b] \to \R$ terintegralkan
Riemann. Tunjukkan bahwa $f$ terintegralkan Riemann untuk setiap
$R = [a,b] \times [c,d]$ dan
\begin{equation*}
\int_R f = (d-c) \int_a^b g .
\end{equation*}
\end{exercise}

\begin{exercise}
Definisikan $f \colon [-1,1] \times [0,1] \to \R$ dengan
\begin{equation*}
f(x,y) \coloneqq
\begin{cases}
x & \text{jika } y \in \Q, \\
0 & \text{selain itu.} 
\end{cases}
\end{equation*}
\begin{enumerate}[a)]
\item
Tunjukkan bahwa
$\int_0^1 \int_{-1}^1 f(x,y) \, dx \, dy$ ada, tetapi
$\int_{-1}^1 \int_0^1 f(x,y) \, dy \, dx$ tidak ada.
\item
Hitung
$\int_{-1}^1 \overline{\int_0^1} f(x,y) \, dy \, dx$ dan
$\int_{-1}^1 \underline{\int_0^1} f(x,y) \, dy \, dx$.
\item
Tunjukkan bahwa $f$ tidak terintegralkan Riemann pada $[-1,1] \times [0,1]$
(gunakan Fubini).
\end{enumerate}
\end{exercise}

\begin{exercise}
Definisikan $f \colon [0,1] \times [0,1] \to \R$ dengan ketentuan bahwa
representasi $y=\nicefrac{p}{q}$ di bawah selalu memenuhi
$p\in\Z$, $q\in\N$, $q>0$, dan $\gcd(|p|,q)=1$ (khususnya, $0=\nicefrac{0}{1}$):
\begin{equation*}
f(x,y) \coloneqq
\begin{cases}
\nicefrac{1}{q} & \text{jika } x \in \Q, y \in \Q, \text{ dan }
  y=\nicefrac{p}{q} \text{ dalam bentuk pecahan paling sederhana,} \\
0               & \text{selain itu.} 
\end{cases}
\end{equation*}
\begin{enumerate}[a)]
\item
Tunjukkan bahwa $f$ terintegralkan Riemann pada $[0,1] \times [0,1]$.
\item
Tentukan
$\overline{\int_0^1} f(x,y) \, dx$ dan
$\underline{\int_0^1} f(x,y) \, dx$ untuk setiap $y \in [0,1]$, dan tunjukkan
bahwa keduanya tidak sama untuk setiap $y \in \Q \cap [0,1]$.
\item
Tunjukkan bahwa
$\int_0^1 \int_0^1 f(x,y) \, dy \, dx$ ada, tetapi
$\int_0^1 \int_0^1 f(x,y) \, dx \, dy$ tidak ada.
\end{enumerate}
Catatan: Menurut Fubini,
$\int_0^1 \overline{\int_0^1} f(x,y) \, dx \, dy$ dan
$\int_0^1 \underline{\int_0^1} f(x,y) \, dx \, dy$ ada dan sama dengan
$\int_{[0,1]\times[0,1]} f$.
\end{exercise}

%%%%%%%%%%%%%%%%%%%%%%%%%%%%%%%%%%%%%%%%%%%%%%%%%%%%%%%%%%%%%%%%%%%%%%%%%%%%%%

\sectionnewpage
\section{Ukuran luar dan himpunan nol}
\label{sec:outermeasure}

%mbxINTROSUBSECTION

\sectionnotes{2 kuliah}

\subsection{Ukuran luar dan himpunan nol}

Sebelum mencirikan semua fungsi yang terintegralkan Riemann, kita perlu sedikit
menyimpang dari pembahasan utama. Kita akan memperkenalkan cara mengukur besar
suatu himpunan di $\R^n$.

\begin{defn}
Kita definisikan \emph{\myindex{ukuran luar}} suatu himpunan
$S \subset \R^n$ sebagai
\glsadd{not:outermeasure}
\begin{equation*}
m^*(S)
\coloneqq
\inf\,
\sum_{j=1}^\infty V(R_j) ,
\end{equation*}
di mana infimum diambil atas semua barisan persegi panjang terbuka
$\{ R_j \}_{j=1}^\infty$ yang memenuhi
$S \subset \bigcup_{j=1}^\infty R_j$. Baik jumlah maupun infimum diperbolehkan
bernilai $\infty$. Lihat \figureref{fig:outermeasure}.
Secara khusus, $S$ disebut \emph{\myindex{berukuran nol}} atau
\emph{\myindex{himpunan nol}} jika $m^*(S) = 0$.
\end{defn}

\begin{myfigureht}
\myincludepdft{outermeasure}{Himpunan S diarsir abu-abu dan ditutupi oleh
banyak persegi panjang yang dibatasi garis tipis bertitik-titik. Tiga persegi
panjang diberi tanda R sub 1, R sub 2, dan R sub 3.}
\caption{Konstruksi ukuran luar; dalam kasus ini $S \subset R_1 \cup R_2
\cup R_3 \cup \cdots$, sehingga $m^*(S) \leq V(R_1) + V(R_2)+V(R_3) +
\cdots$.\label{fig:outermeasure}}
\end{myfigureht}

Salah satu konsekuensi langsung definisi tersebut
(\exerciseref{exercise:outermeasuremono}) adalah bahwa jika $A \subset B$,
maka $m^*(A) \leq m^*(B)$. Tidak sulit pula untuk menunjukkan
(\exerciseref{exercise:allowfiniteseqsinoutermeasure}) bahwa kita memperoleh
bilangan $m^*(S)$ yang sama jika keluarga persegi panjang berhingga maupun
barisan tak hingga diizinkan dalam definisi. Namun, keluarga berhingga saja
tidak cukup.

Teori ukuran pada $\R^n$ merupakan bidang yang sangat rumit. Kita hanya akan
memerlukan himpunan berukuran nol, sehingga pembahasan difokuskan pada himpunan
tersebut. Himpunan $S$ berukuran nol jika untuk setiap $\epsilon > 0$ terdapat
barisan persegi panjang terbuka $\{ R_j \}_{j=1}^\infty$ sedemikian sehingga
\begin{equation} \label{mv:eq:nullR}
S \subset \bigcup_{j=1}^\infty R_j \qquad \text{dan} \qquad
\sum_{j=1}^\infty V(R_j) < \epsilon.
\end{equation}
Jika $S$ berukuran nol dan $S' \subset S$, maka $S'$ berukuran nol. Kita dapat
menggunakan persegi panjang yang sama.

\pagebreak[2]
Kadang-kadang lebih mudah menggunakan bola daripada persegi panjang. Selain
itu, kita dapat memilih bola yang jari-jarinya lebih kecil daripada suatu batas
tetap.

\begin{prop} \label{mv:prop:ballsnull}
Misalkan $\delta > 0$ diberikan. Himpunan $S \subset \R^n$ berukuran nol jika
dan hanya jika, untuk setiap $\epsilon > 0$, terdapat barisan bola terbuka
$\{ B_k \}_{k=1}^\infty$, dengan $B_k$ berjari-jari $r_k < \delta$,
sedemikian sehingga
\begin{equation*}
S \subset \bigcup_{k=1}^\infty B_k \qquad \text{dan} \qquad
\sum_{k=1}^\infty r_k^n < \epsilon.
\end{equation*}
\end{prop}

Perhatikan bahwa \myquote{volume} $B_k$ sebanding dengan $r_k^n$.

\begin{proof}
Jika $C$ merupakan kubus tertutup (persegi panjang yang semua panjang sisinya
sama) dengan panjang sisi $s$, maka menurut
\propref{prop:diameterrectangle}, $C$ termuat dalam bola tertutup berjari-jari
$\sqrt{n}\,s$, sehingga juga termuat dalam bola terbuka berjari-jari
$2\sqrt{n}\,s$.\footnote{%
Bola tertutup berjari-jari $\frac{\sqrt{n}\,s}{2}$ sebenarnya dapat digunakan,
tetapi hal itu memerlukan argumen tambahan, sedangkan $\sqrt{n}\,s$ sudah
memadai.}

Misalkan $R$ merupakan persegi panjang bervolume positif. Pilih bilangan
$s>0$ yang lebih kecil daripada panjang sisi terpendek $R$ dan memenuhi
$2\sqrt{n}\,s<\delta$. Jika setiap panjang sisi $R$ merupakan kelipatan bulat
dari $s$, maka $R$ termuat dalam gabungan kubus tertutup
$C_1,C_2,\ldots,C_m$ yang panjang sisinya $s$ dan memenuhi
$\sum_{k=1}^m V(C_k)=V(R)$. Sekarang, andaikan tidak semua panjang sisi $R$
merupakan kelipatan bulat dari $s$. Tinjau suatu sisi yang panjangnya
$(\ell+\alpha)s$, dengan $\ell$ bilangan bulat dan $0\leq\alpha<1$. Karena
$s$ lebih kecil daripada panjang sisi terpendek, $\ell\geq1$, sehingga
$(\ell+\alpha)s\leq2\ell s$. Dengan memperpanjang sisi tersebut menjadi
$2\ell s$, dan memperpanjang setiap sisi $R$ dengan cara serupa, kita
memperoleh persegi panjang baru yang lebih besar, yang volumenya paling besar
$2^n$ kali volume semula dan panjang sisi-sisinya merupakan kelipatan dari
$s$. Lihat \figureref{fig:nullrectcube}. Jadi, $R$ termuat dalam gabungan
kubus tertutup $C_1,C_2,\ldots,C_m$ yang panjang sisinya $s$ dan memenuhi
\begin{equation*}
\sum_{k=1}^m V(C_k) \leq 2^n V(R) .
\end{equation*}

\begin{myfigureht}
\myincludepdft{nullrectcube}{%
Sebuah persegi panjang ditandai dengan garis tebal. Kisi persegi yang panjang
sisinya s, terdiri dari 2 baris dan 4 kolom, disejajarkan dengan sudut kiri atas
persegi panjang bergaris tebal tersebut. Persegi panjang bergaris tebal
memanjang sedikit lebih dari 2 s ke kanan dan sedikit lebih dari s ke bawah.
Dua petak pertama pada arah mendatar diberi tanda ell s sama dengan 2 s,
sedangkan keseluruhan empat petak pada arah mendatar diberi tanda 2 ell s sama
dengan 4 s.}
\caption{Menutupi persegi panjang dengan kubus-kubus yang volume totalnya
paling besar $2^n V(R)$.\label{fig:nullrectcube}}
\end{myfigureht}

Sekarang, misalkan $S$ merupakan himpunan nol dan terdapat barisan persegi
panjang terbuka $\{R_j\}_{j=1}^\infty$ yang gabungannya memuat $S$ dan memenuhi
\eqref{mv:eq:nullR}. Pilih barisan kubus tertutup
$\{C_k\}_{k=1}^\infty$, dengan $C_k$ mempunyai panjang sisi $s_k$, seperti di
atas, yang menutupi semua persegi panjang $\{R_j\}_{j=1}^\infty$ dan memenuhi
\begin{equation*}
\sum_{k=1}^\infty s_k^n =
\sum_{k=1}^\infty V(C_k) \leq
2^n \sum_{j=1}^\infty V(R_j)
< 2^n \epsilon.
\end{equation*}
Dengan menutupi setiap $C_k$ menggunakan bola $B_k$ berjari-jari
$r_k=2\sqrt{n}\,s_k<\delta$, kita memperoleh
\begin{equation*}
\sum_{k=1}^\infty r_k^n
=
\sum_{k=1}^\infty {\bigl(2\sqrt{n}\bigr)}^n s_k^n
<
{\bigl(4\sqrt{n}\bigr)}^n \epsilon .
\end{equation*}
Karena $S \subset\bigcup_j R_j \subset\bigcup_k C_k \subset\bigcup_k B_k$
dan ${\bigl(4\sqrt{n}\bigr)}^n\epsilon$ dapat dibuat sekecil apa pun, arah maju
terbukti.

Untuk arah sebaliknya, misalkan $S$ ditutupi oleh bola $B_j$ berjari-jari
$r_j$ sedemikian sehingga $\sum_{j=1}^\infty r_j^n<\epsilon$, seperti dalam
pernyataan proposisi. Setiap $B_j$ termuat dalam kubus terbuka $R_j$ yang
panjang sisinya $2r_j$. Jadi, $V(R_j)={(2r_j)}^n=2^n r_j^n$. Oleh karena itu,
\begin{equation*}
S \subset \bigcup_{j=1}^\infty R_j \qquad \text{dan} \qquad
\sum_{j=1}^\infty V(R_j)
\leq
\sum_{j=1}^\infty 2^n r_j^n < 2^n \epsilon. \qedhere
\end{equation*}
\end{proof}

Definisi ukuran luar (bukan hanya himpunan nol) juga dapat dirumuskan dengan
bola terbuka. Generalisasi ini diserahkan kepada pembaca.

\subsection{Contoh dan sifat dasar}

\begin{example}
Himpunan $\Q^n \subset \R^n$ yang terdiri atas titik-titik berkoordinat
rasional berukuran nol.

Bukti:
Himpunan $\Q^n$ terhitung, sehingga dapat ditulis sebagai barisan
$q_1,q_2,\ldots$. Tetapkan sembarang $\epsilon>0$. Untuk setiap $q_j$, pilih
persegi panjang terbuka $R_j$ yang memenuhi $q_j \in R_j$ dan
$V(R_j)<\epsilon 2^{-j}$. Dengan demikian,
\begin{equation*}
\Q^n \subset \bigcup_{j=1}^\infty R_j \qquad \text{dan} \qquad
\sum_{j=1}^\infty V(R_j) <
\sum_{j=1}^\infty \epsilon 2^{-j} = \epsilon .
\end{equation*}
Karena $\epsilon$ sembarang, $\Q^n$ berukuran nol.
\end{example}

Contoh tersebut mengarah pada hasil yang lebih umum.

\begin{prop}
Gabungan terhitung dari himpunan-himpunan berukuran nol juga berukuran nol.
\end{prop}

\begin{proof}
Misalkan
\begin{equation*}
S = \bigcup_{j=1}^\infty S_j ,
\end{equation*}
di mana setiap $S_j$ berukuran nol. Misalkan $\epsilon>0$ diberikan. Untuk
setiap $j$, terdapat barisan persegi panjang terbuka
$\{R_{j,k}\}_{k=1}^\infty$ sedemikian sehingga
\begin{equation*}
S_j \subset \bigcup_{k=1}^\infty R_{j,k}
\qquad \text{dan} \qquad
\sum_{k=1}^\infty V(R_{j,k}) < 2^{-j} \epsilon .
\end{equation*}
Maka
\begin{equation*}
S \subset \bigcup_{j=1}^\infty \bigcup_{k=1}^\infty R_{j,k} .
\end{equation*}
Karena $\N\times\N$ terhitung, keluarga berindeks ganda
$\{R_{j,k}\}_{j,k=1}^\infty$ dapat didaftarkan sebagai satu barisan persegi
panjang. Semua $V(R_{j,k})$ tak negatif, sehingga penjumlahan terhadap semua
indeks $j$ dan $k$ dapat dilakukan dengan terlebih dahulu menjumlahkan terhadap
$k$, kemudian terhadap $j$; lihat
\volIref{\exerciseref*{vI-exercise:tonellifubiniforsums} dalam jilid I}{\exerciseref{exercise:tonellifubiniforsums}}.
Secara khusus,
\begin{equation*}
\sum_{j=1}^\infty \sum_{k=1}^\infty V(R_{j,k}) <
\sum_{j=1}^\infty 2^{-j} \epsilon = \epsilon . \qedhere
\end{equation*}
\end{proof}

Contoh berikut bukan hanya menarik, tetapi juga akan berguna nanti.

\begin{example} \label{mv:example:planenull}
Misalkan $n \in \N$, $k=1,2,\ldots,n$, dan $c \in \R$.
Maka $P \coloneqq \{x \in \R^n : x_k=c\}$ berukuran nol. Perhatikan bahwa
jika $n\geq2$, maka $P$ tidak terhitung.

Bukti:
Pertama, tetapkan $s \in \N$ dan tinjau
\begin{equation*}
P_s \coloneqq \bigl\{x \in \R^n : x_k=c \text{ dan } \sabs{x_j}\leq s
\text{ untuk setiap } j\neq k\bigr\} .
\end{equation*}
Untuk sembarang $\eta>0$, pilih $\epsilon>0$ sedemikian sehingga
$2\epsilon{\bigl(2(s+1)\bigr)}^{n-1}<\eta$, lalu definisikan persegi panjang
terbuka
\begin{equation*}
R \coloneqq \bigl\{x \in \R^n : c-\epsilon<x_k<c+\epsilon \text{ dan }
\sabs{x_j}<s+1 \text{ untuk setiap } j\neq k\bigr\} .
\end{equation*}
Jelas bahwa $P_s\subset R$. Selain itu,
\begin{equation*}
V(R) = 2\epsilon {\bigl(2(s+1)\bigr)}^{n-1} .
\end{equation*}
Berdasarkan pilihan $\epsilon$, $V(R)<\eta$. Jadi, $P_s$ berukuran nol.

Selanjutnya,
\begin{equation*}
P = \bigcup_{s=1}^\infty P_s
\end{equation*}
Karena setiap $P_s$ berukuran nol, proposisi sebelumnya menunjukkan bahwa $P$
berukuran nol.
\end{example}

\begin{example}
Jika $a < b$, maka $m^*\bigl([a,b]\bigr) = b-a$.

Bukti:
Di $\R$, persegi panjang terbuka adalah interval terbuka. Karena
$[a,b]\subset(a-\epsilon,b+\epsilon)$ untuk setiap $\epsilon>0$, dan panjang
interval ini adalah $b-a+2\epsilon$, diperoleh
$m^*\bigl([a,b]\bigr)\leq b-a$.

Ketaksamaan sebaliknya lebih sulit. Misalkan
$\bigl\{(a_j,b_j)\bigr\}_{j=1}^\infty$ adalah interval-interval terbuka
sedemikian sehingga
\begin{equation*}
[a,b] \subset \bigcup_{j=1}^\infty (a_j,b_j) .
\end{equation*}
Kita hendak memberikan batas bawah bagi
$\sum_{j=1}^\infty(b_j-a_j)$. Karena $[a,b]$ kompak, suatu subselimut hingga
dari interval-interval terbuka tersebut masih menutupi $[a,b]$. Membuang
sebagian interval hanya memperkecil jumlah, sehingga cukup ditinjau suatu
subselimut hingga. Buang satu demi satu setiap interval yang dapat dihilangkan
tanpa kehilangan sifat menutupi $[a,b]$; dengan demikian, tersisa subselimut
hingga yang minimal.

Tidak ada interval dalam subselimut minimal ini yang termuat dalam interval lain,
dan setiap interval memotong $[a,b]$. Beri nomor ulang interval-interval yang
tersisa sehingga
$a_1<a_2<\cdots<a_k$. Titik-titik ujung kanannya juga memenuhi
$b_1<b_2<\cdots<b_k$; jika tidak, salah satu interval akan termuat di interval
lain. Interval-interval yang berurutan harus bertumpang tindih, sebab jika
$b_j\leq a_{j+1}$, suatu titik di $[a,b]$ tidak akan tertutupi. Jadi,
$a_j<a_{j+1}<b_j<b_{j+1}$ untuk setiap $j=1,2,\ldots,k-1$. Selain itu,
$a_1<a$ dan $b_k>b$. Lihat \figureref{fig:measureinterval} untuk contoh
konfigurasi. Karena $b_j-a_j>a_{j+1}-a_j$, diperoleh
\begin{equation*}
\sum_{j=1}^k (b_j-a_j)
\geq
\sum_{j=1}^{k-1} (a_{j+1}-a_j)
+
(b_k-a_k)
=
b_k-a_1 > b-a .
\end{equation*}
Jadi, $m^*\bigl([a,b]\bigr)\geq b-a$.
\begin{myfigureht}
\myincludepdft{measureinterval}{%
Diagram interval dari a hingga b pada garis bilangan, yang ditandai dengan
ruas garis tebal gelap. Empat interval terbuka lain yang saling bertumpang
tindih menutupi interval utama. Titik-titik ujung kirinya berlabel a subskrip 1
hingga a subskrip 4, dan titik-titik ujung kanannya berlabel b subskrip 1 hingga
b subskrip 4. Urutan semua titik dari kiri ke kanan adalah a subskrip 1, a,
a subskrip 2, b subskrip 1, a subskrip 3, a subskrip 4, b subskrip 2,
b subskrip 3, b, dan b subskrip 4.}
\caption{Interval-interval terbuka yang menutupi $[a,b]$ dan memenuhi
$a_j<a_{j+1}<b_j<b_{j+1}$ untuk $j=1,2,\ldots,k-1$.
\label{fig:measureinterval}}
\end{myfigureht}
\end{example}

\begin{prop} \label{mv:prop:compactnull}
Misalkan $E\subset\R^n$ merupakan himpunan kompak berukuran nol. Maka, untuk
setiap $\epsilon>0$, terdapat berhingga banyak persegi panjang terbuka
$R_1,R_2,\ldots,R_k$ sedemikian sehingga
\begin{equation*}
E \subset R_1 \cup R_2 \cup \cdots \cup R_k
\qquad \text{dan} \qquad
\sum_{j=1}^k V(R_j) < \epsilon.
\end{equation*}
Selain itu, untuk setiap $\epsilon>0$ dan setiap $\delta>0$, terdapat berhingga
banyak bola terbuka $B_1,B_2,\ldots,B_{\ell}$, masing-masing berjari-jari
$r_j<\delta$, $j=1,2,\ldots,\ell$, sedemikian sehingga
\begin{equation*}
E \subset B_1 \cup B_2 \cup \cdots \cup B_{\ell}
\qquad \text{dan} \qquad
\sum_{j=1}^{\ell} r_j^n < \epsilon.
\end{equation*}
\end{prop}

\begin{proof}
Karena $E$ berukuran nol, terdapat barisan persegi panjang terbuka
$\{R_j\}_{j=1}^\infty$ sedemikian sehingga
\begin{equation*}
E \subset \bigcup_{j=1}^\infty R_j
\qquad \text{dan} \qquad
\sum_{j=1}^\infty V(R_j) < \epsilon.
\end{equation*}
Berdasarkan kekompakan $E$, berhingga banyak persegi panjang ini masih menutupi
$E$.
Setelah mengubah penomoran jika perlu, terdapat $k$ sedemikian sehingga
$E\subset R_1\cup R_2\cup\cdots\cup R_k$. Oleh karena itu,
\begin{equation*}
\sum_{j=1}^k V(R_j) \leq
\sum_{j=1}^\infty V(R_j) < \epsilon.
\end{equation*}

Bukti bahwa kita dapat memilih bola sebagai pengganti persegi panjang
diserahkan kepada pembaca sebagai latihan.
\end{proof}

\begin{example} \label{example:cantor}
Agar pembaca tidak mengira bahwa hanya ada sedikit himpunan berukuran nol dan
bahwa semuanya sederhana, berikut adalah subhimpunan $[0,1]$ yang tak
terhitung, kompak, berukuran nol, dan tidak memuat interval apa pun yang
panjangnya positif. Setiap $x\in[0,1]$ memiliki ekspansi terner:
\begin{equation*}
x = \sum_{n=1}^\infty d_n 3^{-n} ,
\qquad \text{dengan } d_n=0, 1, \text{ atau } 2.
\end{equation*}
Lihat \volIref{\sectionref*{vI-sec:decimals} dalam jilid I\@, khususnya
\exerciseref*{vI-exercise:decimalpropbaseb}}{\sectionref{sec:decimals},
khususnya \exerciseref{exercise:decimalpropbaseb}}. Definisikan
\emph{\myindex{himpunan Cantor}} $C$ sebagai
\begin{equation*}
C \coloneqq \Bigl\{ x \in [0,1] : x = \sum_{n=1}^\infty d_n 3^{-n},
\text{ dengan } d_n = 0 \text{ atau } d_n = 2
\text{ untuk setiap } n \Bigr\} .
\end{equation*}
Dengan kata lain, $x$ berada di $C$ jika memiliki ekspansi terner yang hanya
menggunakan digit $0$ dan $2$. Jika $x$ memiliki dua ekspansi, maka $x$ berada
di $C$ asalkan salah satunya tidak mengandung digit $1$. Definisikan
$C_0\coloneqq[0,1]$ dan
\begin{equation*}
C_k \coloneqq \Bigl\{ x \in [0,1] : x = \sum_{n=1}^\infty d_n 3^{-n},
\text{ dengan } d_n = 0 \text{ atau } d_n = 2
\text{ untuk setiap } n\in\{1,2,\ldots,k\} \Bigr\} .
\end{equation*}
Setiap $x\in C$ jelas berada di setiap $C_k$. Sebaliknya, misalkan $x$ berada
di setiap $C_k$. Setiap $x\in[0,1]$ memiliki paling banyak dua ekspansi terner.
Jika setiap ekspansi terner $x$ mengandung digit $1$, pada masing-masing
ekspansi terdapat posisi pertama tempat digit $1$ muncul. Pilih $k$ yang
sekurang-kurangnya sebesar semua posisi tersebut. Maka tidak satu pun ekspansi
$x$ hanya menggunakan digit $0$ dan $2$ pada $k$ digit pertamanya, bertentangan
dengan $x\in C_k$. Jadi, salah satu ekspansi $x$ hanya menggunakan digit $0$
dan $2$, sehingga $x\in C$. Dengan demikian,
\begin{equation*}
C = \bigcap_{k=1}^\infty C_k .
\end{equation*}
Lihat \figureref{fig:cantor}.

Hal-hal berikut diserahkan sebagai latihan untuk dibuktikan:
\begin{enumerate}[(i)]
\item Setiap $C_k$ merupakan gabungan berhingga dari interval-interval tertutup.
Himpunan ini diperoleh dari $C_{k-1}$ dengan membuang
\myquote{sepertiga tengah terbuka} dari setiap interval tertutup.
\item Setiap $C_k$ tertutup, sehingga $C$ tertutup.
\item 
$m^*(C_k) = {\left(\frac{2}{3}\right)}^k$.
\item Oleh karena itu,
$m^*(C) = 0$.
\item Terdapat korespondensi satu-satu antara $C$ dan $[0,1]$; dengan kata lain,
$C$ tak terhitung.
\end{enumerate}
\begin{myfigureht}
\myincludepdft{cantorfig}{%
Diagram beberapa himpunan pada garis bilangan yang disusun dalam beberapa baris.
Baris pertama adalah C subskrip 0, berupa satu ruas garis utuh. C subskrip 1
pada baris kedua diperoleh dari ruas pada baris pertama dengan menghapus
sepertiga tengah terbukanya, sehingga tersisa dua ruas garis utuh. Baris
berikutnya adalah C subskrip 2; sepertiga tengah terbuka dari kedua ruas kembali
dihapus sehingga tersisa empat ruas.
Pola ini berlanjut dan digambarkan hingga C subskrip 4, yang terdiri atas
16 ruas garis utuh.}
\caption{Konstruksi himpunan Cantor.\label{fig:cantor}}
\end{myfigureht}
\end{example}


\subsection{Citra himpunan nol oleh fungsi diferensiabel}

Sebelum meninjau citra himpunan berukuran nol, mari kita lihat bagaimana fungsi
diferensiabel secara kontinu memetakan sebuah bola.

\begin{lemma} \label{lemma:ballmapder}
Misalkan $U\subset\R^n$ terbuka, $B\subset U$ merupakan bola terbuka atau bola
tertutup berjari-jari tidak lebih dari $r$, dan
$f\colon U\to\R^n$ diferensiabel secara kontinu. Misalkan pula $M>0$ dan
$\bnorm{f'(x)}\leq M$ untuk setiap $x\in B$. Maka $f(B)\subset B'$, dengan
$B'$ suatu bola yang terbuka jika $B$ terbuka dan tertutup jika $B$ tertutup,
serta berjari-jari tidak lebih dari $Mr$.
\end{lemma}

\begin{proof}
Misalkan mula-mula $B=B(y,\rho)$ terbuka, dengan $0<\rho\leq r$. Karena bola
$B$ konveks, \propref{mv:prop:convexlip} menyatakan bahwa
$\bnorm{f(u)-f(v)}\leq M\snorm{u-v}$ untuk setiap $u,v\in B$. Jadi, untuk
$x\in B$, $\snorm{x-y}<\rho$ dan $\bnorm{f(x)-f(y)}<M\rho$. Dengan kata lain,
$f(B)\subset B\bigl(f(y),M\rho\bigr)$, dan bola terakhir berjari-jari tidak
lebih dari $Mr$.

Jika $B$ tertutup, tulis $B=\overline{B(y,\rho)}$ dengan $0<\rho\leq r$.
Karena $f$ kontinu,
$f(B)=f\bigl(\overline{B(y,\rho)}\bigr)\subset
\overline{f\bigl(B(y,\rho)\bigr)}\subset
\overline{B\bigl(f(y),M\rho\bigr)}$, sebab
$f(\widebar{A})\subset\overline{f(A)}$ untuk setiap $A\subset U$ yang memenuhi
$\widebar{A}\subset U$.
% dan untuk setiap f kontinu (limit barisan di A dipetakan ke limit barisan di f(A)).
\end{proof}

Citra himpunan berukuran nol oleh pemetaan kontinu belum tentu berukuran nol,
meskipun menunjukkan hal ini memerlukan usaha (lihat latihan). Namun, jika
pemetaan tersebut diferensiabel secara kontinu, pemetaan itu tidak dapat
\myquote{meregangkan} himpunan secara berlebihan.

\begin{prop} \label{prop:imagenull}
Misalkan $U\subset\R^n$ terbuka dan $f\colon U\to\R^n$ diferensiabel secara
kontinu. Jika $E\subset U$ berukuran nol, maka $f(E)$ berukuran nol.
\end{prop}

\begin{proof}
Kita membuktikan proposisi untuk $E$ kompak dan menyerahkan kasus umum sebagai
latihan. Misalkan $E$ kompak dan berukuran nol. Jika $E=\varnothing$, hasilnya
langsung benar; jadi, andaikan $E\neq\varnothing$.

Pertama, kita akan mengganti $U$ dengan himpunan terbuka yang lebih kecil agar
$\bnorm{f'(x)}$ terbatas. Untuk setiap $x\in E$, pilih bola terbuka $B(x,r_x)$
sedemikian sehingga bola tertutup $C(x,r_x)\subset U$. Berdasarkan kekompakan,
terdapat berhingga banyak titik $x_1,x_2,\ldots,x_q$ sedemikian sehingga
bola-bola $B(x_j,r_{x_j})$ menutupi $E$. Definisikan
\begin{equation*}
U' \coloneqq \bigcup_{j=1}^q B(x_j,r_{x_j}), \qquad
K \coloneqq \bigcup_{j=1}^q C(x_j,r_{x_j}).
\end{equation*}
Kita mempunyai $E\subset U'\subset K\subset U$. Sebagai gabungan berhingga
himpunan-himpunan kompak, $K$ kompak. Fungsi yang memetakan $x$ ke
$\bnorm{f'(x)}$ kontinu. Oleh karena itu, terdapat $M>0$ sedemikian sehingga
$\bnorm{f'(x)}\leq M$ untuk setiap $x\in K$. Tanpa mengurangi keumuman, ganti
$U$ dengan $U'$ dan batasi $f$ pada $U'$. Mulai sekarang, andaikan
$\bnorm{f'(x)}\leq M$ untuk setiap $x\in U$.

Himpunan $U$ yang baru merupakan gabungan berhingga dari bola-bola terbatas,
sehingga $U\neq\R^n$. Untuk setiap $x\in E$, ambil jari-jari maksimum
$\delta_x>0$ sedemikian sehingga
$B(x,\delta_x)\subset U$, lalu tetapkan
$\delta\coloneqq\inf_{x\in E}\delta_x$. Kita tunjukkan bahwa $\delta>0$.
Ambil barisan $\{x_j\}_{j=1}^\infty$ di $E$ sedemikian sehingga
$\delta_{x_j}\to\delta$. Berdasarkan kekompakan $E$, barisan ini memiliki
subbarisan yang konvergen ke suatu $y\in E$; namai kembali subbarisan tersebut
sebagai $\{x_j\}$. Untuk $j$ cukup besar,
$\snorm{x_j-y}<\frac{\delta_y}{2}$, sehingga ketaksamaan segitiga memberikan
$\delta_{x_j}>\frac{\delta_y}{2}$. Jadi,
$\delta\geq\frac{\delta_y}{2}>0$.

Tetapkan sembarang $\eta>0$. Terdapat bola-bola terbuka
$B_1,B_2,\ldots,B_k$, masing-masing berjari-jari
$r_j<\nicefrac{\delta}{2}$, $j=1,2,\ldots,k$, sedemikian sehingga
\begin{equation*}
E \subset B_1 \cup B_2 \cup \cdots \cup B_k
\qquad \text{dan} \qquad
\sum_{j=1}^k r_j^n < \frac{\eta}{M^n}.
\end{equation*}
Kita boleh membuang setiap bola yang tidak memotong $E$. Untuk setiap bola
$B_j$ yang tersisa, pilih $z_j\in B_j\cap E$. Jika $w\in B_j$, maka
$\snorm{w-z_j}<2r_j<\delta\leq\delta_{z_j}$, sehingga $B_j\subset U$.
Berdasarkan \lemmaref{lemma:ballmapder}, terdapat bola-bola terbuka
$B_1',B_2',\ldots,B_k'$, dengan masing-masing $B_j'$ berjari-jari $Mr_j$ dan
$f(B_j)\subset B_j'$ untuk setiap $j$. Maka,
\begin{equation*}
\begin{aligned}
f(E)
&
\subset f(B_1) \cup f(B_2) \cup \cdots \cup f(B_k)
\\
&
\subset B_1' \cup B_2' \cup \cdots \cup B_k'
\end{aligned}
\qquad \text{dan} \qquad
\sum_{j=1}^k {(Mr_j)}^n
 < \eta. \qedhere
\end{equation*}
\end{proof}

\subsection{Latihan}

\begin{exercise}
Lengkapi bukti \propref{mv:prop:compactnull}: Tunjukkan bahwa bola dapat
digunakan sebagai pengganti persegi panjang.
\end{exercise}

\begin{exercise} \label{exercise:outermeasuremono}
Jika $A\subset B$, maka $m^*(A)\leq m^*(B)$.
\end{exercise}

\begin{exercise}
Misalkan $X\subset\R^n$ merupakan himpunan sedemikian sehingga, untuk setiap
$\epsilon>0$, terdapat himpunan $Y$ yang memenuhi $X\subset Y$ dan
$m^*(Y)\leq\epsilon$. Buktikan bahwa $X$ berukuran nol.
\end{exercise}

\begin{exercise} \label{exercise:rectanglemeasure}
Tunjukkan bahwa jika $R\subset\R^n$ merupakan persegi panjang tertutup, maka
$m^*(R)=V(R)$.
\end{exercise}

\begin{exercise}
Penutupan himpunan berukuran nol dapat sangat besar. Temukan contoh himpunan
$S\subset\R^n$ yang berukuran nol, tetapi penutupannya memenuhi
$\widebar{S}=\R^n$.
\end{exercise}

\begin{exercise}
\pagebreak[2]
Buktikan kasus umum \propref{prop:imagenull} tanpa menggunakan kekompakan:
\begin{enumerate}[a)]
\item
Tiru bukti tersebut untuk menunjukkan bahwa proposisi berlaku jika $E$
\emph{\myindex{relatif kompak}}; himpunan $E\subset U$ disebut relatif kompak
jika penutupan $E$ dalam topologi subruang pada $U$ bersifat kompak. Dengan kata
lain, terdapat himpunan kompak $K$ sedemikian sehingga $E\subset K\subset U$.\\
Petunjuk: Batas atas bagi norma turunan masih berlaku, tetapi pada bagian kedua
bukti Anda perlu menggunakan keluarga bola yang terhitung. Berhati-hatilah
karena penutupan $E$ tidak lagi harus berukuran nol.
\item
Sekarang, buktikan proposisi itu untuk setiap himpunan nol $E$.\\
Petunjuk: Pertama, tunjukkan bahwa
$\{x\in U : \snorm{x-y}\geq
\nicefrac{1}{m} \text{ untuk setiap } y\notin U
\text{ dan } \snorm{x}\leq m\}$ kompak untuk setiap $m>0$.
\end{enumerate}
\end{exercise}

\begin{exercise}
Misalkan $U\subset\R^n$ terbuka dan $f\colon U\to\R$ diferensiabel secara
kontinu. Misalkan
$G\coloneqq\bigl\{(x,y)\in U\times\R:y=f(x)\bigr\}$ merupakan grafik $f$.
Tunjukkan bahwa $G$ berukuran nol di $\R^{n+1}$.
\end{exercise}

\begin{exercise}
Diberikan persegi panjang tertutup $R\subset\R^n$, tunjukkan bahwa untuk
setiap $\epsilon>0$, terdapat bilangan $s>0$ dan berhingga banyak kubus terbuka
$C_1,C_2,\ldots,C_k$ dengan panjang sisi $s$ sedemikian sehingga
$R\subset C_1\cup C_2\cup\cdots\cup C_k$ dan
\begin{equation*}
\sum_{j=1}^k V(C_j) \leq V(R) + \epsilon .
\end{equation*}
\end{exercise}

\begin{exercise}
Tunjukkan bahwa terdapat bilangan $k=k(n,r,\delta)$ yang hanya bergantung pada
$n$, $r$, dan $\delta$, dengan sifat berikut: Diberikan
$B(x,r)\subset\R^n$ dan $\delta>0$, terdapat $k$ bola terbuka
$B_1,B_2,\ldots,B_k$, masing-masing berjari-jari tidak lebih dari $\delta$,
sedemikian sehingga $B(x,r)\subset B_1\cup B_2\cup\cdots\cup B_k$.
Perhatikan bahwa $k$ dapat dipilih agar hanya bergantung pada $n$ dan rasio
$\nicefrac{\delta}{r}$.
\end{exercise}

\begin{exercise}[Menantang]
Buktikan pernyataan-pernyataan dalam \exampleref{example:cantor}, yaitu:
\begin{enumerate}[a)]
\item
Setiap $C_k$ merupakan gabungan berhingga dari interval-interval tertutup,
sehingga $C$ tertutup.
\item
$m^*(C_k) = {\left(\frac{2}{3}\right)}^k$.
\item
$m^*(C) = 0$.
\item
Terdapat korespondensi satu-satu antara $C$ dan $[0,1]$.
\end{enumerate}
\end{exercise}

\begin{exercise}
\pagebreak[3]
Buktikan bahwa himpunan Cantor dalam \exampleref{example:cantor} tidak memuat
interval yang panjangnya positif. Artinya, jika $a<b$, terdapat titik
$x\notin C$ sedemikian sehingga $a<x<b$.
\\
Perhatikan suatu akibat dari pernyataan ini. Setiap himpunan terbuka di $\R$
merupakan gabungan terhitung dari interval-interval yang saling lepas, tetapi
himpunan tertutup (meskipun merupakan komplemen himpunan terbuka) tidak harus
merupakan gabungan terhitung dari interval-interval ataupun titik-titik.
\end{exercise}

\begin{samepage}
\begin{exercise}[Menantang]
Mari kita bangun \emph{\myindex{fungsi Cantor}}, yang juga disebut
\emph{\myindex{tangga Iblis}}. Misalkan $C$ adalah himpunan Cantor dan $C_k$
seperti dalam \exampleref{example:cantor}. Tuliskan $x\in[0,1]$ dalam
representasi terner
$x=\sum_{n=1}^{\infty}d_n3^{-n}$. Jika $d_n\neq1$ untuk setiap $n$, tetapkan
$c_n\coloneqq\frac{d_n}{2}$ untuk setiap $n$. Jika tidak, misalkan $k$ adalah
bilangan bulat terkecil yang memenuhi $d_k=1$. Tetapkan
$c_n\coloneqq\frac{d_n}{2}$ jika $n<k$, $c_k\coloneqq1$, dan
$c_n\coloneqq0$ jika $n>k$. Definisikan
\begin{equation*}
\varphi(x) \coloneqq \sum_{n=1}^\infty c_n \, 2^{-n} .
\end{equation*}
\begin{enumerate}[a)]
\item
Buktikan bahwa $\varphi$ kontinu dan monoton naik (lihat
\figureref{fig:cantorfunction}).
\item
Buktikan bahwa untuk $x\notin C$, fungsi $\varphi$ diferensiabel di $x$ dan
$\varphi'(x)=0$. (Perhatikan bahwa $\varphi'$ ada dan bernilai nol kecuali pada
suatu himpunan berukuran nol, tetapi fungsi tersebut tetap dapat naik dari $0$
ke $1$.)
\item
Definisikan $\psi\colon[0,1]\to[0,2]$ dengan
$\psi(x)\coloneqq\varphi(x)+x$. Tunjukkan bahwa $\psi$ kontinu, naik tegas, dan
bijektif.
\item
Buktikan bahwa meskipun $m^*(C)=0$, berlaku
$m^*\bigl(\psi(C)\bigr)\neq0$. Dengan kata lain, fungsi kontinu tidak harus
memetakan himpunan berukuran nol ke himpunan berukuran nol. Petunjuk:
$m^*\bigl(\psi([0,1]\setminus C)\bigr)=1$, sedangkan
$m^*\bigl([0,2]\bigr)=2$.
\end{enumerate}
\end{exercise}
\end{samepage}

\begin{myfigureht}
\myincludegraphics{cantorfunction}{%
Grafik suatu fungsi pada interval 0 hingga 1, yang ditampilkan pada sumbu
horizontal. Fungsi ini monoton, berawal dari 0 di sisi kiri, dan berakhir pada
1 di sisi kanan. Grafik memiliki banyak ruas mendatar yang bersesuaian dengan
interval-interval sepertiga tengah yang dibuang selama konstruksi himpunan
Cantor. Sebagai contoh, fungsi bernilai satu per dua dari satu per tiga hingga
dua per tiga. Fungsi bernilai satu per empat dari satu per sembilan hingga dua
per sembilan, dan bernilai tiga per empat dari tujuh per sembilan hingga delapan
per sembilan. Pola yang sama berlaku untuk semua interval sepertiga tengah
lainnya.}
\caption{Fungsi Cantor atau tangga Iblis (fungsi $\varphi$ dari latihan).
\label{fig:cantorfunction}}
\end{myfigureht}

\begin{exercise} \label{exercise:allowfiniteseqsinoutermeasure}
Buktikan bahwa ukuran luar yang sama diperoleh jika definisinya mengizinkan
barisan hingga maupun tak berhingga. Definisikan
$\mu^*(S)\coloneqq\inf\,\sum_{j\in I}V(R_j)$, dengan infimum diambil atas semua
himpunan indeks terhitung $I$ (hingga atau tak berhingga) dan semua keluarga
persegi panjang terbuka $\{R_j\}_{j\in I}$ yang memenuhi
$S\subset\bigcup_{j\in I}R_j$. Buktikan bahwa untuk setiap $S\subset\R^n$,
$\mu^*(S)=m^*(S)$.
\end{exercise}

\begin{exercise}
Buktikan bahwa untuk sebarang dua subhimpunan $A,B\subset\R^n$, berlaku
$m^*(A\cup B)\leq m^*(A)+m^*(B)$.
\end{exercise}

\begin{exercise}
Misalkan $A,B\subset\R^n$ dan $m^*(B)=0$. Buktikan bahwa
$m^*(A\cup B)=m^*(A)$.
\end{exercise}

\begin{exercise}[Menantang] \label{exercise:outermeasureofsumofrectangles}
Misalkan $R_1,R_2,\ldots,R_n$ merupakan persegi panjang terbuka yang saling
lepas. Buktikan bahwa
$m^*(R_1\cup R_2\cup\cdots\cup R_n)
=m^*(R_1)+m^*(R_2)+\cdots+m^*(R_n)$.
Petunjuk: Beberapa latihan di atas mungkin sangat berguna.
\end{exercise}


%%%%%%%%%%%%%%%%%%%%%%%%%%%%%%%%%%%%%%%%%%%%%%%%%%%%%%%%%%%%%%%%%%%%%%%%%%%%%%

\sectionnewpage
\section{Himpunan fungsi yang terintegralkan secara Riemann}
\label{sec:riemannlebesgue}

%mbxINTROSUBSECTION

\sectionnotes{1 kuliah}

\subsection{Osilasi dan kekontinuan}

Misalkan $D \subset \R^n$ dan $f \colon D \to \R$.
Alih-alih sekadar mengatakan bahwa $f$ kontinu atau tidak kontinu di
suatu titik $x \in D$,
kita hendak mengukur seberapa diskontinu $f$ di $x$.
Untuk setiap $\delta > 0$, definisikan \emph{osilasi}
$f$ pada bola berjari-jari $\delta$ dalam topologi subruang,
$B_D(x,\delta) = B_{\R^n}(x,\delta) \cap D$, sebagai
\glsadd{not:oscillation}
\begin{equation*}
o(f,x,\delta) \coloneqq
{\sup_{y \in B_D(x,\delta)} f(y)}
-
{\inf_{y \in B_D(x,\delta)} f(y)}
= 
\sup_{y_1,y_2 \in B_D(x,\delta)} \bigl(f(y_1)-f(y_2)\bigr) .
\end{equation*}
Artinya, $o(f,x,\delta)$ adalah panjang interval terkecil yang memuat citra
$f\bigl(B_D(x,\delta)\bigr)$; panjang ini mungkin tak hingga.
Jika supremum dan infimum dipahami sebagai bilangan real diperluas,
definisi ini juga bermakna untuk fungsi tak terbatas, yang osilasinya dapat
bernilai $\infty$, meskipun kita terutama akan membahas fungsi terbatas.
Jelas bahwa $o(f,x,\delta) \geq 0$ dan
$o(f,x,\delta) \leq o(f,x,\delta')$ apabila $\delta < \delta'$.
Oleh karena itu, limit saat $\delta \to 0$ dari kanan ada sebagai bilangan
real diperluas, dan kita mendefinisikan \emph{\myindex{osilasi}} $f$
di $x$ sebagai
\begin{equation*}
o(f,x) \coloneqq
\lim_{\delta \to 0^+}
o(f,x,\delta) =
\inf_{\delta > 0}
o(f,x,\delta) .
\end{equation*}

Kita akan membuktikan bahwa suatu fungsi kontinu di $x$ jika dan hanya jika
$o(f,x) = 0$.
Sebagai contoh, misalkan $f \colon \R \to \R$ adalah fungsi Dirichlet, yakni
$f(x)=1$ jika $x \in \Q$ dan $f(x)=0$ selainnya. Maka $o(f,x) = 1$ untuk setiap
$x$, sebab setiap interval terbuka memuat bilangan rasional maupun irasional.
Oleh sebab itu, $f$ tidak kontinu di titik mana pun.
Sebagai contoh lain, yang barangkali menjadi asal istilah ini,
misalkan $g \colon \R \to \R$ diberikan oleh
$g(x) = \sin(\nicefrac{1}{x})$ untuk $x\neq 0$ dan $g(0)=0$;
lihat \figureref{fig:sin1xosc}.
Pada titik diskontinuitas $x=0$, kita memperoleh $o(g,0) = 2$, sebab
dalam setiap lingkungan $0$, fungsi tersebut mengambil baik nilai $1$ maupun $-1$.
Untuk setiap $x \neq 0$, fungsi itu kontinu sehingga, seperti yang akan kita lihat,
$o(g,x) = 0$.

\begin{myfigureht}
% juga digunakan dalam 3.1
\myincludegraphics{sin1xfig}{%
Grafik sebuah fungsi pada interval di sekitar titik asal.
Fungsi tersebut berosilasi
di antara minus 1 dan 1
makin cepat ketika x mendekati 0, sedemikian cepat sehingga di dekat titik asal
grafik tersebut tampak hitam pekat.}
\caption{Grafik $\sin(\nicefrac{1}{x})$.\label{fig:sin1xosc}}
\end{myfigureht}

\begin{prop}
\pagebreak[2]
Suatu fungsi
$f \colon D \to \R$ kontinu di $x \in D$ jika dan hanya jika $o(f,x) = 0$.
\end{prop}

\begin{proof}
Pertama-tama, andaikan $f$ kontinu di $x \in D$. Untuk sebarang $\epsilon > 0$,
terdapat $\delta > 0$ sedemikian sehingga untuk $y \in B_D(x,\delta)$,
berlaku $\babs{f(x)-f(y)} < \epsilon$. Oleh karena itu, jika $y_1,y_2 \in
B_D(x,\delta)$, maka
\begin{equation*}
f(y_1)-f(y_2) =
\bigl(f(y_1)-f(x)\bigr)-\bigl(f(y_2)-f(x)\bigr) < \epsilon + \epsilon = 2 \epsilon .
\end{equation*}
Dengan mengambil supremum atas $y_1$ dan $y_2$, diperoleh
\begin{equation*}
o(f,x,\delta) = 
\sup_{y_1,y_2 \in B_D(x,\delta)} \bigl(f(y_1)-f(y_2)\bigr)
\leq
2 \epsilon .
\end{equation*}
Karena $o(f,x) \leq o(f,x,\delta) \leq 2\epsilon$ dan $\epsilon > 0$ dapat
dipilih sebarang, maka $o(f,x) = 0$.

Sebaliknya, andaikan $o(f,x) = 0$. Untuk sebarang $\epsilon > 0$,
pilih $\delta > 0$ sedemikian sehingga $o(f,x,\delta) < \epsilon$. Jika
$y \in B_D(x,\delta)$, maka
\begin{equation*}
\babs{f(x)-f(y)}
\leq
\sup_{y_1,y_2 \in B_D(x,\delta)} \bigl(f(y_1)-f(y_2)\bigr)
=
o(f,x,\delta) < \epsilon. \qedhere
\end{equation*}
\end{proof}

\begin{prop} \label{prop:seclosed}
Misalkan $D \subset \R^n$ tertutup,
$f \colon D \to \R$, dan $\epsilon > 0$.
Himpunan $\bigl\{ x \in D : o(f,x) \geq \epsilon \bigr\}$ tertutup di $\R^n$.
\end{prop}

\begin{proof}
Secara ekuivalen, kita hendak menunjukkan bahwa
$G \coloneqq \bigl\{ x \in D : o(f,x) < \epsilon \bigr\}$ terbuka dalam topologi subruang pada $D$.
Ambil $x \in G$.
Karena $\inf_{\delta > 0} o(f,x,\delta) < \epsilon$, pilih $\delta > 0$
sedemikian sehingga
\begin{equation*}
o(f,x,\delta) < \epsilon .
\end{equation*}
Ambil sebarang $\xi \in B_D(x,\nicefrac{\delta}{2})$. Perhatikan bahwa
$B_D(\xi,\nicefrac{\delta}{2}) \subset B_D(x,\delta)$. Oleh karena itu,
\begin{equation*}
o(f,\xi,\nicefrac{\delta}{2}) =
\sup_{y_1,y_2 \in B_D(\xi,\nicefrac{\delta}{2})} \bigl(f(y_1)-f(y_2)\bigr) 
\leq
\sup_{y_1,y_2 \in B_D(x,\delta)} \bigl(f(y_1)-f(y_2)\bigr) = o(f,x,\delta) <
\epsilon .
\end{equation*}
Jadi, $o(f,\xi) < \epsilon$ pula. Karena hal ini berlaku untuk setiap $\xi \in
B_D(x,\nicefrac{\delta}{2})$, himpunan $G$ terbuka dalam topologi subruang
pada $D$, sehingga $D \setminus G$ tertutup dalam $D$.
Karena $D$ tertutup di $\R^n$, himpunan $D \setminus G$ juga tertutup di
$\R^n$, seperti yang diklaim.
\end{proof}


\subsection{Himpunan fungsi yang terintegralkan secara Riemann}

Kita telah melihat bahwa fungsi kontinu terintegralkan secara Riemann, tetapi
beberapa jenis diskontinuitas tetap diperbolehkan.
Ternyata, suatu fungsi terintegralkan secara Riemann jika dan hanya jika
titik-titik diskontinuitasnya membentuk himpunan berukuran nol.

\begin{thm}[Riemann--Lebesgue atau Lebesgue--Vitali\footnote{%
\href{https://en.wikipedia.org/wiki/Giuseppe_Vitali}{Giuseppe Vitali}
(1875--1932) adalah matematikawan Italia. Perhatikan pula bahwa
nama Riemann--Lebesgue sering
merujuk pada hasil seperti
\volIref{\exerciseref*{vI-exercise:rieleblem} dari jilid I}{\exerciseref{exercise:rieleblem}}.}]%
\index{teorema Riemann--Lebesgue}%
\index{teorema Lebesgue--Vitali}%
Misalkan $R \subset \R^n$ persegi panjang tertutup dan $f \colon R \to \R$
fungsi terbatas. Maka $f$ terintegralkan secara Riemann jika dan hanya jika
himpunan titik diskontinuitas $f$ berukuran nol.
\end{thm}

\begin{proof}
Misalkan $S \subset R$ adalah himpunan titik diskontinuitas $f$, yakni
$S = \bigl\{ x \in R : o(f,x) > 0 \bigr\}$.
Andaikan $S$ berukuran nol: $m^*(S) = 0$.
Untuk membuktikan bahwa $f$ terintegralkan, kita mengurung himpunan bermasalah
itu di dalam suatu kumpulan subpersegi panjang partisi yang volume totalnya
kecil. Karena suatu partisi hanya memiliki berhingga banyak
subpersegi panjang, kita memerlukan kekompakan. Seandainya $S$ tertutup, maka
$S$ kompak dan dapat ditutupi oleh berhingga banyak persegi panjang kecil.
Sayangnya, $S$ sendiri pada umumnya tidak tertutup, tetapi himpunan berikut
tertutup. Untuk $\epsilon > 0$, definisikan
\begin{equation*}
S_\epsilon \coloneqq \bigl\{ x \in R : o(f,x) \geq \epsilon \bigr\} .
\end{equation*}
Berdasarkan \propref{prop:seclosed}, $S_\epsilon$ tertutup.
Karena $S_\epsilon\subset R$ dan $R$ terbatas, $S_\epsilon$ kompak. Selain itu,
$S_\epsilon \subset S$ dan $S$ berukuran nol,
sehingga $S_\epsilon$ berukuran nol.
Menurut \propref{mv:prop:compactnull}, terdapat berhingga banyak persegi panjang
terbuka $O_1,O_2,\ldots,O_k$ yang menutupi $S_\epsilon$ dan memenuhi
$\sum_{j=1}^k V(O_j) < \epsilon$.

Himpunan $T \coloneqq R \setminus ( O_1 \cup \cdots \cup O_k )$ tertutup dan
terbatas, sehingga kompak. Karena $o(f,x) < \epsilon$ untuk setiap $x \in T$,
untuk setiap $x \in T$ terdapat $\delta > 0$ sedemikian sehingga
$o(f,x,\delta) < \epsilon$. Oleh sebab itu, terdapat persegi panjang tertutup
$T_x \subset B(x,\delta)$ yang cukup kecil, dengan $x$ berada di bagian dalam
$T_x$, sedemikian sehingga
\begin{equation*}
\sup_{y\in T_x \cap R} f(y) - \inf_{y\in T_x \cap R} f(y) < \epsilon.
\end{equation*}
Bagian dalam semua persegi panjang $T_x$ menutupi $T$. Karena $T$ kompak,
terdapat berhingga banyak di antaranya, $T_1,T_2,\ldots,T_m$, yang menutupi
$T$. Bentuk partisi $P$ dari titik-titik ujung interval koordinat
$T_1,T_2,\ldots,T_m$ dan $O_1,O_2,\ldots,O_k$, dengan mengabaikan titik-titik
yang berada di luar interval koordinat $R$. Subpersegi panjang
$R_1,R_2,\ldots,R_p$ dari $P$ memiliki sifat bahwa setiap $R_j$ termuat dalam
suatu $T_\ell$ atau dalam penutupan suatu $O_\ell$. Urutkan subpersegi panjang
tersebut sedemikian sehingga $R_1,R_2,\ldots,R_q$ merupakan subpersegi panjang
yang termuat dalam suatu $T_\ell$, sedangkan $R_{q+1},R_{q+2},\ldots,R_{p}$
merupakan sisanya.
Lihat \figureref{fig:nullsetintegrate}. Dengan demikian,
\begin{equation*}
\sum_{j=1}^q V(R_j) \leq V(R)
\qquad \text{dan} \qquad
\sum_{j=q+1}^p V(R_j)
\leq
\sum_{\ell=1}^k V(O_\ell)
< \epsilon .
\end{equation*}
Ketaksamaan kedua berlaku karena, untuk setiap $j>q$, kita dapat memilih satu
$\ell(j)$ sedemikian sehingga $R_j \subset \widebar{O}_{\ell(j)}$. Untuk setiap
$\ell$ tetap, subpersegi panjang yang dipasangkan dengan $\ell$ memiliki bagian
dalam yang saling lepas dan semuanya termuat dalam $\widebar{O}_\ell$, sehingga
jumlah volumenya paling besar $V(O_\ell)$. Dengan menjumlahkan atas semua
$\ell$, diperoleh ketaksamaan tersebut.
Misalkan $m_j$ dan $M_j$ masing-masing adalah infimum dan supremum $f$ pada
$R_j$, seperti biasa. Jika $R_j \subset T_\ell$ untuk suatu $\ell$, maka
$M_j-m_j < \epsilon$.
Pilih $B > 0$ sedemikian sehingga
$\babs{f(x)} \leq B$ untuk setiap $x \in R$. Maka $M_j-m_j \leq 2B$ pada semua
subpersegi panjang, sehingga
\begin{equation*}
\begin{split}
U(P,f)-L(P,f)
& =
\sum_{j=1}^p (M_j-m_j) V(R_j)
\\
& =
\left(
\sum_{j=1}^q (M_j-m_j) V(R_j)
\right)
+
\left(
\sum_{j=q+1}^p (M_j-m_j) V(R_j)
\right)
\\
& \leq
\left(
\sum_{j=1}^q \epsilon\, V(R_j)
\right)
+
\left(
\sum_{j=q+1}^p 2 B\, V(R_j)
\right)
\\
& <
\epsilon\, V(R)
+
2B \epsilon = \epsilon \bigl(V(R)+2B\bigr) .
\end{split}
\avoidbreak
\end{equation*}
Konstruksi di atas berlaku untuk setiap $\epsilon>0$, sedangkan
$V(R)+2B>0$ tetap. Jadi, untuk sebarang $\eta>0$, pilih sejak awal
$\epsilon<\eta/\bigl(V(R)+2B\bigr)$; ruas kanan menjadi lebih kecil daripada
$\eta$. Dengan demikian, $f$ terintegralkan.

\begin{myfigureht}
\myincludepdft{nullsetintegrate}{%
Diagram sebuah persegi panjang yang dibagi oleh banyak garis vertikal dan
horizontal bertitik-titik. Di tengah diagram terdapat garis tebal gelap yang
bergerigi. Garis ini ditutupi oleh persegi panjang-persegi panjang berarsir,
yang sisi-sisinya berimpit dengan beberapa garis bertitik-titik.}
\caption{Persegi panjang $R$, dengan $S_\epsilon$ ditandai oleh garis hitam
tebal dan setiap $O_\ell$ oleh persegi panjang berarsir. Partisi ditunjukkan
oleh garis-garis bertitik-titik. Perhatikan bagaimana garis-garis partisi
membagi bagian dari setiap $O_\ell$ yang berada di dalam $R$ menjadi
subpersegi panjang $R_j$.\label{fig:nullsetintegrate}}
\end{myfigureht}

\medskip

\pagebreak[2]
Untuk arah sebaliknya, andaikan $f$ terintegralkan secara Riemann
pada $R$.
Misalkan $S$ kembali menyatakan himpunan titik diskontinuitas $f$. Tinjau
barisan himpunan
\begin{equation*}
S_{1/k} = \bigl\{ x \in R : o(f,x) \geq \nicefrac{1}{k} \bigr\}.
\end{equation*}
Tetapkan $k \in \N$.
Untuk sebarang $\epsilon > 0$, pilih partisi $P$ dengan subpersegi panjang
$R_1,R_2,\ldots,R_p$ sedemikian sehingga
\begin{equation*}
U(P,f)-L(P,f) =
\sum_{j=1}^p (M_j-m_j) V(R_j)
< \epsilon .
\end{equation*}
Urutkan $R_1,R_2,\ldots,R_p$ sedemikian sehingga bagian dalam masing-masing
$R_1,R_2,\ldots,R_{q}$ beririsan dengan $S_{1/k}$, sedangkan bagian dalam
masing-masing $R_{q+1},R_{q+2},\ldots,R_p$ tidak beririsan dengan $S_{1/k}$.
Tuliskan $R_j^\circ$ untuk bagian dalam $R_j$.
Andaikan $j \leq q$ dan ambil
$x \in R_j^\circ \cap S_{1/k}$.
Pilih $\delta > 0$ yang cukup kecil sehingga $B(x,\delta) \subset R_j$.
Karena $x \in S_{1/k}$, diperoleh
$o(f,x,\delta) \geq o(f,x) \geq \nicefrac{1}{k}$.
Bersama dengan $B(x,\delta) \subset R_j$, hal ini mengakibatkan
$M_j-m_j \geq \nicefrac{1}{k}$.
Maka
\begin{equation*}
\epsilon >
\sum_{j=1}^p (M_j-m_j) V(R_j)
\geq
\sum_{j=1}^q (M_j-m_j) V(R_j)
\geq
\frac{1}{k}
\sum_{j=1}^q V(R_j) .
\end{equation*}
Dengan kata lain,
$\sum_{j=1}^q V(R_j) < k \epsilon$.
Misalkan $G$ adalah gabungan batas semua subpersegi panjang dalam $P$.
Himpunan $G$ berukuran nol karena merupakan gabungan berhingga himpunan sejenis
dengan yang dibahas dalam \exampleref{mv:example:planenull}.
Dengan demikian,
\begin{equation*}
S_{1/k} \subset R_1^\circ \cup R_2^\circ \cup \cdots \cup R_q^\circ \cup G .
\end{equation*}
Untuk menunjukkan secara eksplisit bahwa $S_{1/k}$ berukuran nol, ambil
sebarang $\eta>0$. Dalam argumen di atas, pilih
$\epsilon<\eta/(2k)$, sehingga jumlah volume $R_1^\circ,\ldots,R_q^\circ$
kurang daripada $\eta/2$. Tutupi pula $G$ dengan persegi panjang terbuka yang
volume totalnya kurang daripada $\eta/2$. Jadi, $S_{1/k}$ dapat ditutupi dengan
volume total kurang daripada $\eta$, sehingga $S_{1/k}$ berukuran nol. Karena
\begin{equation*}
S = \bigcup_{k=1}^\infty S_{1/k}
\end{equation*}
dan gabungan terhitung himpunan-himpunan berukuran nol juga berukuran nol,
maka $S$ berukuran nol.
\end{proof}

\begin{cor} \label{cor:closednessofriemannintegrable}
\pagebreak[2]
Misalkan $R \subset \R^n$ merupakan persegi panjang tertutup.  Misalkan
$\sR(R)$ adalah himpunan fungsi yang terintegralkan secara Riemann pada $R$.
Maka
\begin{enumerate}[(i)]
\item
$\sR(R)$ adalah aljabar atas $\R$:
Jika $f,g \in \sR(R)$ dan $a \in \R$, maka $af \in \sR(R)$,
$f+g \in \sR(R)$, dan $fg \in \sR(R)$.
\item
Jika $f,g \in \sR(R)$ dan
\begin{equation*}
\varphi(x) \coloneqq \max \bigl\{ f(x) , g(x) \bigr\} ,
\qquad
\psi(x) \coloneqq \min \bigl\{ f(x) , g(x) \bigr\} ,
\end{equation*}
maka $\varphi, \psi \in \sR(R)$.
\item
Jika $f \in \sR(R)$, maka $\sabs{f} \in \sR(R)$, dengan
$\sabs{f}(x) \coloneqq \babs{f(x)}$.
\item
Jika $R' \subset \R^n$ merupakan persegi panjang tertutup lain,
$U \subset \R^n$ dan $U' \subset \R^n$ merupakan himpunan terbuka sedemikian
sehingga $R \subset U$ dan $R' \subset U'$,
$g \colon U \to U'$ diferensiabel secara kontinu dan bijektif,
$g^{-1}$ diferensiabel secara kontinu,
$g(R) \subset R'$,
dan $f \in \sR(R')$,
maka komposisi $f \circ g$ terintegralkan secara Riemann pada $R$.
\end{enumerate}
\end{cor}

Pembuktiannya termuat dalam latihan-latihan.

\subsection{Latihan}

\begin{exercise}
Misalkan $f \colon (a,b) \times (c,d) \to \R$ merupakan fungsi kontinu yang
terbatas. Tunjukkan bahwa integral $f$ pada $R = [a,b] \times [c,d]$
terdefinisi dengan baik dan nilainya ditentukan secara tunggal. Dengan kata
lain, perluas $f$ ke $R$ dengan menetapkan nilai-nilai sembarang tetapi terbatas
pada batas $R$, lalu hitung integralnya dan tunjukkan bahwa hasilnya tidak
bergantung pada nilai-nilai $f$ pada batas.
\end{exercise}

\begin{exercise}
Misalkan $R \subset \R^n$ merupakan persegi panjang tertutup. Tunjukkan bahwa
$\sR(R)$, yakni himpunan fungsi yang terintegralkan secara Riemann, merupakan
aljabar atas $\R$. Artinya, tunjukkan bahwa jika $f,g \in \sR(R)$ dan
$a \in \R$, maka $af \in \sR(R)$, $f+g \in \sR(R)$, dan $fg \in \sR(R)$.
\end{exercise}

\begin{exercise}
Misalkan $R \subset \R^n$ merupakan persegi panjang tertutup dan
$f \colon R \to \R$ merupakan fungsi terbatas yang bernilai nol di luar suatu
himpunan tertutup berukuran nol $E \subset R$. Tunjukkan bahwa $\int_R f$ ada
dan hitung nilainya.
\end{exercise}

\begin{exercise}
Misalkan $R \subset \R^n$ merupakan persegi panjang tertutup, sedangkan
$f \colon R \to \R$ dan $g \colon R \to \R$ merupakan dua fungsi yang
terintegralkan secara Riemann. Andaikan $f = g$ di luar suatu himpunan tertutup
berukuran nol $E \subset R$. Tunjukkan bahwa $\int_R f = \int_R g$.
\end{exercise}

\begin{samepage}
\begin{exercise}
Misalkan $R \subset \R^n$ merupakan persegi panjang tertutup dan
$f \colon R \to \R$ merupakan fungsi terbatas.
\begin{enumerate}[a)]
\item
Andaikan terdapat himpunan tertutup berukuran nol $E \subset R$ sedemikian
sehingga pembatasan $f|_{R\setminus E}$ kontinu. Maka $f \in \sR(R)$.
\item
Temukan contoh himpunan berukuran nol $E \subset R$ yang tidak tertutup
sedemikian sehingga pembatasan $f|_{R\setminus E}$ kontinu dan
$f \notin \sR(R)$.
\end{enumerate}
\end{exercise}
\end{samepage}

\begin{exercise}
Misalkan $R \subset \R^n$ merupakan persegi panjang tertutup dan
$f \colon R \to \R$ serta $g \colon R \to \R$ terintegralkan secara Riemann.
Tunjukkan bahwa
\begin{equation*}
\varphi(x) \coloneqq \max \bigl\{ f(x) , g(x) \bigr\}
\qquad
\text{dan}
\qquad
\psi(x) \coloneqq \min \bigl\{ f(x) , g(x) \bigr\}
\end{equation*}
terintegralkan secara Riemann.
\end{exercise}

\begin{exercise}
Misalkan $R \subset \R^n$ merupakan persegi panjang tertutup dan
$f \colon R \to \R$ terintegralkan secara Riemann. Tunjukkan bahwa $\sabs{f}$
terintegralkan secara Riemann. Petunjuk: Definisikan
$f_+(x) \coloneqq \max \bigl\{ f(x) , 0 \bigr\}$ dan
$f_-(x) \coloneqq \max \bigl\{ -f(x) , 0 \bigr\}$, lalu nyatakan $\sabs{f}$
menggunakan $f_+$ dan $f_-$.
\end{exercise}

\begin{exercise}
\leavevmode
\begin{enumerate}[a)]
\item
Misalkan $R \subset \R^n$ dan $R' \subset \R^n$ merupakan persegi panjang
tertutup, sedangkan $U \subset \R^n$ dan $U' \subset \R^n$ merupakan himpunan
terbuka sedemikian sehingga $R \subset U$ dan $R' \subset U'$.
Misalkan pula $g \colon U \to U'$ diferensiabel secara kontinu dan bijektif,
$g^{-1}$ diferensiabel secara kontinu, $g(R) \subset R'$, dan
$f \in \sR(R')$. Tunjukkan bahwa komposisi $f \circ g$ terintegralkan secara
Riemann pada $R$.
\item
Temukan contoh tandingan ketika $g$ tidak injektif.
Petunjuk: Cobalah $g(x,y) \coloneqq (x,0)$ dan
$R=R'=[0,1] \times [0,1]$.
\end{enumerate}
\end{exercise}

\begin{exercise}
Misalkan
$f \colon [0,1]^2 \to \R$ didefinisikan oleh
\begin{equation*}
f(x,y) \coloneqq
\begin{cases}
\frac{1}{kq} &
\begin{aligned}
&\text{jika } x,y \in \Q, \quad x = \frac{\ell}{k}, \quad y=\frac{p}{q}, \\
&\text{dengan kedua pecahan tersebut ditulis} \\
&\text{dalam bentuk paling sederhana,}
\end{aligned} \\
0            & \text{selain itu.}
\end{cases}
\end{equation*}
Tunjukkan bahwa $f \in \sR\bigl({[0,1]}^2\bigr)$.
\end{exercise}

\begin{exercise}
Hitung osilasi $o\bigl(f,(x,y)\bigr)$ untuk setiap $(x,y) \in \R^2$ bagi fungsi
\begin{equation*}
f(x,y) \coloneqq
\begin{cases}
\frac{xy}{x^2+y^2} & \text{jika } (x,y) \neq (0,0), \\
0                  & \text{jika } (x,y) = (0,0) .
\end{cases}
\end{equation*}
\end{exercise}

\begin{exercise}
Perhatikan fungsi popcorn $f \colon [0,1] \to \R$,
\begin{equation*}
f(x) \coloneqq
\begin{cases}
\frac{1}{q} & \text{jika } x \in \Q \text{ dan } x = \frac{p}{q}
 \text{ dalam bentuk pecahan paling sederhana,} \\
0           & \text{selain itu.}
\end{cases}
\end{equation*}
Hitung $o(f,x)$ untuk setiap $x \in [0,1]$.
\end{exercise}

\begin{exercise}
Misalkan $f \colon [a,b] \to \R$ dan $g \colon [c,d] \to \R$ terintegralkan
secara Riemann. Tunjukkan bahwa fungsi
$h \colon [a,b] \times [c,d] \to \R$ yang didefinisikan oleh
$h(x,y) \coloneqq f(x) g(y)$ terintegralkan secara Riemann dan memenuhi
\begin{equation*}
\int_{[a,b] \times [c,d]} h =
\left( \int_a^b f \right)
\left( \int_c^d g \right) .
\end{equation*}
\end{exercise}

\begin{exercise}
Misalkan $R \subset \R^n$ merupakan persegi panjang tertutup dan
$f \colon R \to \R$ merupakan fungsi yang terintegralkan secara Riemann dengan
$f(x) \geq 0$ untuk setiap $x \in R$. Tunjukkan bahwa jika $\int_R f = 0$,
maka terdapat himpunan berukuran nol $E \subset R$ sedemikian sehingga
$f(x) = 0$ untuk setiap $x \in R \setminus E$
(dikatakan bahwa \myquote{$f=0$ \myindex{hampir di mana-mana}}).
Catatan: Latihan ini secara khusus menyiratkan pernyataan yang cukup subtil
berikut: jika $V(R)>0$ dan $f(x) > 0$ untuk setiap $x \in R$, maka
$\int_R f > 0$.
\end{exercise}


%%%%%%%%%%%%%%%%%%%%%%%%%%%%%%%%%%%%%%%%%%%%%%%%%%%%%%%%%%%%%%%%%%%%%%%%%%%%%%

\sectionnewpage
\section{Himpunan terukur Jordan}
\label{sec:jordansets}

%mbxINTROSUBSECTION

\sectionnotes{1--1,5 kuliah}

\subsection{Volume dan himpunan terukur Jordan}

Untuk suatu himpunan
$S \subset \R^n$, \emph{\myindex{fungsi karakteristik}}
atau \emph{\myindex{fungsi indikator}} $\chi_S \colon \R^n \to \R$
didefinisikan oleh
\glsadd{not:indicatorfunction}
\begin{equation*}
\chi_S(x) \coloneqq
\begin{cases}
1 & \text{jika } x \in S, \\
0 & \text{jika } x \notin S.
\end{cases}
\end{equation*}
Himpunan terbatas $S$ disebut \emph{\myindex{terukur Jordan}}\footnote{Dinamai
menurut matematikawan Prancis
\href{https://en.wikipedia.org/wiki/Camille_Jordan}{Marie Ennemond Camille Jordan}
(1838--1922).}
jika terdapat persegi panjang tertutup $R$ dengan $S \subset R$ sedemikian
sehingga $\chi_S$ terintegralkan secara Riemann, yakni
$\chi_S \in \sR(R)$.
Kondisi ini tidak bergantung pada persegi panjang yang dipilih. Jika
$S=\varnothing$, maka $\chi_S$ identik nol, sehingga terintegralkan pada setiap
persegi panjang tertutup yang memuat $S$ dan setiap integralnya bernilai nol.
Sekarang andaikan $S \neq \varnothing$. Ambil dua persegi panjang tertutup
$R$ dan $R'$ dengan $S \subset R$ dan $S \subset R'$. Karena $S$ tidak kosong,
$R \cap R'$ merupakan persegi panjang tertutup yang juga memuat $S$.
Berdasarkan \propref{mv:prop:integralsmallerset} dan
\exerciseref{mv:zerooutside}, $\chi_S \in \sR(R \cap R')$, sehingga
$\chi_S \in \sR(R')$. Jadi,
\begin{equation*}
\int_R \chi_S = \int_{R'} \chi_S = \int_{R \cap R'} \chi_S.
\end{equation*}
Kita mendefinisikan
\emph{volume berdimensi $n$}%
\index{volume berdimensi n@volume berdimensi $n$!himpunan terukur Jordan}%
\index{volume}
untuk himpunan terbatas terukur Jordan $S$ dengan
\glsadd{not:ndimvolume}
\begin{equation*}
V(S) \coloneqq \int_R \chi_S ,
\end{equation*}
dengan $R$ sebarang persegi panjang tertutup yang memuat $S$.

\begin{prop}
Himpunan terbatas $S \subset \R^n$ terukur Jordan jika dan hanya jika
batas $\partial S$ berukuran nol.
\end{prop}

\begin{proof}
Misalkan $R$ merupakan persegi panjang tertutup yang memuat $S$ di bagian
dalamnya. Jika $x \in \partial S$, maka untuk setiap $\delta > 0$, baik
$S \cap B(x,\delta)$ (tempat $\chi_S$ bernilai 1) maupun
$(R \setminus S) \cap B(x,\delta)$ (tempat $\chi_S$ bernilai 0) tidak kosong.
Jadi, $\chi_S$ tidak kontinu di $x$. Jika $x$ berada di bagian dalam $S$ atau
di komplemen penutupan $\widebar{S}$, maka $\chi_S$ identik satu atau identik
nol pada suatu lingkungan dari $x$, sehingga $\chi_S$ kontinu di $x$.
Dengan demikian, himpunan titik diskontinuitas $\chi_S$ tepat sama dengan
batas $\partial S$. Proposisi pun terbukti.
\end{proof}

\begin{prop} \label{prop:jordanmeas}
Misalkan $S$ dan $T$ merupakan himpunan terbatas terukur Jordan. Maka
\begin{enumerate}[(i)]
\item Penutupan $\widebar{S}$ terukur Jordan.
\item Bagian dalam $S^\circ$ terukur Jordan.
\item $S \cup T$ terukur Jordan.
\item $S \cap T$ terukur Jordan.
\item $S \setminus T$ terukur Jordan.
\end{enumerate}
\end{prop}

Pembuktian proposisi ini diserahkan sebagai latihan.
Selanjutnya, kita menunjukkan bahwa volume yang didefinisikan di atas sama
dengan ukuran luar yang didefinisikan sebelumnya.

\begin{prop}
Jika $S \subset \R^n$ terukur Jordan, maka $V(S) = m^*(S)$.
\end{prop}

\begin{proof}
Jika $S=\varnothing$, maka $V(S)=m^*(S)=0$, sehingga kesimpulannya langsung.
Selanjutnya, andaikan $S \neq \varnothing$. Untuk $\epsilon > 0$, misalkan
$R$ merupakan persegi panjang tertutup yang memuat $S$. Pilih suatu partisi
$P$ dari $R$ sedemikian sehingga
\begin{equation*}
U(P,\chi_S) \leq \left( \int_R \chi_S \right) + \epsilon = V(S) + \epsilon
\qquad \text{dan} \qquad
L(P,\chi_S) \geq \left( \int_R \chi_S \right) - \epsilon = V(S)-\epsilon.
\end{equation*}
Misalkan $R_1,R_2,\ldots,R_k$ merupakan semua subpersegi panjang dalam $P$
tempat $\chi_S$ tidak identik nol. Artinya, untuk setiap $j$ terdapat suatu
titik $x \in R_j$ dengan $x \in S$ (yakni $\chi_S(x)=1$). Karena $S$ tidak
kosong, $k \geq 1$. Misalkan $O_j$ merupakan persegi panjang terbuka dengan
$R_j \subset O_j$ dan
$V(O_j) < V(R_j) + \nicefrac{\epsilon}{k}$. Perhatikan bahwa
$S \subset \bigcup_j O_j$. Maka
\begin{equation*}
U(P,\chi_S) = \sum_{j=1}^k V(R_j) > 
\left(\sum_{j=1}^k V(O_j)\right) - \epsilon \geq m^*(S) - \epsilon .
\end{equation*}
Karena $U(P,\chi_S) \leq V(S) + \epsilon$, diperoleh
$m^*(S) - \epsilon \leq V(S) + \epsilon$. Karena $\epsilon>0$ sebarang,
$m^*(S) \leq V(S)$.

Misalkan $R'_1,R'_2,\ldots,R'_\ell$ merupakan semua subpersegi panjang dalam
$P$ tempat $\chi_S$ identik satu. Dengan kata lain, subpersegi panjang tersebut
termuat dalam $S$. Bagian dalam $R'^\circ_j$ dari subpersegi panjang tersebut
saling lepas dan $V(R'^\circ_j) = V(R'_j)$. Berdasarkan
\exerciseref{exercise:outermeasureofsumofrectangles},
\begin{equation*}
m^*\Bigl(\bigcup_{j=1}^\ell R'^\circ_j\Bigr)
=
\sum_{j=1}^\ell
V(R'^\circ_j) .
\end{equation*}
Oleh karena itu,
\begin{equation*}
m^*(S) \geq
m^*\Bigl(\bigcup_{j=1}^\ell R'_j\Bigr)
\geq
m^*\Bigl(\bigcup_{j=1}^\ell R'^\circ_j\Bigr)
=
\sum_{j=1}^\ell
V(R'^\circ_j)
=
\sum_{j=1}^\ell
V(R'_j)
=
L(P,\chi_S) \geq V(S) - \epsilon .
\end{equation*}
Karena $\epsilon>0$ sebarang, $m^*(S) \geq V(S)$. Dengan demikian,
$V(S)=m^*(S)$.
\end{proof}

\subsection{Pengintegralan pada himpunan terukur Jordan}

Di $\R$, pada dasarnya hanya ada satu jenis himpunan yang wajar sebagai daerah
pengintegralan, yaitu interval. Di $\R^n$ terdapat banyak jenis himpunan yang
wajar. Himpunan yang sesuai dengan integral Riemann adalah himpunan terukur
Jordan.

\begin{defn}
Misalkan $S \subset \R^n$ merupakan himpunan terbatas terukur Jordan.
Fungsi terbatas $f \colon S \to \R$ dikatakan
\emph{\myindex{terintegralkan secara Riemann pada $S$}}%
\index{terintegralkan pada $S$}, atau\glsadd{not:integrablefuncR}
$f \in \sR(S)$, jika untuk suatu persegi panjang tertutup $R$ dengan
$S \subset R$, fungsi $\widetilde{f} \colon R \to \R$ yang didefinisikan oleh
\begin{equation*}
\widetilde{f}(x) \coloneqq
\begin{cases}
f(x) & \text{jika } x \in S, \\
0 & \text{selain itu},
\end{cases}
\end{equation*}
berada di $\sR(R)$. Kita menulis
\glsadd{not:riemannintR}%
\begin{equation*}
\int_S f \coloneqq \int_R \widetilde{f}.
\end{equation*}
\end{defn}

Seperti di atas, \propref{mv:prop:integralsmallerset} dan
\exerciseref{mv:zerooutside} menunjukkan bahwa pilihan persegi panjang $R$
tidak memengaruhi hasil dan integral tersebut terdefinisi dengan baik.

Jika $f$ didefinisikan pada himpunan yang lebih besar dan kita ingin
mengintegralkannya pada $S$, definisi tersebut diterapkan pada pembatasan
$f|_S$. Secara khusus, misalkan $R$ merupakan persegi panjang tertutup,
$f \in \sR(R)$, dan $S \subset R$ terukur Jordan. Perluasan nol dari $f|_S$
ke $R$ adalah hasil kali $f\chi_S$. Karena hasil kali fungsi-fungsi yang
terintegralkan secara Riemann juga terintegralkan secara Riemann,
$f|_S \in \sR(S)$ dan
\begin{equation*}
\int_S f = \int_R f \chi_S .
\end{equation*}

\begin{prop}
Jika $S \subset \R^n$ merupakan himpunan terbatas terukur Jordan dan
$f \colon S \to \R$ kontinu dan terbatas, maka $f$ terintegralkan secara
Riemann pada $S$.
\end{prop}

\begin{proof}
Definisikan $\widetilde{f}$ seperti di atas untuk suatu persegi panjang
tertutup $R$ dengan $S \subset R$. Jika
$x \in R \setminus \widebar{S}$, maka $\widetilde{f}$ identik nol pada suatu
lingkungan dari $x$. Jika $x$ berada di bagian dalam $S$, maka
$\widetilde{f}=f$ pada suatu lingkungan dari $x$ dan kontinu di $x$. Jadi,
$\widetilde{f}$ hanya mungkin diskontinu pada $\partial S$, yang berukuran nol,
sehingga pembuktian selesai.
\end{proof}

Suatu sifat dikatakan berlaku
\emph{\myindex{untuk hampir setiap} $x$}\index{hampir setiap}, atau berlaku
\emph{\myindex{hampir di mana-mana}}, jika sifat itu berlaku untuk semua $x$
kecuali pada suatu himpunan berukuran nol. Sebagai contoh,
$f \colon S \to \R$ dan $g \colon S \to \R$ sama hampir di mana-mana jika
terdapat himpunan berukuran nol $E \subset S$ sedemikian sehingga
$f(x) = g(x)$ untuk setiap $x \in S \setminus E$.

\medskip

Banyak sifat baku integral tetap berlaku karena pada dasarnya kita masih
mengintegralkan pada persegi panjang. Bahkan, beberapa pernyataan dapat
diperkuat dengan mengganti syarat ``di setiap titik'' menjadi
\emph{hampir di mana-mana}. Pembuktian ketiga proposisi berikut diserahkan
sebagai latihan.

\begin{prop} \label{prop:jordanintbasic}
Misalkan $S \subset \R^n$ merupakan himpunan terbatas terukur Jordan,
$f \colon S \to \R$ dan $g \colon S \to \R$ terintegralkan secara Riemann
pada $S$, serta $\alpha \in \R$. Maka
\begin{enumerate}[(i)]
\item
Jika $f = 0$ hampir di mana-mana, maka $\int_S f = 0$.
\item
Jika $f = g$ hampir di mana-mana, maka $\int_S f = \int_S g$.
\item
$f+g$ terintegralkan secara Riemann pada $S$ dan
$\int_S (f+g) = \int_S f + \int_S g$.
\item
$\alpha f$ terintegralkan secara Riemann pada $S$ dan
$\int_S \alpha f = \alpha \int_S f$.
\item
Jika $f(x) \leq g(x)$ untuk hampir setiap $x$, maka
$\int_S f \leq \int_S g$.
\end{enumerate}
\end{prop}

Integral juga bersifat aditif.

\begin{prop} \label{prop:jordanintadd}
Misalkan $A \subset \R^n$ dan $B \subset \R^n$ merupakan himpunan terbatas
terukur Jordan yang saling lepas, dan $f \colon A \cup B \to \R$ sedemikian
sehingga pembatasan $f|_A$ dan $f|_B$ masing-masing terintegralkan secara
Riemann pada $A$ dan $B$. Maka $f$ terintegralkan secara Riemann pada
$A \cup B$ dan
\begin{equation*}
\int_{A \cup B} f = \int_A f + \int_B f .
\end{equation*}
\end{prop}

Akhirnya, untuk mengintegralkan pada daerah yang bukan persegi panjang dengan
teorema Fubini, cara yang lazim adalah membagi daerah itu menjadi bagian-bagian
lebih sederhana yang dapat dibatasi oleh dua grafik. Kita menyatakan hasilnya
di bidang, tetapi pernyataan serupa dapat dibuat dalam lebih banyak peubah.
Pembuktiannya sekali lagi diserahkan sebagai latihan.

\begin{prop} \label{prop:intovertypeIset}
Misalkan $f \colon [a,b] \to \R$ dan $g \colon [a,b] \to \R$ merupakan
fungsi kontinu sedemikian sehingga $f(x) < g(x)$ untuk setiap
$x \in (a,b)$. Misalkan
\begin{equation*}
U \coloneqq \bigl\{ (x,y) \in \R^2 : a < x < b \text{ dan } f(x) < y < g(x) \bigr\} .
\end{equation*}
Lihat \figureref{fig:typeIdomain}.
Maka $U$ terukur Jordan. Jika $\varphi \colon U \to \R$ kontinu dan terbatas
pada $U$, dan $\widetilde{\varphi}$ menyatakan perluasan nol $\varphi$ ke suatu
persegi panjang tertutup yang memuat $U$, maka
\begin{equation*}
\int_U \varphi =
\int_a^b \int_{f(x)}^{g(x)} \widetilde{\varphi}(x,y) \, dy \, dx .
\end{equation*}
\end{prop}

Proposisi ini juga berlaku jika $\varphi$ hanya diasumsikan terintegralkan
secara Riemann. Namun, integral dalam terhadap $y$ harus ditafsirkan sebagai
integral Darboux atas atau bawah dari fungsi satu peubah
$y \mapsto \widetilde{\varphi}(x,y)$ pada interval $[f(x),g(x)]$.

\begin{myfigureht}
\myincludegraphics{typeIdomain}{%
Diagram suatu daerah pada bidang untuk x di antara garis
x sama dengan a di kiri dan x sama dengan b di kanan.
Batas bawahnya adalah grafik y sama dengan f dari x,
dan batas atasnya adalah grafik y sama dengan g dari x.
Daerah tersebut diarsir dan garis-garis vertikal digambar untuk
menunjukkan bahwa irisan setiap garis vertikal dengan daerah itu
merupakan interval terhubung.}
\caption{Daerah di antara dua grafik.\label{fig:typeIdomain}}
\end{myfigureht}

\subsection{Citra himpunan bagian terukur Jordan}

Akhirnya, citra himpunan terukur Jordan juga terukur Jordan di bawah
pemetaan yang cukup teratur.  Untuk menyederhanakan pembahasan, kita
mengasumsikan bahwa determinan Jacobian tidak pernah bernilai nol.

\begin{prop} \label{prop:imagejordanmeas}
Misalkan $U \subset \R^n$ terbuka dan
$S \subset U$ merupakan himpunan kompak yang terukur Jordan.
Misalkan
$g \colon U \to \R^n$ merupakan pemetaan injektif
yang diferensiabel secara kontinu sedemikian sehingga
determinan Jacobian $J_g$ tidak pernah bernilai nol pada $S$.
Maka $g(S)$ terbatas dan terukur Jordan.
\end{prop}

\begin{proof}
Tetapkan $T \coloneqq g(S)$.  %Karena $S \subset \R^n$ tertutup dan terbatas, $S$ kompak.
Berdasarkan
\volIref{\lemmaref*{vI-lemma:continuouscompact} dari jilid I}{\lemmaref{lemma:continuouscompact}},
himpunan
$T$ juga kompak sehingga tertutup dan terbatas.
Kita mengklaim bahwa $\partial T \subset g(\partial S)$.  Andaikan klaim ini
telah dibuktikan.  Karena $S$ terukur Jordan,
$\partial S$ berukuran nol.  Selanjutnya, $g(\partial S)$ berukuran
nol berdasarkan \propref{prop:imagenull}.  Karena $\partial T \subset g(\partial
S)$, himpunan $T$ terukur Jordan.

Tinggal membuktikan klaim tersebut.
Karena $T$ tertutup, $\partial T \subset T$.
Misalkan $y \in \partial T$.  Maka terdapat
$x \in S$
sedemikian sehingga $g(x) = y$, dan menurut hipotesis $J_g(x) \neq 0$.
Kita gunakan teorema fungsi invers (\thmref{thm:inverse}).  Terdapat
lingkungan $V \subset U$ dari $x$ dan himpunan terbuka $W$ sedemikian sehingga
pembatasan $g|_V$ merupakan fungsi bijektif dari $V$ ke $W$
dengan fungsi invers yang diferensiabel secara kontinu.  Khususnya, $g(x) = y \in W$.
Karena $y \in \partial T$, terdapat barisan $\{ y_k \}_{k=1}^\infty$ di $W$ dengan
$\lim_{k\to\infty} y_k = y$ dan $y_k \notin T$.  Karena fungsi invers
$(g|_V)^{-1}$ kontinu, terdapat
barisan $\{ x_k \}_{k=1}^\infty$ di $V$ sedemikian sehingga $g(x_k) = y_k$ dan
$\lim_{k\to\infty} x_k = x$.
Karena $y_k \notin T = g(S)$, jelas bahwa $x_k \notin S$.  Karena $x \in S$, kita
menyimpulkan bahwa $x \in \partial S$.  Klaim tersebut terbukti: $\partial T \subset
g(\partial S)$.
\end{proof}

\subsection{Latihan}

\begin{exercise}
Buktikan \propref{prop:jordanmeas}.
\end{exercise}

\begin{exercise}
Buktikan bahwa setiap himpunan konveks terbatas terukur Jordan.  Petunjuk:
Induksi terhadap dimensi.
\end{exercise}

\begin{exercise} \label{exercise:intovertypeIset}
Buktikan \propref{prop:intovertypeIset}.  Dengan kata lain,
\begin{enumerate}[a)]
\item
Tunjukkan bahwa $U$ terukur Jordan.
\item
Buktikan bahwa
$\int_U \varphi =
\int_a^b \int_{f(x)}^{g(x)} \widetilde{\varphi}(x,y) \, dy \, dx$,
dengan $\widetilde{\varphi}$ menyatakan perluasan nol $\varphi$ ke suatu
persegi panjang tertutup yang memuat $U$.
\end{enumerate}
\end{exercise}

\begin{exercise}
Kita akan mengonstruksi contoh himpunan terbuka yang tidak terukur Jordan.
Mulailah dalam satu dimensi.  Misalkan $\{ r_j \}_{j=1}^\infty$ merupakan enumerasi
semua bilangan rasional dalam $(0,1)$.  Misalkan $(a_j,b_j)$ merupakan interval terbuka
sedemikian sehingga $(a_j,b_j) \subset (0,1)$ untuk setiap $j$, $r_j \in (a_j,b_j)$,
dan $\sum_{j=1}^\infty (b_j-a_j) < \nicefrac{1}{2}$.  Sekarang tetapkan $U \coloneqq
\bigcup_{j=1}^\infty (a_j,b_j)$.
\begin{enumerate}[a)]
\item
Tunjukkan bahwa interval-interval terbuka $(a_j,b_j)$ seperti di atas memang ada.
\item
Buktikan bahwa $\partial U = [0,1] \setminus U$.
\item
Buktikan bahwa $\partial U$ tidak berukuran nol, dan karena itu $U$ tidak terukur Jordan.
\item
Tunjukkan bahwa $W \coloneqq
\bigl( U \times (0,2) \bigr) \cup \bigl( (0,1) \times (1,2) \bigr)$
merupakan himpunan terbuka yang terhubung dan terbatas di $\R^2$
serta tidak terukur Jordan.
\end{enumerate}
\end{exercise}

\begin{exercise}
Misalkan $K \subset \R^n$ merupakan himpunan tertutup berukuran nol.
\begin{enumerate}[a)]
\item
Jika $K$ terbatas, buktikan bahwa $K$ terukur Jordan.
\item
Jika $S \subset \R^n$ terbatas dan terukur Jordan, buktikan bahwa
$S \setminus K$ terukur Jordan.
\item
Konstruksikan himpunan terbatas terukur Jordan $S \subset \R^n$
dan himpunan terbatas berukuran nol $T \subset \R^n$ sedemikian sehingga
baik $T$ maupun $S \setminus T$ tidak terukur Jordan.
\end{enumerate}
\end{exercise}

\begin{exercise}
Misalkan $U \subset \R^n$ terbuka dan $K \subset U$ kompak.
Temukan himpunan kompak terukur Jordan $S$ sedemikian sehingga $S \subset U$
dan $K \subset S^\circ$ ($K$ berada di bagian dalam $S$).
\end{exercise}

\begin{exercise} \label{exercise:closednessofriemannintegrable}
Buktikan versi \corref{cor:closednessofriemannintegrable} dengan mengganti semua
persegi panjang tertutup dengan himpunan terukur Jordan yang tertutup dan terbatas.
\end{exercise}

\begin{exercise}
Buktikan \propref{prop:jordanintbasic}.
\end{exercise}

\begin{exercise}
Buktikan \propref{prop:jordanintadd}.
\end{exercise}

%%%%%%%%%%%%%%%%%%%%%%%%%%%%%%%%%%%%%%%%%%%%%%%%%%%%%%%%%%%%%%%%%%%%%%%%%%%%%%

\sectionnewpage
\section{Teorema Green}
\label{sec:mvgreenstheorem}

\sectionnotes{1 kuliah, memerlukan \chapterref{path:chapter}}

Salah satu teorema terpenting dalam kalkulus multivariabel
adalah teorema Stokes umum,
suatu generalisasi dari teorema dasar kalkulus.
Versi dua dimensinya
disebut teorema Green%
\footnote{Dinamai menurut fisikawan matematis Britania
\href{https://en.wikipedia.org/wiki/George_Green_(mathematician)}{George Green}
(1793--1841).}.  Kita akan menyatakan teorema ini secara umum, tetapi
hanya membuktikan satu kasus khusus yang penting.

\begin{defn}
Misalkan $U \subset \R^2$ merupakan himpunan terbuka, terhubung, dan terbatas.
Andaikan batas
$\partial U$ merupakan gabungan saling lepas dari citra berhingga banyak
lintasan tertutup sederhana yang mulus sepotong-sepotong, sedemikian sehingga
setiap $p \in \partial U$ termasuk dalam penutupan
$\R^2 \setminus \widebar{U}$.
Maka $U$ disebut
\emph{\myindex{domain terbatas dengan batas mulus sepotong-sepotong}}%
\index{batas mulus sepotong-sepotong}
di $\R^2$.
\end{defn}

Kondisi tersebut menyatakan bahwa secara lokal $\partial U$ memisahkan $\R^2$
menjadi sisi \myquote{dalam} dan sisi \myquote{luar}.  Kondisi ini memastikan
bahwa $\partial U$ bukan sekadar \myquote{sayatan} di dalam $U$.
Ketika kita menelusuri setiap lintasan dengan
orientasi tertentu, sisi kiri dan kanan terdefinisi dengan baik; $U$ berada di
sisi kiri dan komplemen $U$ di sisi kanan, atau sebaliknya.
Orientasi pada $U$ adalah arah penelusuran lintasan-lintasan batas tersebut.
Kita dapat membalik orientasi jika perlu dengan mereparametrisasi lintasan.

\begin{defn}
Misalkan $U \subset \R^2$ merupakan domain terbatas dengan batas mulus
sepotong-sepotong, dan setiap komponen batas $\partial U$ telah diberi orientasi.
Untuk setiap komponen batas, misalkan $\gamma \colon [a,b] \to \R^2$
merupakan parametrisasi komponen tersebut yang memberikan orientasinya.
Tuliskan $\gamma(t) = \big(x(t),y(t)\bigr)$.  Jika pada setiap komponen dan
setiap nilai parameter reguler di setiap bagian mulus, vektor
$n(t) \coloneqq \bigl(-y'(t),x'(t)\bigr)$ mengarah ke dalam domain,
yaitu $\epsilon n(t) + \gamma(t) \in U$ untuk semua $\epsilon > 0$ yang cukup
kecil, maka $\partial U$ \emph{\myindex{berorientasi positif}}.
Lihat \figureref{fig:positiveorient}.
Sebaliknya, jika $n(t)$ mengarah ke luar domain pada setiap komponen dan setiap
nilai parameter reguler, maka $\partial U$
\emph{\myindex{berorientasi negatif}}.
\end{defn}

\begin{myfigureht}
\myincludepdft{positiveorient}{%
Dua diagram domain berarsir dengan batas yang ditunjukkan oleh panah.
Di sebelah kiri terdapat domain sederhana U yang batasnya ditandai oleh
panah-panah yang mengarah mengelilingi domain berlawanan arah jarum jam.  Pada sebuah titik batas
ditampilkan vektor gamma prima dari t sama dengan x prima dari t, y prima dari
t, yang mengarah sepanjang batas.  Vektor yang ortogonal terhadap gamma prima
dan mengarah ke dalam domain juga ditampilkan serta diberi tanda n dari t sama
dengan minus y prima dari t, x prima dari t.  Di sebelah kanan terdapat diagram
domain yang memiliki lubang.  Batas luar diorientasikan berlawanan arah jarum
jam, sedangkan batas dalam di sekitar lubang diorientasikan searah jarum jam.}
\caption{Domain berorientasi positif (kiri) dan domain berlubang yang
berorientasi positif (kanan).\label{fig:positiveorient}}
\end{myfigureht}

Vektor $n(t)$ diperoleh dengan memutar $\gamma'(t)$
sebesar $90^\circ$ berlawanan arah jarum jam, yaitu ke kiri.
Ketika kita menelusuri batas berorientasi positif
sesuai arah orientasinya,
domain berada \myquote{di sebelah kiri kita.}
Sebagai contoh,
jika $U$ merupakan domain terbatas \myquote{tanpa lubang,} yaitu $\partial U$
terhubung, maka orientasi positif berarti kita menelusuri $\partial U$
berlawanan arah jarum jam.  Jika terdapat \myquote{lubang,} kita menelusuri
batas di sekelilingnya searah jarum jam.

\begin{prop}
Misalkan $U \subset \R^2$ merupakan domain terbatas dengan batas mulus
sepotong-sepotong.  Maka $U$ terukur Jordan.
\end{prop}

\begin{proof}
Kita harus menunjukkan bahwa
$\partial U$ merupakan himpunan nol.  Karena $\partial U$ merupakan gabungan
berhingga lintasan mulus sepotong-sepotong, yang masing-masing merupakan
gabungan berhingga lintasan mulus, cukup kita tunjukkan bahwa citra suatu
lintasan mulus di $\R^2$ merupakan himpunan nol.
Misalkan $\gamma \colon [a,b] \to \R^2$ merupakan lintasan mulus.
Cukup kita tunjukkan bahwa
$\gamma\bigl((a,b)\bigr)$ merupakan himpunan nol, sebab menambahkan
titik-titik $\gamma(a)$ dan $\gamma(b)$
ke himpunan nol tetap menghasilkan himpunan nol.
Definisikan
\begin{equation*}
f \colon (a,b) \times (-1,1) \to \R^2,
\qquad \text{dengan} \qquad
f(x,y) \coloneqq \gamma(x) .
\end{equation*}
Himpunan $(a,b) \times \{ 0 \}$ merupakan himpunan nol di $\R^2$ dan
$\gamma\bigl((a,b)\bigr) = f\bigl( (a,b) \times \{ 0 \} \bigr)$.
Berdasarkan \propref{prop:imagenull},
$\gamma\bigl((a,b)\bigr)$ merupakan himpunan nol di $\R^2$;
karena itu
$\gamma\bigl([a,b]\bigr)$ merupakan himpunan nol, sehingga
$\partial U$ merupakan himpunan nol.
\end{proof}

\begin{thm}[Green]\index{teorema Green}
Misalkan $U \subset \R^2$ merupakan domain terbatas dengan batas mulus
sepotong-sepotong dan batasnya berorientasi positif.  Misalkan $P$ dan $Q$
merupakan fungsi yang diferensiabel secara kontinu pada suatu himpunan terbuka
yang memuat penutupan $\widebar{U}$.  Maka
\begin{equation*}
\int_{\partial U}
P \, dx + Q\,  dy
=
\int_{U}
\left(\frac{\partial Q}{\partial x} - \frac{\partial P}{\partial y} \right)
.
\end{equation*}
\end{thm}

Kita menyatakan teorema Green dalam bentuk umum, meskipun kita hanya akan
membuktikan versi khusus untuk jenis domain tertentu.  Versi umum mengikuti
dari kasus khusus dengan menggunakan argumen geometris tambahan: domain umum
dibagi menjadi domain-domain yang lebih kecil, kemudian kasus khusus diterapkan
pada setiap domain tersebut.
%Kita tidak akan membuktikan kasus umum.

Misalkan $U \subset \R^2$ merupakan domain terbatas dengan batas mulus
sepotong-sepotong.  Kita katakan bahwa $U$ \emph{bertipe I}%
\index{domain tipe I}
jika terdapat bilangan $a < b$ serta fungsi kontinu
$f \colon [a,b] \to \R$ dan $g \colon [a,b] \to \R$ sedemikian sehingga
\begin{equation*}
U \coloneqq \bigl\{ (x,y) \in \R^2 : a < x < b \text{ dan } f(x) < y < g(x) \bigr\} .
\end{equation*}
Demikian pula, $U$ \emph{bertipe II}\index{domain tipe II}
jika terdapat bilangan $c < d$ serta fungsi kontinu
$h \colon [c,d] \to \R$ dan $k \colon [c,d] \to \R$ sedemikian sehingga
\begin{equation*}
U \coloneqq \bigl\{ (x,y) \in \R^2 : c < y < d \text{ dan } h(y) < x < k(y) \bigr\} .
\end{equation*}
Terakhir, $U \subset \R^2$ \emph{bertipe III}\index{domain tipe III}
jika bertipe I dan bertipe II sekaligus\@.  Lihat \figureref{fig:greenstypes}.

\begin{myfigureht}
\myincludegraphics{greenstypes}{%
Diagram tiga domain.  Di sebelah kiri terdapat domain tipe I di antara dua
grafik; setiap garis vertikal memotong domain dalam suatu interval terhubung,
dan beberapa garis vertikal digambar di dalam domain.
Di tengah terdapat domain tipe II, dengan gagasan yang sama tetapi menggunakan
garis-garis horizontal.
Di sebelah kanan terdapat domain tipe III; setiap garis horizontal maupun
vertikal memotong domain dalam suatu interval terhubung, dan beberapa garis
horizontal dan vertikal digambar.}
\caption{Tipe-tipe domain untuk teorema Green.\label{fig:greenstypes}}
\end{myfigureht}

Domain yang umum digunakan untuk menerapkan teorema Green antara lain persegi
panjang dan cakram; keduanya merupakan domain tipe III.
Kita hanya akan membuktikan teorema Green untuk domain tipe III.

\begin{proof}[Bukti teorema Green untuk $U$ bertipe III]
Misalkan $f,g,h,k$ adalah fungsi-fungsi yang didefinisikan di atas.
Berdasarkan \propref{prop:intovertypeIset},
$U$ terukur Jordan; karena $U$ bertipe I\@,
\begin{equation*}
\begin{split}
\int_U 
\left(- \frac{\partial P}{\partial y} \right)
& =
\int_a^b \int_{f(x)}^{g(x)}
\left(- \frac{\partial P}{\partial y} (x,y) \right)
\, dy \, dx 
\\
& =
\int_a^b \Bigl(
- P\bigl(x,g(x)\bigr) +
P\bigl(x,f(x)\bigr)
\Bigr) \, dx
\\
& =
\int_a^b P\bigl(x,f(x)\bigr) \, dx 
-
\int_a^b P\bigl(x,g(x)\bigr) \, dx .
\end{split}
\end{equation*}
Kita mengintegralkan $P\,dx$ sepanjang batas.
Integral bentuk-1 $P\,dx$ di sepanjang ruas-ruas garis vertikal pada batas
(jika ada) bernilai nol. Oleh karena itu, bentuk ini hanya diintegralkan sepanjang
bagian atas dan bagian bawah. Sebagai parameter, $x$ bergerak dari kiri ke
kanan. Jika kita menggunakan parametrisasi yang memetakan $x$ ke
$\bigl(x,f(x)\bigr)$ dan ke $\bigl(x,g(x)\bigr)$, kita memperoleh kedua integral
lintasan di atas. Akan tetapi, arah integral lintasan kedua terbalik: bagian atas
seharusnya dilintasi dari kanan ke kiri. Setelah membalik orientasi ini, kita
peroleh
\begin{equation*}
\int_{\partial U} P \, dx
=
\int_a^b P\bigl(x,f(x)\bigr) \, dx 
+
\int_b^a P\bigl(x,g(x)\bigr) \, dx
=
\int_U 
\left(- \frac{\partial P}{\partial y} \right) .
\end{equation*}
Demikian pula, $U$ juga bertipe II\@. Integral bentuk-1 $Q\,dy$ di sepanjang
garis-garis horizontal bernilai nol. Jadi,
\begin{equation*}
\int_U 
\frac{\partial Q}{\partial x}
=
\int_c^d \int_{h(y)}^{k(y)}
\frac{\partial Q}{\partial x}(x,y)
\, dx \, dy 
=
\int_c^d \Bigl(
Q\bigl(k(y),y\bigr) 
-
Q\bigl(h(y),y\bigr)
\Bigr) \, dy 
=
\int_{\partial U} Q \, dy .
\end{equation*}
Dengan menggabungkan kedua perhitungan tersebut, kita peroleh
\begin{equation*}
\int_{\partial U} P\, dx + Q \, dy 
=
\int_{\partial U} P\, dx + \int_{\partial U} Q \, dy 
=
\int_U 
\Bigl(-\frac{\partial P}{\partial y}\Bigr)
+
\int_U 
\frac{\partial Q}{\partial x}
=
\int_U 
\Bigl(
\frac{\partial Q}{\partial x}
-\frac{\partial P}{\partial y}
\Bigr) . \qedhere
\end{equation*}
\end{proof}

Mari kita lihat bagaimana versi sederhana teorema Green (hanya untuk domain
tipe III) dapat digunakan pada lintasan yang lebih rumit.

\begin{example}
Misalkan $P(x,y) = \frac{-y}{x^2+y^2}$, 
$Q(x,y) = \frac{x}{x^2+y^2}$. Jika $(P,Q)$ kita pandang sebagai sebuah vektor,
kita memperoleh suatu \emph{\myindex{medan vektor}}. Medan vektor $(P,Q)$ ini disebut
\emph{\myindex{medan vektor pusaran}} karena memberikan kecepatan partikel
yang bergerak dalam pusaran mengelilingi titik asal. Variasi medan vektor ini
sering muncul dalam penerapan.
Misalkan $\gamma$ adalah lintasan yang menelusuri batas sebuah persegi panjang
berlawanan arah jarum jam, dan bagian dalam persegi panjang itu memuat titik asal.
Kita akan membuktikan bahwa
\begin{equation*}
\int_{\gamma} \frac{-y}{x^2+y^2} \, dx + \frac{x}{x^2+y^2} \, dy = 2 \pi .
\end{equation*}

Pertama, kita gambarkan sebuah lingkaran $C$ berjari-jari $r > 0$ yang berpusat di
titik asal, seluruhnya berada di dalam $\gamma$, dan berorientasi searah jarum
jam. Ambil $U$ sebagai domain di antara $\gamma$ dan $C$.
Lihat \figureref{fig:vortexbox}.
Integral sepanjang $\partial U$ adalah jumlah integral sepanjang
$\gamma$ dan integral sepanjang $C$.
Domain $U$ bukan bertipe III\@, sehingga kita tidak dapat langsung menerapkan
versi teorema Green yang telah kita buktikan.
Namun, jika persegi panjang tersebut kita potong sepanjang sumbu-sumbu seperti
ditunjukkan oleh garis putus-putus pada gambar, keempat domain yang dihasilkan,
sebut saja $U_1,U_2,U_3,U_4$,
bertipe III\@.
Garis-garis putus-putus itu memiliki orientasi yang berlawanan
pada dua $U_j$ yang berbagi garis tersebut.
Karena itu, ketika kita mengintegralkan sepanjang keduanya,
integral-integralnya saling meniadakan. Dengan kata lain,
\begin{multline*}
\int_{\partial U} P\, dx + Q \, dy
=
\\
\int_{\partial U_1} P\, dx + Q \, dy
+
\int_{\partial U_2} P\, dx + Q \, dy
+
\int_{\partial U_3} P\, dx + Q \, dy
+
\int_{\partial U_4} P\, dx + Q \, dy .
\end{multline*}
Sekarang kita dapat menerapkan teorema Green pada setiap $U_j$.
Pembaca diminta memeriksa bahwa di luar titik asal,
$\frac{\partial Q}{\partial x} - \frac{\partial P}{\partial y} = 0$.
Maka, untuk setiap $j$, diperoleh
\begin{equation*}
\int_{\partial U_j} P\, dx + Q \, dy 
=
\int_{U_j} \left( \frac{\partial Q}{\partial x} - \frac{\partial P}{\partial
y} \right)
=
\int_{U_j} 0 = 0 ,
\end{equation*}
dan dengan demikian $\int_{\partial U} P\, dx + Q \, dy = 0$.
Karena $\partial U$ merupakan gabungan $C$ dan $\gamma$, diperoleh
\begin{equation*}
\int_C P \, dx + Q \, dy
+
\int_{\gamma} P \, dx + Q \, dy
=
\int_{\partial U} P \, dx + Q \, dy = 0 .
\end{equation*}
Jadi, integral sepanjang $C$ adalah negatif dari integral sepanjang
$\gamma$. Integral sepanjang $C$ mudah dihitung karena pada $C$ berlaku
$x^2+y^2 = r^2$, sehingga $P(x,y) = \frac{-y}{r^2}$ dan
$Q(x,y) = \frac{x}{r^2}$.
Pembaca diminta menghitung
\begin{equation*}
\int_C P\, dx + Q \, dy
=
\int_C \frac{-y}{r^2}\, dx + \frac{x}{r^2} \, dy
=
- 2 \pi ,
\end{equation*}
dengan memperhatikan bahwa $C$ berorientasi searah jarum jam, bukan berlawanan
arah jarum jam. Dengan demikian, klaim tersebut terbukti.
\begin{myfigureht}
\myincludepdft{vortexbox}{%
Diagram domain U berupa persegi panjang dengan lubang berbentuk lingkaran di
sekitar titik asal. Batas persegi panjang dilintasi berlawanan arah jarum jam
dan diberi label gamma. Batas lingkaran di bagian dalam diberi label C dan
dilintasi searah jarum jam. Domain U dibagi menjadi empat bagian oleh empat
potongan sepanjang sumbu yang ditunjukkan sebagai garis putus-putus. Keempat
bagian diberi label U subskrip 1 hingga U subskrip 4. Batas setiap bagian
ditandai dengan arah penelusuran berlawanan arah jarum jam. Di sepanjang
garis-garis potong, batas-batas itu ditelusuri dalam kedua arah sehingga
kontribusi integralnya saling meniadakan.}
\caption{Mengubah integral sepanjang batas persegi panjang menjadi integral
sepanjang lingkaran kecil di sekitar titik asal. Domain $U$ adalah seluruh
daerah berarsir di antara lingkaran dan persegi panjang.\label{fig:vortexbox}}
\end{myfigureht}

Perhatikan bahwa jika titik asal tidak berada di dalam $\gamma$, maka
$\int_\gamma
P\,dx+Q\,dy = 0$, sebab kita dapat menerapkan teorema Green secara
langsung pada $\gamma$.
Jadi, integral ini dapat mendeteksi apakah titik asal berada di dalam $\gamma$.
\end{example}

Sebagai contoh kedua, kita akan menunjukkan kegunaan teorema Green dalam membuktikan
suatu hasil mendasar tentang fungsi harmonik.

\begin{example}
Misalkan $U \subset \R^2$ terbuka dan
$f \colon U \to \R$ harmonik; artinya, $f$ diferensiabel secara kontinu hingga
orde dua dan memenuhi \emph{\myindex{persamaan Laplace}},
$\frac{\partial^2 f}{\partial x^2} +
\frac{\partial^2 f}{\partial y^2} = 0$.
Sebagai contoh, fungsi harmonik dapat berupa distribusi panas pada keadaan tunak
atau potensial listrik di antara muatan-muatan.
Kita akan membuktikan salah satu sifat paling mendasar dari fungsi-fungsi ini.

Misalkan $D_r \coloneqq B(p,r)$ adalah cakram dengan penutupan
$\overline{D_r} = C(p,r) \subset U$. Tuliskan
$p = (x_0,y_0)$. Berilah $\partial D_r$ orientasi positif.
Lihat \exerciseref{green:balltype3orient}.
Dengan menggunakan teorema Green dan diferensiasi di bawah tanda integral,
diperoleh
\begin{equation*}
\begin{split}
0
& =
\frac{1}{2\pi r}
\int_{D_r}
\left(
\frac{\partial^2 f}{\partial x^2} +
\frac{\partial^2 f}{\partial y^2}
\right)
\\
& 
=
\frac{1}{2\pi r}
\int_{\partial D_r}
- \frac{\partial f}{\partial y} \, dx + 
\frac{\partial f}{\partial x} \, dy
\\
&
=
\frac{1}{2\pi r}
\int_0^{2\pi}
\biggl(
- \frac{\partial f}{\partial y} \bigl(x_0+r\cos(t),y_0+r\sin(t)\bigr) \bigl(-r\sin(t)\bigr)
\\
& \hspace{1.2in}
+ \frac{\partial f}{\partial x} \bigl(x_0+r\cos(t),y_0+r\sin(t)\bigr) \, r\cos(t)
\biggr) \, dt
\\
&
=
\frac{d}{dr}
\left[
\frac{1}{2\pi}
\int_0^{2\pi}
f\bigl(x_0+r\cos(t),y_0+r\sin(t)\bigr) \, dt
\right] .
\end{split}
\end{equation*}
Definisikan $g(r) \coloneqq 
\frac{1}{2\pi}
\int_0^{2\pi}
f\bigl(x_0+r\cos(t),y_0+r\sin(t)\bigr) \, dt$ untuk $r \geq 0$ yang cukup kecil.
Fungsi ini kontinu di $r=0$ (latihan), dan baru saja kita buktikan bahwa
$g'(r) = 0$ untuk semua $r > 0$.
Karena itu, $g(0) = g(r)$ untuk semua $r > 0$, dan
\begin{equation*}
g(r) = g(0) = 
\frac{1}{2\pi}
\int_0^{2\pi}
f\bigl(x_0+0\cos(t),y_0+0\sin(t)\bigr) \, dt
=
f(x_0,y_0).
\end{equation*}
Kita telah membuktikan \emph{\myindex{sifat nilai rata-rata}} fungsi harmonik:
\begin{equation*}
f(x_0,y_0) = 
\frac{1}{2\pi}
\int_0^{2\pi}
f\bigl(x_0+r\cos(t),y_0+r\sin(t)\bigr) \, dt 
=
\frac{1}{2\pi r}
\int_{\partial D_r} f \, ds .
\end{equation*}
Artinya, untuk fungsi harmonik, nilai di $p = (x_0,y_0)$ sama dengan nilai
rata-rata fungsi tersebut pada lingkaran berjari-jari sembarang $r$ yang berpusat
di $(x_0,y_0)$, asalkan $r$ cukup kecil sehingga seluruh cakram tertutup itu
termuat di dalam $U$.
\end{example}

\subsection{Latihan}

\begin{exercise} \label{green:balltype3orient}
Buktikan bahwa cakram $B(p,r) \subset \R^2$ merupakan domain tipe III, dan
buktikan bahwa orientasi yang diberikan oleh parametrisasi $\gamma(t) =
\bigl(x_0+r\cos(t),y_0+r\sin(t)\bigr)$, dengan $p = (x_0,y_0)$, merupakan
orientasi positif pada batas $\partial B(p,r)$.
\\
Catatan: Anda boleh menggunakan pengetahuan tentang sinus dan kosinus dari kalkulus.
\end{exercise}

\begin{exercise}
Buktikan bahwa domain konveks terbatas dengan batas mulus sepotong-sepotong
merupakan domain tipe III.
\end{exercise}

\begin{exercise}
Misalkan $V \subset \R^2$ merupakan domain terbatas bertipe III dengan batas
mulus sepotong-sepotong, dan misalkan $U \subset \R^2$ adalah himpunan terbuka
sedemikian sehingga $\widebar{V} \subset U$. Misalkan $f \colon U \to \R$
merupakan fungsi yang diferensiabel secara kontinu hingga orde dua. Buktikan bahwa
$\int_{\partial V}
\frac{\partial f}{\partial x} \, dx + 
\frac{\partial f}{\partial y} \, dy = 0$.
\end{exercise}

\begin{samepage}
\begin{exercise}
Untuk cakram $B(p,r) \subset \R^2$, berikan orientasi positif pada batas
$\partial B(p,r)$.
\begin{enumerate}[a)]
\item
Hitung $\displaystyle \int_{\partial B(p,r)} -y \, dx$.
\item
Hitung $\displaystyle \int_{\partial B(p,r)} x \, dy$.
\item
Hitung $\displaystyle \int_{\partial B(p,r)} \frac{-y}{2} \, dx +
\frac{x}{2} \, dy$.
\end{enumerate}
\end{exercise}
\end{samepage}

\begin{exercise}
Gunakan teorema Green untuk menunjukkan bahwa luas segitiga dengan
titik-titik sudut
$(x_1,y_1)$,
$(x_2,y_2)$,
$(x_3,y_3)$ adalah
$\frac{1}{2}\sabs{x_1y_2 + x_2 y_3 + x_3 y_1 - y_1x_2 - y_2x_3 - y_3x_1}$.
Petunjuk: Lihat latihan sebelumnya.
\end{exercise}

\begin{exercise}
Gunakan sifat nilai rata-rata untuk membuktikan
\emph{prinsip maksimum}\index{prinsip maksimum!fungsi harmonik}
bagi fungsi harmonik:
Misalkan $U \subset \R^2$ adalah himpunan terbuka terhubung dan
$f \colon U \to \R$ harmonik. Buktikan bahwa jika $f$ mencapai maksimum di
$p \in U$, maka $f$ konstan.
\end{exercise}

\begin{exercise}
Misalkan $f(x,y) \coloneqq \ln \sqrt{x^2+y^2}$.
\begin{enumerate}[a)]
\item
Tunjukkan bahwa $f$ harmonik pada domain definisinya.
\item
Tunjukkan bahwa $\lim\limits_{(x,y) \to 0} f(x,y) = -\infty$.
\item
Dengan menggunakan lingkaran $C_r$ berjari-jari
$r$ yang berpusat di titik asal, hitung $\frac{1}{2\pi r} \int_{C_r} f \, ds$.
Apa yang terjadi ketika $r \to 0$?
\item
Mengapa teorema Green tidak dapat digunakan pada bagian c)?
\end{enumerate}
\end{exercise}


%%%%%%%%%%%%%%%%%%%%%%%%%%%%%%%%%%%%%%%%%%%%%%%%%%%%%%%%%%%%%%%%%%%%%%%%%%%%%%

\sectionnewpage
\section{Penggantian variabel}
\label{sec:mvchangeofvars}

\sectionnotes{1 kuliah}

Dalam satu variabel, kita mempunyai rumus penggantian variabel yang telah
dikenal, yaitu substitusi $u$,
\begin{equation*}
\int_a^b f\bigl(g(x)\bigr) g'(x)\, dx = 
\int_{g(a)}^{g(b)} f(u) \, du .
\end{equation*}
Padanannya dalam dimensi yang lebih tinggi jauh lebih rumit.  Kerumitan pertama
berkaitan dengan orientasi.  Jika kita menggunakan definisi integral dari bab
ini, kita tidak memiliki pembedaan antara $\int_a^b$ dan $\int_b^a$.  Kita hanya
mengintegralkan pada interval $[a,b]$.  Dengan notasi ini, dan dengan
mengasumsikan bahwa $g$ injektif, rumus penggantian variabel menjadi
\begin{equation*}
\int_{[a,b]} f\bigl(g(x)\bigr) \, \babs{g'(x)}\, dx = 
\int_{g([a,b])} f(u) \, du .
\end{equation*}
Dalam bagian ini kita akan memperoleh padanan bentuk tersebut untuk beberapa
variabel.

Perhatikan peran $\babs{g'(x)}$ dalam rumus tersebut.  Secara umum, integral
mengukur volume, sehingga dalam satu dimensi integral mengukur panjang.
Perhatikan bahwa $\babs{g'(x)}$ menskalakan $dx$ dan dengan demikian juga
menskalakan panjang.  Jika $g$ linear, yaitu $g(x)=Lx$, maka $g'(x) = L$ dan
panjang interval $g([a,b])$ tepat sama dengan $\sabs{L}(b-a)$.  Hal ini karena
$g([a,b])$ adalah $[La,Lb]$ atau $[Lb,La]$.  Sifat ini berlaku pula dalam
dimensi yang lebih tinggi, dengan $\sabs{L}$ digantikan oleh nilai mutlak
determinan.

\begin{prop} \label{prop:volrectdet}
Andaikan $R \subset \R^n$ adalah suatu persegi panjang
dan $A \colon \R^n \to \R^n$ linear.  Maka
$A(R)$ terukur Jordan dan $V\bigl(A(R)\bigr) = \babs{\det (A)} \, V(R)$.
\end{prop}

\begin{proof}
Cukup membuktikan pernyataan ini untuk matriks elementer.  Buktinya diserahkan
sebagai latihan.
\end{proof}

Mari kita buktikan bahwa
nilai mutlak determinan Jacobian
$\babs{J_g(x)} = \babs{\det \bigl(g'(x)\bigr)}$
berperan sebagai pengganti $\babs{g'(x)}$ dalam rumus penggantian variabel
untuk banyak dimensi.
Teorema berikut sebenarnya berlaku dalam bentuk yang lebih umum,
tetapi pernyataan ini memadai untuk banyak keperluan.

\begin{thm}
Andaikan $U \subset \R^n$ terbuka,
$S \subset U$ merupakan himpunan kompak yang terukur Jordan, dan
$g \colon U \to \R^n$ merupakan pemetaan injektif
yang diferensiabel secara kontinu, sedemikian sehingga
$J_g$ tidak pernah bernilai nol pada $S$.
Andaikan $f \colon g(S) \to \R$ terintegralkan secara Riemann.
Maka $f \circ g$ terintegralkan secara Riemann pada $S$ dan
\begin{equation*}
\int_S f\bigl(g(x)\bigr) \, \babs{J_g(x)} \, dx
=
\int_{g(S)} f(u) \, du . 
\end{equation*}
\end{thm}

Himpunan $g(S)$ terukur Jordan berdasarkan \propref{prop:imagejordanmeas},
sehingga ruas kiri memang bermakna.
Bahwa ruas kanan bermakna menyusul dari
\corref{cor:closednessofriemannintegrable} (sebenarnya
\exerciseref{exercise:closednessofriemannintegrable}).

\begin{proof}
Himpunan $S$ dapat ditutupi oleh berhingga banyak persegi panjang tertutup
$P_1,P_2,\ldots,P_k$ yang
interiornya saling lepas dan masing-masing memenuhi $P_j \subset U$
(\exerciseref{mv:changeofvarcoverbyrects}).
Cukup membuktikan teorema untuk $P_j \cap S$ sebagai pengganti $S$
karena kita cukup menjumlahkan integral-integralnya.
Definisikan $f(y) \coloneqq 0$ untuk setiap $y \notin g(S)$.
Fungsi $f$ yang baru terintegralkan secara Riemann karena $g(S)$ terukur Jordan.
Sekarang kita dapat mengganti integral-integral di atas $S$ dengan integral-integral
di atas seluruh persegi panjang.
Oleh karena itu, tanpa mengurangi keumuman, kita andaikan bahwa $S$ merupakan persegi panjang tertutup.

Matriks $g'(x)$ invertibel untuk setiap $x \in S$,
dan matriks ini kontinu.  Oleh karena itu, ${\bigl(g'(x)\bigr)}^{-1}$
dan dengan demikian
$\bnorm{{\bigl(g'(x)\bigr)}^{-1}}$ kontinu dan tidak pernah bernilai nol.
Karena $S$ kompak, terdapat $M > 1$ sedemikian sehingga
$\bnorm{{\bigl(g'(x)\bigr)}^{-1}} \leq M$ untuk setiap $x \in S$.

Diberikan $\epsilon > 0$.
Untuk setiap $x \in S$, definisikan
\begin{equation*}
W_x \coloneqq
\left\{ y \in U : \bnorm{g'(x)-g'(y)} < \frac{\epsilon}{2M} \right\} .
\end{equation*}
Menurut \exerciseref{mv:changeofvarWxopen},
$W_x$ terbuka.
Karena $x \in W_x$ untuk setiap $x$, kita memperoleh suatu selimut terbuka.
Berdasarkan lema selimut Lebesgue
(\volIref{\lemmaref*{vI-ms:lebesgue} dari jilid I}{\lemmaref{ms:lebesgue}}),
terdapat $\delta > 0$ sedemikian sehingga
untuk setiap $y \in S$, terdapat $x$ sedemikian sehingga $B(y,\delta) \subset W_x$.
Dengan kata lain, jika $Q$ merupakan persegi panjang yang panjang sisi terpanjangnya kurang
dari $\frac{\delta}{\sqrt{n}}$ dan $y \in Q$, maka $Q \subset
B(y,\delta) \subset W_x$.  Berdasarkan ketaksamaan segitiga,
$\bnorm{g'(\xi)-g'(\eta)} < \nicefrac{\epsilon}{M}$ untuk setiap $\xi, \eta \in Q$.

Definisikan $\varphi(x) \coloneqq f\bigl(g(x)\bigr) \babs{J_g(x)}$.
Terdapat partisi $P$ dari $S$ sedemikian sehingga
$\epsilon + \int_S \varphi \geq U(P,\varphi)$.
Kita dapat mengandaikan bahwa $\delta$ cukup kecil dibanding panjang sisi
setiap subpersegi panjang dalam partisi $P$, sehingga setiap subpersegi panjang
tersebut dapat kita bagi lagi menjadi subpersegi panjang yang setiap sisinya $s$ memenuhi
$\frac{\delta}{2\sqrt{n}} \leq s \leq \frac{\delta}{\sqrt{n}}$
(\exerciseref{mv:changeofvarrectside}).
Nyatakan subpersegi panjang ini sebagai
$R_1,R_2,\ldots,R_N$.
Untuk setiap $j=1,2,\dots,N$, pilih $x_j \in R_j$ sedemikian sehingga
$\sabs{J_g(x_j)} \leq \sabs{J_g(x)}$ untuk setiap $x \in R_j$,
yang mungkin dilakukan karena $\sabs{J_g(x)}$ kontinu dan $R_j$
kompak.

Perhatikan suatu $R_j$.
Pertama, andaikan $x_j=0$, $g(0) = 0$, dan $g'(0) = I$.
Kita mengklaim bahwa
$g(R_j)$ termuat dalam suatu
persegi panjang yang volumenya paling besar $V(R_j) \, {\bigl(1+4\sqrt{n} \, \epsilon\bigr)}^n$.
Mari kita buktikan klaim ini.
Untuk sebarang $y \in R_j$,
terapkan teorema dasar kalkulus
pada fungsi $t \mapsto g(ty)$ untuk memperoleh
$g(y) = \int_0^1 g'(ty)y \,dt$.  Karena panjang setiap
sisi $R_j$ paling besar $\frac{\delta}{\sqrt{n}}$,
kita mempunyai $\snorm{y} \leq \delta$.
Perhatikan bahwa $\bnorm{g'(x)-I} < \epsilon$ karena $M > 1$,
sehingga
\begin{equation*}
%mbxlatex \begin{aligned}
\bnorm{g(y)-y}
=
\norm{\int_0^1 \bigl(g'(ty) y - y\bigr) \,dt}
%mbxlatex &
\leq
\int_0^1 \bnorm{g'(ty) y - y} \,dt
%mbxlatex \\
%mbxlatex &
\leq
\snorm{y} \int_0^1 \bnorm{g'(ty) - I} \,dt
\leq
\delta \epsilon .
%mbxlatex \end{aligned}
\end{equation*}
Oleh karena itu, $g(R_j) \subset \widetilde{R}_j$, dengan
$\widetilde{R}_j$ merupakan persegi panjang yang diperoleh dari
$R_j$ dengan memperluasnya sejauh
$\delta \epsilon$ pada setiap sisi.  Lihat \figureref{changeofvarssq:fig}.

\begin{myfigureht}
\myincludepdft{changeofvarssq}{%
Diagram sebuah persegi panjang abu-abu R sub j di dalam persegi panjang
bergaris putus-putus tilde R sub j.
Citra g dari R sub j digambar tebal sebagai versi terdeformasi dari persegi
panjang abu-abu dan seluruhnya terletak di dalam persegi panjang bergaris putus-putus.
Sisi horizontal persegi panjang abu-abu diberi label
s sub 1 dan sisi vertikalnya diberi label s sub 2.  Jarak antara
sisi-sisi persegi panjang abu-abu dan persegi panjang bergaris putus-putus diberi label delta
epsilon.  Sebuah titik di dalam semua persegi panjang diberi label
x sub j sama dengan nol sama dengan g dari x sub j.
Sebuah titik y diberi label pada sisi persegi panjang abu-abu, dan titik di dekatnya
dalam jarak delta epsilon pada sisi bangun hitam ditandai sebagai g dari y.}
\caption{Citra $R_j$ di bawah $g$ terletak di dalam
$\widetilde{R}_j$.  Sebuah titik contoh $y \in R_j$ (sebenarnya pada batas $R_j$) ditandai,
dan $g(y)$ harus terletak dalam jarak $\delta\epsilon$ dari $y$
(juga ditandai).\label{changeofvarssq:fig}}
\end{myfigureht}


Jika panjang sisi-sisi $R_j$ adalah $s_1,s_2,\ldots,s_n$, maka
$V(R_j) = s_1 s_2 \cdots s_n$.   Ingat bahwa $\delta \leq 2\sqrt{n} \, s_j$.
Jadi,
\begin{equation*}
\begin{split}
V(\widetilde{R}_j) & =
(s_1+2\delta \epsilon )
(s_2+2\delta \epsilon )
\cdots
(s_n+2\delta \epsilon )
\\
& \leq
\bigl(s_1+4 \sqrt{n}\,s_1 \epsilon \bigr)
\bigl(s_2+4 \sqrt{n}\,s_2 \epsilon \bigr)
\cdots
\bigl(s_n+4 \sqrt{n}\,s_n \epsilon \bigr)
\\
& =
s_1 \bigl(1+4 \sqrt{n}\, \epsilon \bigr)
\,
s_2 \bigl(1+4 \sqrt{n}\, \epsilon \bigr)
\cdots
s_n \bigl(1+4 \sqrt{n}\, \epsilon \bigr)
=
V(R_j) \, {\bigl(1+4\sqrt{n} \, \epsilon\bigr)}^n .
\end{split}
\end{equation*}
Klaim tersebut terbukti: $g(R_j) \subset \widetilde{R}_j$ dan
\begin{equation*}
V\bigl(g(R_j)\bigr) \leq V(\widetilde{R}_j) \leq V(R_j) \,
{\bigl(1+4\sqrt{n} \, \epsilon\bigr)}^n .
\end{equation*}

Selanjutnya, andaikan $A \coloneqq g'(0)$ tidak harus merupakan matriks identitas.
Tuliskan $g = A \circ \widetilde{g}$ dengan $\widetilde{g}'(0) = I$.
Kita mempunyai $\snorm{A^{-1}} \leq M$, sehingga
\begin{equation*}
\bnorm{\widetilde{g}'(x)-I} =
\bnorm{A^{-1}g'(x)-A^{-1}A} \leq
\bnorm{A^{-1}}\,\bnorm{g'(x)-A} < M \frac{\epsilon}{M} = \epsilon .
\end{equation*}
Dengan kata lain, klaim di atas berlaku untuk $\widetilde{g}$ karena yang
kita perlukan dalam buktinya tepatlah $\bnorm{\widetilde{g}'(x)-I} < \epsilon$.
Oleh karena itu,
$\widetilde{g}(R_j)$
termuat dalam suatu persegi panjang $\widetilde{R}_j$ dengan
$V(\widetilde{R}_j) \leq V(R_j) \, {\bigl(1+4\sqrt{n} \, \epsilon\bigr)}^n$.

Berdasarkan \propref{prop:volrectdet},
$V\bigl(A(\widetilde{R}_j)\bigr) = \babs{\det(A)} \, V(\widetilde{R}_j)$, sehingga
\begin{equation*}
\begin{split}
V\bigl(g(R_j)\bigr) & =
V\bigl(A\bigl(\widetilde{g}(R_j)\bigr)\bigr)
\\
& \leq
V\bigl(A(\widetilde{R}_j)\bigr)
\\
& \leq
\babs{\det(A)} \, V(R_j) \,
{\bigl(1+4\sqrt{n} \, \epsilon\bigr)}^n \\
& =
\babs{J_g(0)} \, V(R_j) \,
{\bigl(1+4\sqrt{n} \, \epsilon\bigr)}^n .
\end{split}
\end{equation*}
Translasi tidak mengubah volume, sehingga
untuk setiap $R_j$ dan $x_j \in R_j$, termasuk ketika $x_j \neq 0$ dan $g(x_j)
\neq 0$, kita peroleh
\begin{equation*}
V\bigl(g(R_j)\bigr) \leq
\babs{J_g(x_j)} \, V(R_j) \,
{\bigl(1+4\sqrt{n} \, \epsilon\bigr)}^n .
\end{equation*}

Tuliskan $f$ sebagai
$f = f_+ - f_-$ untuk dua fungsi nonnegatif yang terintegralkan secara Riemann,
yaitu $f_+$ dan $f_-$:
\begin{equation*}
f_+(u) \coloneqq \max \bigl\{ f(u) , 0 \bigr\}, \qquad
f_-(u) \coloneqq \max \bigl\{ -f(u) , 0 \bigr\} .
\end{equation*}
Jadi, jika kita membuktikan teorema untuk $f$ yang nonnegatif,
kita memperoleh teorema untuk $f$ sebarang.
Oleh karena itu, tanpa mengurangi keumuman, andaikan bahwa
$f(u) \geq 0$ untuk setiap $u \in g(S)$.

Karena persegi panjang $R_1,R_2,\ldots,R_N$ memberikan suatu penghalusan dari
partisi $P$,
\begin{equation*}
\begin{split}
\epsilon + \int_S f\bigl(g(x)\bigr) \, \babs{J_g(x)} \, dx
& \geq
\sum_{j=1}^N \biggl(\sup_{x \in R_j} f\bigl(g(x)\bigr) \, \babs{J_g(x)} \biggr) \, V(R_j)
\\
& \geq
\sum_{j=1}^N \biggl(\sup_{x \in R_j} f\bigl(g(x)\bigr) \biggr) \, \babs{J_g(x_j)} \, V(R_j)
\\
& \geq
\sum_{j=1}^N \biggl(\sup_{u \in g(R_j)} f(u) \biggr) \,
V\bigl(g(R_j)\bigr)
\frac{1}{{(1+4\sqrt{n} \, \epsilon)}^n}
\\
& \geq
\sum_{j=1}^N \left(\int_{g(R_j)}f(u) \,du \right)
\frac{1}{{(1+4\sqrt{n} \, \epsilon)}^n}
\\
& =
\frac{1}{{(1+4\sqrt{n} \, \epsilon)}^n}
\int_{g(S)} f(u) \,du .
\end{split}
\end{equation*}
Kesamaan terakhir berlaku karena tumpang tindih persegi panjang-persegi panjang
tersebut berada pada batas-batasnya, yang berukuran nol, sehingga citra
batas-batas tersebut juga berukuran nol.
Dengan membiarkan $\epsilon$ menuju nol, kita peroleh
\begin{equation*}
\int_S f\bigl(g(x)\bigr) \, \babs{J_g(x)} \, dx \geq \int_{g(S)} f(u) \,du .
\end{equation*}

Ingat bahwa $g^{-1}$ ada dan $g^{-1}\bigl(g(S)\bigr) = S$.
Selain itu, $1 = J_{g\circ g^{-1}} = J_g\bigl(g^{-1}(u)\bigr) \,J_{g^{-1}}(u)$ untuk $u \in g(S)$.
Jadi,
\begin{equation*}
\begin{split}
\int_{g(S)} f(u) \, du
& =
\int_{g(S)} f\Bigl(g\bigl(g^{-1}(u)\bigr)\Bigr) \,
\Babs{J_g\bigl(g^{-1}(u)\bigr)} \, \babs{J_{g^{-1}}(u)} \, du
\\
& \geq
\int_{g^{-1}(g(S))} f\bigl(g(x)\bigr) \, \babs{J_g(x)} \, dx
=
\int_{S} f\bigl(g(x)\bigr) \, \babs{J_g(x)} \, dx .
\qedhere
\end{split}
\end{equation*}
\end{proof}

\subsection{Latihan}

\begin{exercise}
Buktikan \propref{prop:volrectdet}.
\end{exercise}

\begin{exercise} \label{mv:changeofvarcoverbyrects}
Misalkan $U \subset \R^n$ terbuka
dan $S \subset U$ merupakan himpunan kompak yang terukur Jordan.
Tunjukkan bahwa terdapat berhingga banyak persegi panjang tertutup
$P_1,P_2, \ldots, P_k$ sedemikian sehingga $P_j \subset U$,
$S \subset P_1 \cup P_2 \cup \cdots \cup P_k$, dan
bagian dalamnya saling lepas, yakni
$P_j^\circ \cap P^\circ_\ell = \emptyset$ apabila $j \neq \ell$.
\end{exercise}

\begin{exercise} \label{mv:changeofvarWxopen}
Misalkan $U \subset \R^n$ terbuka, $x \in U$,
dan $g \colon U \to \R^n$ merupakan pemetaan yang diferensiabel secara kontinu.
Untuk setiap $c > 0$, tunjukkan bahwa
\begin{equation*}
\bigl\{ y \in U : \bnorm{g'(x)-g'(y)} < c \bigr\}
\end{equation*}
merupakan himpunan terbuka.
\end{exercise}

\begin{exercise} \label{mv:changeofvarrectside}
Misalkan $R \subset \R^n$ merupakan persegi panjang tertutup.
Tunjukkan bahwa jika $\delta > 0$ cukup kecil dibandingkan
dengan panjang sisi-sisi $R$, maka $R$ dapat dipartisi
menjadi subpersegi panjang sedemikian sehingga panjang setiap sisi $s$
dari setiap subpersegi panjang memenuhi $\frac{\delta}{2} \leq s \leq \delta$.
\end{exercise}

\begin{exercise}
Buktikan versi teorema berikut:
\emph{Misalkan $f \colon \R^n \to \R$ merupakan fungsi yang terintegralkan
secara Riemann dan bertumpuan kompak. Misalkan $K \subset \R^n$
merupakan tumpuan $f$, $S$ merupakan himpunan kompak,
dan $g \colon \R^n \to \R^n$ merupakan
fungsi yang pembatasannya pada suatu lingkungan $U$ dari
$S$ bersifat injektif dan diferensiabel secara kontinu,
$g(S) = K$, dan $J_g$ tidak pernah bernilai nol pada $S$ (dalam rumus,
anggap $J_g(x) = 0$ apabila $g$ dievaluasi di $x$ di luar lingkungan tersebut,
yakni ketika $x \notin
U$). Maka}
\begin{equation*}
\int_{\R^n} f\bigl(g(x)\bigr) \, \babs{J_g(x)} \, dx
=
\int_{\R^{n}} f(u) \, du .
\end{equation*}
\end{exercise}

\begin{exercise}
\pagebreak[2]
Buktikan versi teorema berikut:
\emph{Misalkan $S \subset \R^n$ merupakan himpunan terbuka, terbatas, dan terukur Jordan,
$g \colon S \to \R^n$ merupakan pemetaan injektif
yang diferensiabel secara kontinu sedemikian sehingga
$J_g$ tidak pernah bernilai nol pada $S$, dan $g(S)$ terbatas serta
terukur Jordan (himpunan ini juga terbuka).
Misalkan $f \colon g(S) \to \R$ terintegralkan secara Riemann.
Maka $f \circ g$ terintegralkan secara Riemann pada $S$ dan}
\begin{equation*}
\int_S f\bigl(g(x)\bigr) \, \babs{J_g(x)} \, dx
=
\int_{g(S)} f(u) \, du .
\end{equation*}
Petunjuk: Nyatakan $S$ sebagai gabungan menaik dari himpunan-himpunan kompak
yang terukur Jordan, lalu terapkan teorema di bagian ini pada himpunan-himpunan
tersebut. Kemudian buktikan bahwa kita dapat mengambil limit.
\end{exercise}

\end{document}
